% Environment Wellbeing and Health — Anglais 1re Tech — Cours signé OnBuch+
% Programme officiel MINESEC. Compiler avec : tectonic N03-environment-wellbeing-and-health.tex
\newcommand{\DOCMATIERE}{Anglais}
\newcommand{\DOCNIVEAU}{1re Tech}
\newcommand{\DOCMODULE}{Anglais – classe de Première technique}
\newcommand{\DOCLECON}{N03}
\newcommand{\DOCTITRE}{Environment Wellbeing and Health}
% OnBuch+ — préambule « cours signés » ENRICHI (compilé avec tectonic / XeLaTeX)
% Système visuel « Pop » d'OnBuch+ : fond crème (jamais blanc), bordures encre
% épaisses, ombres « dures » décalées, titres Archivo Black en capitales.
% Version MATHÉMATIQUES : graphes (pgfplots/tikz), boîtes pédagogiques
% (définition, propriété, méthode, attention, le savais-tu, exercice résolu,
% exercice, TP, bilan), page de garde et signature OnBuch+ ancrée en bas.
%
% Macros attendues dans main.tex AVANT \input{preamble} :
%   \DOCMATIERE  \DOCNIVEAU  \DOCMODULE  \DOCLECON (numéro)  \DOCTITRE
\documentclass[11pt]{article}

\usepackage{fontspec}
\usepackage[french]{babel}
\usepackage[a4paper,margin=2cm,top=3.4cm,bottom=2.8cm,headheight=1.4cm,footskip=1.3cm]{geometry}
\usepackage{amsmath,amssymb,amsfonts}
\usepackage{mathtools} % \xmapsto, \coloneqq... souvent utilisés par les modèles
\usepackage{cancel} % simplifications barrées (bilans de forces)
% \abs{z} (module, valeur absolue) et \norm{u} (norme), utilisés par les modèles
\providecommand{\abs}{}\let\abs\relax\DeclarePairedDelimiter{\abs}{\lvert}{\rvert}
\providecommand{\norm}{}\let\norm\relax\DeclarePairedDelimiter{\norm}{\lVert}{\rVert}
\usepackage[table]{xcolor} % [table] : fournit aussi \rowcolors (alternance de lignes)
\usepackage{pagecolor}
\usepackage{tikz}
\usetikzlibrary{arrows.meta,positioning,calc,shapes.geometric,shapes.misc,shapes.arrows,decorations.pathmorphing,decorations.pathreplacing,decorations.markings,patterns,shadows,mindmap,trees,fit,backgrounds,matrix,intersections,angles,quotes}
\usepackage{pgfplots}
\usepackage[version=4]{mhchem} % équations nucléaires \ce{^{235}_{92}U}
\usepackage[european,siunitx]{circuitikz} % schémas électriques
\pgfplotsset{compat=1.18}
\usepgfplotslibrary{fillbetween,groupplots} % aires entre courbes (calcul intégral)
\usepackage{enumitem}
\usepackage{fancyhdr}
% [most] charge toutes les bibliothèques tcolorbox (listings, minted, raster,
% documentation...) et gonfle inutilement la mémoire TeX sur un document long
% (~300 boîtes) : on ne charge que ce qui est réellement utilisé.
\usepackage[skins,breakable]{tcolorbox}
\usepackage{graphicx}
\usepackage{colortbl}
\usepackage{booktabs}
\usepackage{tabularx}
\usepackage{multirow}
\usepackage{diagbox} % case d'en-tête diagonale des tableaux à double entrée
% Colonnes extensibles centrée / gauche / droite (C, L, R), souvent utilisées par les modèles
\newcolumntype{C}{>{\centering\arraybackslash}X}
\newcolumntype{L}{>{\raggedright\arraybackslash}X}
\newcolumntype{R}{>{\raggedleft\arraybackslash}X}
\usepackage{array}
\usepackage{multicol}
\usepackage{siunitx}
% parse-numbers=false : accepte aussi \SI{2,5 \times 10^{-3}}{...}
\sisetup{locale=FR,output-decimal-marker={,},group-minimum-digits=4,parse-numbers=false}
\DeclareSIUnit\ohm{\ensuremath{\symup{\Omega}}} % U+2126 absent de la police
\DeclareSIUnit\division{div} % réglages d'oscilloscope (ms/div, V/div)
\DeclareSIUnit\tr{tr} % tours (vitesse de rotation, ex. \SI{1500}{\tr\per\minute})
\DeclareSIUnit\molar{M} % concentration molaire (ex. \SI{3}{\molar}, \SI{1}{\micro\molar})
\DeclareSIUnit\an{an} % année (le modèle confond parfois avec \year, un compteur TeX) — ex. \SI{5}{\kilo\meter\per\an}
\DeclareSIUnit\FCFA{FCFA} % franc CFA (ex. \SI{500}{\FCFA\per\kilo\gram})
\DeclareSIUnit\degreeBrix{\textdegree Brix} % teneur en sucre d'un sirop/jus (réfractométrie)
\DeclareSIUnit\month{mois} % le modèle utilise parfois \month, un compteur TeX, comme unité de durée
\DeclareSIUnit\microliter{\micro\liter} % le modèle écrit parfois \microliter en un seul token
\DeclareSIUnit\million{millions} % dénombrement (ex. \SI{15}{\million\per\mL} spermatozoïdes)
\usepackage{qrcode}
% \degree : commande fréquemment utilisée par les modèles pour un angle ou une
% température en dehors de \SI{}{} (non fournie par les paquets chargés).
\providecommand{\degree}{^{\circ}}
% \qed (fin de démonstration), utilisé par les modèles sans amsthm
\providecommand{\qed}{\hfill$\square$}
% \barre : double barre des valeurs interdites dans un tableau de variations (tkz-tab)
\providecommand{\barre}{\ensuremath{\|}}
% environnement proof (amsthm non chargé) utilisé par les modèles
\ifdefined\proof\else\newenvironment{proof}[1][Démonstration]{\par\noindent\textit{#1.}\ }{\qed\par}\fi
\usepackage{lastpage}
\usepackage{hyperref}
\hypersetup{hidelinks}

% ============================================================
% Palette « Pop » OnBuch+ — mêmes hex que la classe P (plus_ui.dart)
% ============================================================
\definecolor{poporange}{HTML}{F59321}
\definecolor{popdark}{HTML}{E07A0C}
\definecolor{popcream}{HTML}{FFF6E8}
\definecolor{popink}{HTML}{1C1714}
\definecolor{poppurple}{HTML}{7A5AE0}
\definecolor{popgreen}{HTML}{2E9E6A}
\definecolor{poppink}{HTML}{E0457B}
\definecolor{popblue}{HTML}{2F7DE1}
\definecolor{popgold}{HTML}{C98A12}
\definecolor{popmuted}{HTML}{8A7F76}
\definecolor{popline}{HTML}{EADFD2}
\definecolor{popgray}{HTML}{8A7F76}
\colorlet{popgrey}{popgray}
% Alias tolérants (le modèle nomme parfois les couleurs différemment)
\colorlet{popwhite}{white}
\colorlet{popblack}{popink}
\colorlet{popred}{poppink}
\colorlet{popyellow}{popgold}
% Teintes claires (fonds de tableaux/encadrés) — le modèle nomme parfois
% les couleurs avec un suffixe L (ex. popblueL) sans les avoir définies.
\colorlet{poporangeL}{poporange!20}
\colorlet{popdarkL}{popdark!20}
\colorlet{poppurpleL}{poppurple!20}
\colorlet{popgreenL}{popgreen!20}
\colorlet{poppinkL}{poppink!20}
\colorlet{popblueL}{popblue!20}
\colorlet{popgoldL}{popgold!20}
\colorlet{popredL}{popred!20}
\colorlet{popgrayL}{popgray!20}
% Teintes claires pour fonds de tableaux / zones
\colorlet{popblueL}{popblue!10!white}
\colorlet{poporangeL}{poporange!14!white}
\colorlet{popgreenL}{popgreen!12!white}
\colorlet{poppurpleL}{poppurple!12!white}
\colorlet{poppinkL}{poppink!12!white}
\colorlet{popgoldL}{popgold!14!white}

\pagecolor{popcream}

% ============================================================
% Typographie
% ============================================================
\defaultfontfeatures{Ligatures=TeX}
\setmainfont{PlusJakartaSans-Medium.ttf}[Path=../../../fonts/,BoldFont=PlusJakartaSans-Bold.ttf]
% Maths en sans-serif (Fira Math) avec chiffres et lettres droites de Plus
% Jakarta, pour que les formules se fondent dans le texte.
\usepackage[math-style=ISO,bold-style=ISO]{unicode-math}
\setmathfont{FiraMath-Regular.otf}[Path=../../../fonts/]
\setmathfont{PlusJakartaSans-Medium.ttf}[Path=../../../fonts/,range={up/num,up/{Latin,latin},\mathup}]
\usepackage{icomma} % virgule décimale sans espace en mode math
% Caractères mathématiques Unicode absents des polices de texte (Plus Jakarta,
% Archivo Black) : rendus par leur équivalent mathématique, en texte comme en
% formule — sinon ils s'affichent en carré vide (ex. « Arithmétique dans ℤ »).
\usepackage{newunicodechar}
\newunicodechar{ℝ}{\ensuremath{\mathbb{R}}}
\newunicodechar{ℤ}{\ensuremath{\mathbb{Z}}}
\newunicodechar{ℂ}{\ensuremath{\mathbb{C}}}
\newunicodechar{ℕ}{\ensuremath{\mathbb{N}}}
\newunicodechar{ℚ}{\ensuremath{\mathbb{Q}}}
\newunicodechar{𝔻}{\ensuremath{\mathbb{D}}}
\newunicodechar{α}{\ensuremath{\alpha}}
\newunicodechar{β}{\ensuremath{\beta}}
\newunicodechar{γ}{\ensuremath{\gamma}}
\newunicodechar{δ}{\ensuremath{\delta}}
\newunicodechar{ε}{\ensuremath{\varepsilon}}
\newunicodechar{θ}{\ensuremath{\theta}}
\newunicodechar{λ}{\ensuremath{\lambda}}
\newunicodechar{μ}{\ensuremath{\mu}}
\newunicodechar{σ}{\ensuremath{\sigma}}
\newunicodechar{τ}{\ensuremath{\tau}}
\newunicodechar{φ}{\ensuremath{\varphi}}
\newunicodechar{ψ}{\ensuremath{\psi}}
\newunicodechar{ω}{\ensuremath{\omega}}
\newunicodechar{Σ}{\ensuremath{\Sigma}}
% majuscules produites par \MakeUppercase dans les titres (α-aminés -> Α)
\newunicodechar{Α}{\ensuremath{\alpha}}
\newunicodechar{Β}{\ensuremath{\beta}}
\newunicodechar{Γ}{\ensuremath{\Gamma}}
\newunicodechar{Δ}{\ensuremath{\Delta}}
\newunicodechar{Ε}{\ensuremath{\varepsilon}}
\newunicodechar{Θ}{\ensuremath{\Theta}}
\newunicodechar{Λ}{\ensuremath{\Lambda}}
\newunicodechar{Μ}{\ensuremath{\mu}}
\newunicodechar{Τ}{\ensuremath{\tau}}
\newunicodechar{Φ}{\ensuremath{\Phi}}
\newunicodechar{Ψ}{\ensuremath{\Psi}}
\newunicodechar{Ω}{\ensuremath{\Omega}}
\newunicodechar{↦}{\ensuremath{\mapsto}}
\newunicodechar{∈}{\ensuremath{\in}}
\newunicodechar{∉}{\ensuremath{\notin}}
\newunicodechar{∧}{\ensuremath{\wedge}}
\newunicodechar{⇔}{\ensuremath{\Leftrightarrow}}
\newunicodechar{⟺}{\ensuremath{\Longleftrightarrow}}
\newunicodechar{⇒}{\ensuremath{\Rightarrow}}
\newunicodechar{⟹}{\ensuremath{\Longrightarrow}}
\newunicodechar{∘}{\ensuremath{\circ}}
\newunicodechar{∪}{\ensuremath{\cup}}
\newunicodechar{∩}{\ensuremath{\cap}}
\newunicodechar{≡}{\ensuremath{\equiv}}
\newunicodechar{⊂}{\ensuremath{\subset}}
\newunicodechar{⊥}{\ensuremath{\perp}}
\newunicodechar{∃}{\ensuremath{\exists}}
\newunicodechar{∀}{\ensuremath{\forall}}
\newunicodechar{∛}{\ensuremath{\sqrt[3]{\,}}}
\newunicodechar{∜}{\ensuremath{\sqrt[4]{\,}}}
\newunicodechar{⃗}{\ensuremath{{}^{\to}}}
\newunicodechar{ⁿ}{\ifmmode^{n}\else\textsuperscript{\ensuremath{n}}\fi}
\newunicodechar{ˣ}{\ifmmode^{x}\else\textsuperscript{\ensuremath{x}}\fi}
\newunicodechar{ᵇ}{\ifmmode^{b}\else\textsuperscript{\ensuremath{b}}\fi}
\newunicodechar{ʳ}{\ifmmode^{r}\else\textsuperscript{\ensuremath{r}}\fi}
\newunicodechar{ᵘ}{\ifmmode^{u}\else\textsuperscript{\ensuremath{u}}\fi}
\newunicodechar{ʸ}{\ifmmode^{y}\else\textsuperscript{\ensuremath{y}}\fi}
\newunicodechar{ᵃ}{\ifmmode^{a}\else\textsuperscript{\ensuremath{a}}\fi}
\newunicodechar{ᵉ}{\ifmmode^{e}\else\textsuperscript{\ensuremath{e}}\fi}
\newunicodechar{ᵏ}{\ifmmode^{k}\else\textsuperscript{\ensuremath{k}}\fi}
\newunicodechar{ᵖ}{\ifmmode^{p}\else\textsuperscript{\ensuremath{p}}\fi}
\newunicodechar{ᴺ}{\ifmmode^{N}\else\textsuperscript{\ensuremath{N}}\fi}
\newunicodechar{ᶜ}{\ifmmode^{c}\else\textsuperscript{\ensuremath{c}}\fi}
\newunicodechar{ʰ}{\ifmmode^{h}\else\textsuperscript{\ensuremath{h}}\fi}
\newunicodechar{ᵠ}{\ifmmode^{q}\else\textsuperscript{\ensuremath{q}}\fi}
\newunicodechar{ₙ}{\ifmmode_{n}\else\textsubscript{\ensuremath{n}}\fi}
\newunicodechar{ₐ}{\ifmmode_{a}\else\textsubscript{\ensuremath{a}}\fi}
\newunicodechar{ᵢ}{\ifmmode_{i}\else\textsubscript{\ensuremath{i}}\fi}
\newunicodechar{ᵣ}{\ifmmode_{r}\else\textsubscript{\ensuremath{r}}\fi}
\newunicodechar{ₖ}{\ifmmode_{k}\else\textsubscript{\ensuremath{k}}\fi}
\newunicodechar{ᵦ}{\ifmmode_{\beta}\else\textsubscript{\ensuremath{\beta}}\fi}
\newunicodechar{ₓ}{\ifmmode_{x}\else\textsubscript{\ensuremath{x}}\fi}
\newfontfamily\popdisplay{ArchivoBlack-Regular.ttf}[Path=../../../fonts/]
\newfontfamily\poplogo{SpaceGrotesk-Bold.ttf}[Path=../../../fonts/]
\newfontfamily\popsemi{PlusJakartaSans-SemiBold.ttf}[Path=../../../fonts/]
\newfontfamily\popxbold{PlusJakartaSans-ExtraBold.ttf}[Path=../../../fonts/]
\newfontfamily\popxboldls{PlusJakartaSans-ExtraBold.ttf}[Path=../../../fonts/,LetterSpace=3.0]

\setlist[enumerate]{leftmargin=1.7em,itemsep=5pt,topsep=4pt}
\setlist[itemize]{leftmargin=1.7em,itemsep=4pt,topsep=4pt,label={\color{poporange}\textbullet}}
\setlength{\parindent}{0pt}
\setlength{\parskip}{5pt}
\renewcommand{\arraystretch}{1.25}

% pgfplots : style Pop par défaut
\pgfplotsset{
  every axis/.append style={axis line style={popink,line width=1pt},tick style={popink},
    grid style={popline,line width=0.5pt},label style={font=\small},tick label style={font=\footnotesize}},
  popcurve/.style={poporange,line width=1.8pt,no markers,smooth},
}
% Même style disponible dans un \draw TikZ simple (hors pgfplots)
\tikzset{popcurve/.style={poporange,line width=1.8pt}}
\tikzset{popgrid/.style={popline,very thin}}
\colorlet{popcurve}{poporange} % le nom du style est parfois utilisé comme couleur

% ============================================================
% Pill / squiggle
% ============================================================
\newcommand{\poppill}[3]{%
  \tikz[baseline=-0.55ex]\node[draw=popink,line width=1.2pt,rounded corners=2.6pt,%
    fill=#2,inner xsep=6.5pt,inner ysep=3pt]{%
    \popxboldls\fontsize{7}{7}\selectfont\color{#3}\MakeUppercase{#1}};%
}
\newcommand{\popsquiggle}[1]{%
  \tikz[baseline=-0.4ex]{
    \draw[#1,line width=2.6pt,line cap=round]
      plot [smooth] coordinates {(0,0.09) (0.34,0.01) (0.68,0.21) (1.02,0.01) (1.36,0.10)};
  }%
}
% Mot-clé surligné (marqueur orange sous le texte)
\newcommand{\cle}[1]{{\popxbold\color{popdark}#1}}

% ============================================================
% En-tête / pied
% ============================================================
\pagestyle{fancy}
\fancyhf{}
\renewcommand{\headrulewidth}{0pt}
\renewcommand{\footrulewidth}{0pt}
\lhead{\raisebox{-0.15ex}{\poplogo\fontsize{13}{13}\selectfont\color{popink}OnBuch\color{poporange}+}}
\rhead{\poppill{\DOCMATIERE\ \textbullet\ \DOCNIVEAU\ \textbullet\ Leçon \DOCLECON}{white}{popink}}
\lfoot{\popxbold\fontsize{7.6}{7.6}\selectfont\color{popmuted}\MakeUppercase{Cours signé OnBuch+}\ \textbullet\ {\popsemi\DOCTITRE}}
\rfoot{\tikz[baseline=-0.6ex]\node[rounded corners=8.5pt,fill=popink,minimum height=17pt,minimum width=17pt,inner xsep=5pt,inner sep=0]{\popxbold\fontsize{8}{8}\selectfont\color{popcream}\thepage\,/\,\pageref*{LastPage}};}

% ============================================================
% Titres
% ============================================================
\newcommand{\chaptitre}[2]{%
  \begin{tcolorbox}[enhanced,colback=poporange,colframe=popink,boxrule=2.6pt,arc=11pt,
      drop shadow=popink,left=15pt,right=15pt,top=12pt,bottom=11pt,before skip=3pt,after skip=18pt]
    \poppill{#2}{popink}{popcream}\par\vspace{9pt}
    {\raggedright\hyphenpenalty=10000\exhyphenpenalty=10000\popdisplay\fontsize{22}{24}\selectfont\color{popcream}\MakeUppercase{#1}\par}
  \end{tcolorbox}
}
% Section numérotée : pastille orange avec le numéro + titre Archivo
\newcounter{popsec}
\newcommand{\coursec}[1]{%
  \stepcounter{popsec}\setcounter{popsub}{0}%
  \par\vspace{16pt}\noindent
  \begin{minipage}{\linewidth}
  \tikz[baseline=-0.7ex]\node[draw=popink,line width=1.6pt,fill=poporange,rounded corners=4pt,minimum size=22pt,inner sep=0]{\popdisplay\fontsize{12}{12}\selectfont\color{popcream}\thepopsec};%
  \hspace{8pt}{\popdisplay\fontsize{15}{17}\selectfont\color{popink}\MakeUppercase{#1}}\par
  \vspace{4pt}\noindent{\color{poporange}\rule{36pt}{3.6pt}}
  \end{minipage}\par\nopagebreak\vspace{8pt}
}
\newcounter{popsub}
\newcommand{\courssub}[1]{%
  \stepcounter{popsub}%
  \par\vspace{9pt}\noindent{\popxbold\fontsize{11.5}{13}\selectfont\color{popink}{\color{poporange}\thepopsec.\thepopsub}\hspace{6pt}#1}\par\nopagebreak\vspace{4pt}
}

% ============================================================
% Boîtes pédagogiques — même recette : fond blanc, bordure épaisse colorée,
% ombre dure encre, badge plein. Titre optionnel [#1].
% ============================================================
\tcbset{popbase/.style={breakable,enhanced,colback=white,boxrule=2.3pt,arc=9pt,
  shadow={3pt}{-3pt}{0pt}{popink},left=14pt,right=14pt,top=11pt,bottom=11pt,before skip=11pt,after skip=15pt,
  fontupper=\popsemi\fontsize{10.6}{14.5}\selectfont\color{popink},
  fontlower=\popsemi\fontsize{10.6}{14.5}\selectfont\color{popink}}}
% Coche dessinée (le glyphe ✓ n'existe pas dans les polices du document)
\newcommand{\popcheck}[1]{\tikz[baseline=-0.5ex]\draw[#1,line width=1.6pt,line cap=round,line join=round](-0.13,0.0)--(-0.04,-0.09)--(0.13,0.1);}
\providecommand{\coche}{\popcheck{popgreen}}
\providecommand{\resultat}[1]{\par\smallskip\noindent\fcolorbox{popgreen}{popgreenL}{\parbox{\dimexpr\linewidth-2\fboxsep-2\fboxrule\relax}{\textbf{Résultat.} #1}}\par\smallskip}
\newcommand{\popboxhead}[3]{\poppill{#1}{#2}{white}\ifx\relax#3\relax\else\hspace{6pt}{\popxbold\fontsize{10.6}{12}\selectfont\color{popink}#3}\fi\par\vspace{7pt}}

\newtcolorbox{aretenir}[1][]{popbase,colframe=popgreen,before upper={\popboxhead{À retenir}{popgreen}{#1}}}
\newtcolorbox{exemplebox}[1][]{popbase,colframe=poporange,before upper={\popboxhead{Exemple}{poporange}{#1}}}
\newtcolorbox{definition}[1][]{popbase,colframe=popblue,before upper={\popboxhead{Définition}{popblue}{#1}}}
\newtcolorbox{propriete}[1][]{popbase,colframe=poppurple,before upper={\popboxhead{Propriété}{poppurple}{#1}}}
\newtcolorbox{methode}[1][]{popbase,colframe=popgold,before upper={\popboxhead{Méthode}{popgold}{#1}}}
\newtcolorbox{attention}[1][]{popbase,colframe=poppink,before upper={\popboxhead{Attention}{poppink}{#1}}}
\newtcolorbox{experience}[1][]{popbase,colframe=popblue,colback=popblueL,before upper={\popboxhead{Expérience}{popblue}{#1}}}
% « Le savais-tu ? » : carte violette pleine (contexte, culture, Cameroun)
\newtcolorbox{savaistu}[1][]{popbase,colback=poppurple,colframe=popink,
  fontupper=\popsemi\fontsize{10.4}{14.2}\selectfont\color{white},
  before upper={\poppill{Le savais-tu\,?}{white}{poppurple}\ifx\relax#1\relax\else\hspace{6pt}{\popxbold\color{white}#1}\fi\par\vspace{7pt}}}
% Exercice résolu : énoncé en haut, \tcblower puis la solution sur fond crème
\newcounter{popexo}\newcounter{popexr}
\newtcolorbox{exoresolu}[1][]{popbase,colframe=poporange,colbacklower=popcream,
  segmentation style={popink,line width=1.2pt,dashed},
  before upper={\stepcounter{popexr}\popboxhead{Exercice résolu \thepopexr}{poporange}{#1}},
  before lower={\poppill{Solution}{popink}{popcream}\par\vspace{6pt}}}
% Exercice (non résolu en place) — la correction est dans la partie Corrigés
\newtcolorbox{exercice}[1][]{popbase,colframe=popink,
  before upper={\stepcounter{popexo}\popboxhead{Exercice \thepopexo}{popink}{#1}}}
% Corrigé
\newtcolorbox{corrige}[1][]{popbase,colframe=popgreen,colback=popgreenL,
  before upper={\popboxhead{Corrigé}{popgreen}{#1}}}
% Objectifs / prérequis / bilan
\newtcolorbox{objectifs}{popbase,colframe=popink,colback=poporangeL,
  before upper={\popboxhead{Objectifs de la leçon}{popink}{}}}
\newtcolorbox{prerequis}{popbase,colframe=popink,colback=popblueL,
  before upper={\popboxhead{Prérequis}{popblue}{}}}
\newtcolorbox{bilan}{popbase,colframe=popink,colback=white,boxrule=2.8pt,
  before upper={\popboxhead{Fiche bilan}{poporange}{L'essentiel à connaître pour l'examen}}}
% Figure encadrée avec légende
\newtcolorbox{popfigure}[1][]{enhanced,breakable=false,colback=white,colframe=popline,boxrule=1.4pt,arc=8pt,
  left=10pt,right=10pt,top=10pt,bottom=8pt,before skip=10pt,after skip=12pt,halign=center,
  lower separated=false,halign lower=center,
  fontlower=\popsemi\fontsize{9}{11.5}\selectfont\color{popmuted},
  #1}
\newcommand{\legende}[1]{\par\vspace{6pt}{\popsemi\fontsize{9}{11.5}\selectfont\color{popmuted}#1}}

% ============================================================
% Signature OnBuch+
% ============================================================
\newcommand{\poplogoinline}[1]{{\poplogo\fontsize{#1}{#1}\selectfont\color{popink}OnBuch\color{poporange}+}}
\newcommand{\signatureonbuch}{%
  \begin{tcolorbox}[enhanced,colback=popink,colframe=popink,boxrule=2.2pt,arc=10pt,
      drop shadow=poporange,left=14pt,right=14pt,top=10pt,bottom=10pt,before skip=0pt,after skip=0pt]
    \tikz[baseline=-0.7ex]\node[circle,fill=popgreen,draw=popcream,line width=1.3pt,minimum size=22pt,inner sep=0]{\popcheck{white}};%
    \hspace{9pt}\begin{minipage}[c]{0.62\linewidth}
      {\popxbold\fontsize{12}{13}\selectfont\color{popcream}Signé\ \textbullet\ {\poplogo OnBuch\color{poporange}+}}\par\vspace{2pt}
      {\raggedright\popsemi\fontsize{8}{10.5}\selectfont\color{popline}\MakeUppercase{Équipe pédagogique OnBuch+}\\\MakeUppercase{Conforme au programme officiel MINESEC}\par}
    \end{minipage}\hfill
    {\poplogo\fontsize{22}{22}\selectfont\color{popcream}OnBuch\color{poporange}+}
  \end{tcolorbox}%
}

% ============================================================
% Page de garde
% #1 titre · #2 matière · #3 niveau · #4 module · #5 n° leçon · #6 durée ·
% #7 description / objectifs (texte ou itemize)
% ============================================================
\newcommand{\pagegarde}[7]{%
  \thispagestyle{empty}
  \newgeometry{a4paper,margin=1.7cm,top=1.6cm,bottom=1.4cm}%
  \begin{tikzpicture}[remember picture,overlay]
    \fill[poporange!18] ([xshift=-3.2cm,yshift=-3.6cm]current page.north east) circle (4.6cm);
    \fill[poppurple!14] ([xshift=2.4cm,yshift=7.6cm]current page.south west) circle (3.8cm);
  \end{tikzpicture}%
  {\poplogo\fontsize{30}{30}\selectfont\color{popink}OnBuch\color{poporange}+}\hfill
  \raisebox{0.6ex}{\poppill{Cours signé \textbullet\ Premium}{popink}{popcream}}\par
  \vspace{1pt}\popsquiggle{poppurple}\par
  \vspace{8pt}
  \begin{tcolorbox}[enhanced,colback=poporange,colframe=popink,boxrule=2.8pt,arc=13pt,
      drop shadow=popink,left=18pt,right=18pt,top=12pt,bottom=13pt,after skip=10pt]
    \poppill{#2 \textbullet\ #3}{popink}{popcream}\hspace{5pt}\poppill{#4}{white}{popink}\par\vspace{8pt}
    {\popdisplay\fontsize{14}{14}\selectfont\color{popink}LEÇON #5}\par\vspace{4pt}
    {\raggedright\hyphenpenalty=10000\exhyphenpenalty=10000\popdisplay\fontsize{25}{27}\selectfont\color{popcream}\MakeUppercase{#1}\par}
    \vspace{8pt}
    {\popxbold\fontsize{9.5}{9.5}\selectfont\color{popink}\MakeUppercase{Durée conseillée : #6}}
  \end{tcolorbox}
  \begin{tcolorbox}[enhanced,colback=white,colframe=popink,boxrule=2.2pt,arc=10pt,
      drop shadow=popink,left=16pt,right=16pt,top=10pt,bottom=10pt,after skip=8pt]
    \poppill{Dans cette leçon}{poppurple}{white}\par\vspace{5pt}
    \popsemi\fontsize{9.3}{12.6}\selectfont\color{popink}#7
  \end{tcolorbox}
  \noindent\begin{minipage}[t]{0.487\linewidth}
    \begin{tcolorbox}[enhanced,colback=popgreen,colframe=popink,boxrule=2.2pt,arc=10pt,
        drop shadow=popink,left=11pt,right=11pt,top=10pt,bottom=10pt,height=3.1cm,valign=center]
      \tcbox[colback=white,colframe=popink,boxrule=1.3pt,arc=5pt,boxsep=0pt,nobeforeafter,
        left=3pt,right=3pt,top=3pt,bottom=3pt,valign=center]{\qrcode[height=1.9cm]{https://play.google.com/store/apps/details?id=com.luvvix.onbuch&hl=fr}}%
      \hspace{8pt}\begin{minipage}[c]{0.5\linewidth}
        {\popdisplay\fontsize{11}{12.5}\selectfont\color{white}\MakeUppercase{Télécharge OnBuch}}\par\vspace{4pt}
        {\raggedright\popsemi\fontsize{8.4}{11}\selectfont\color{white}Gratuit \textbullet\ Google Play\\Tuteur IA, annales, concours, résultats\par}
      \end{minipage}
    \end{tcolorbox}
  \end{minipage}\hfill
  \begin{minipage}[t]{0.487\linewidth}
    \begin{tcolorbox}[enhanced,colback=poppurple,colframe=popink,boxrule=2.2pt,arc=10pt,
        drop shadow=popink,left=11pt,right=11pt,top=10pt,bottom=10pt,height=3.1cm,valign=center]
      \tikz[baseline=-0.7ex]\node[circle,fill=white,draw=popink,line width=1.3pt,minimum size=1.6cm,inner sep=0]{\popdisplay\fontsize{24}{24}\selectfont\color{poppurple}+};%
      \hspace{8pt}\begin{minipage}[c]{0.6\linewidth}
        {\popdisplay\fontsize{11}{12.5}\selectfont\color{white}\MakeUppercase{Passe à OnBuch+}}\par\vspace{4pt}
        {\raggedright\popsemi\fontsize{8.4}{11}\selectfont\color{white}Tous les cours signés, Tuteur IA illimité, fiches PDF — onglet OnBuch+ de l'app\par}
      \end{minipage}
    \end{tcolorbox}
  \end{minipage}
  \vspace{8pt}
  \popcontenu
  \vfill
  \signatureonbuch
  \restoregeometry
  \newpage
}

% Rangée « Ce cours contient » (page de garde)
\newcommand{\poptile}[3]{% #1 couleur #2 gros texte #3 légende
  \begin{tcolorbox}[enhanced,colback=#1,colframe=popink,boxrule=2pt,arc=9pt,shadow={2.5pt}{-2.5pt}{0pt}{popink},
      left=6pt,right=6pt,top=9pt,bottom=9pt,halign=center,valign=center,height=1.75cm,width=\linewidth,nobeforeafter]
    {\popdisplay\fontsize{12.5}{13}\selectfont\color{white}\MakeUppercase{#2}}\par\vspace{3pt}
    {\centering\popsemi\fontsize{7.8}{9.5}\selectfont\color{white}#3\par}
  \end{tcolorbox}}
\newcommand{\popcontenu}{%
  \par\noindent{\popxboldls\fontsize{7.5}{7.5}\selectfont\color{popmuted}CE COURS SIGNÉ CONTIENT}\par\vspace{7pt}
  \noindent\begin{minipage}[t]{0.235\linewidth}\poptile{popblue}{Cours}{complet, illustré, pas à pas}\end{minipage}\hfill
  \begin{minipage}[t]{0.235\linewidth}\poptile{poporange}{Méthodes}{et exercices résolus}\end{minipage}\hfill
  \begin{minipage}[t]{0.235\linewidth}\poptile{poppink}{Exercices}{type examen + corrigés}\end{minipage}\hfill
  \begin{minipage}[t]{0.235\linewidth}\poptile{popgreen}{Fiche bilan}{l'essentiel à retenir}\end{minipage}\par}

% Sommaire
\newtcolorbox{sommaire}{enhanced,colback=white,colframe=popink,boxrule=2.2pt,arc=10pt,
  drop shadow=popink,left=18pt,right=18pt,top=13pt,bottom=13pt,before skip=6pt,after skip=20pt,
  before upper={\poppill{Sommaire}{poppurple}{white}\par\vspace{9pt}\popsemi\fontsize{10.6}{14.5}\selectfont\color{popink}}}

% Signature de fin de cours — poussée tout en bas de la dernière page
\newcommand{\findecours}{%
  \par\vfill\nopagebreak
  \begin{center}\popsquiggle{poporange}\end{center}\vspace{4pt}
  \signatureonbuch
}
\AtBeginDocument{\def\llbracket{\mathopen{[\mkern-3.5mu[}}\def\rrbracket{\mathclose{]\mkern-3.5mu]}}} % intervalles d'entiers (glyphe ⟦ absent de la police)
% Glyphes absents de Fira Math : redéfinis à partir de glyphes disponibles
\AtBeginDocument{%
  \def\setminus{\mathbin{\backslash}}\let\smallsetminus\setminus
  \def\vdots{\mathord{\vbox{\baselineskip3.2pt\lineskiplimit0pt\kern4pt\hbox{.}\hbox{.}\hbox{.}}}}%
  \def\ddots{\mathinner{\mkern1mu\raise6.5pt\hbox{.}\mkern2mu\raise3.6pt\hbox{.}\mkern2mu\raise0.7pt\hbox{.}\mkern1mu}}%
  \def\triangle{\mathord{\scalebox{0.9}{$\symup{\Delta}$}}}%
  \def\nabla{\mathord{\raisebox{\depth}{\scalebox{1}[-1]{$\symup{\Delta}$}}}}%
  \def\checkmark{\popcheck{popdark}}%
  \def\star{\mathbin{\ast}}\def\bigstar{\mathord{\ast}}%
  \def\diamond{\mathbin{\scalebox{0.6}{$\Diamond$}}}%
  \def\bigcup{\mathop{\mathchoice{\raisebox{-0.3em}{\scalebox{1.7}{$\cup$}}}{\scalebox{1.25}{$\cup$}}{\cup}{\cup}}\displaylimits}%
  \def\bigcap{\mathop{\mathchoice{\raisebox{-0.3em}{\scalebox{1.7}{$\cap$}}}{\scalebox{1.25}{$\cap$}}{\cap}{\cap}}\displaylimits}%
  \def\lesseqqgtr{\mathrel{\lessgtr}}\def\bigoplus{\mathop{\oplus}\displaylimits}%
}
\providecommand{\ssi}{si et seulement si} % abréviation fréquente
\AtBeginDocument{\providecommand{\wideparen}{\overparen}} % arc de cercle

%% FIGFIX-BEGIN v5 — correctifs globaux des figures (docs/onbuch-generation/tools/figfix)
\usepackage{adjustbox}\usepackage{etoolbox}\tcbuselibrary{raster}
% toute figure TikZ/pgfplots plus large que la ligne est réduite à la largeur disponible
\BeforeBeginEnvironment{tikzpicture}{\begin{adjustbox}{max width=\linewidth}}
\AfterEndEnvironment{tikzpicture}{\end{adjustbox}}
% légendes pgfplots lisibles par-dessus les courbes
\pgfplotsset{every axis/.append style={legend style={fill=white,fill opacity=0.92,text opacity=1,draw=black!15,inner sep=3pt}}}
% étiquettes d'arêtes (syntaxe edge["..."]) avec fond blanc
\tikzset{every edge quotes/.append style={fill=white,inner sep=1.2pt}}
%% FIGFIX-END
\begin{document}

\pagegarde{Environment Wellbeing and Health}{Anglais}{1re Tech}{Anglais – classe de Première technique}{N03}{13 séances (fiche de progression harmonisée)}{Cette leçon développe les compétences de communication en anglais autour de trois thèmes liés : la visite des sites touristiques, la protection de l'environnement (plantation d'arbres, jardinage, habitats naturels) et la prévention des maladies transmissibles. Elle mobilise des structures grammaticales clés (used to, tag questions, compound nouns, direct/indirect speech) à travers des situations concrètes du quotidien camerounais.
\vspace{4pt}
\begin{multicols}{2}\small
\begin{itemize}
\item À la fin de cette leçon, je sais échanger des informations sur la visite de sites touristiques camerounais en utilisant le vocabulaire approprié.
\item À la fin de cette leçon, je sais employer correctement used to, didn't use to et to be used to pour décrire des habitudes passées et présentes.
\item À la fin de cette leçon, je sais comprendre des textes audio sur la plantation d'arbres, le jardinage et la protection des habitats naturels.
\item À la fin de cette leçon, je sais former et utiliser les tag questions pour confirmer une information.
\item À la fin de cette leçon, je sais lire et exploiter des textes sur la prévention des maladies transmissibles.
\item À la fin de cette leçon, je sais former les noms composés et leurs pluriels.
\end{itemize}\end{multicols}}

\begin{sommaire}
\begin{multicols}{2}
\begin{enumerate}
\item Visiting tourist sites: speaking and grammar (used to)
\item Protecting the environment: listening, vocabulary and tag questions
\item Preventing communicable diseases: reading and compound nouns
\item Reporting what people say: direct and indirect speech (Part 1)
\item Reporting questions and commands: indirect speech (Part 2)
\item Writing and integration: communicating on health issues
\item Activité d'intégration
\item Méthodes et exercices résolus
\item Exercices
\item Corrigés des exercices
\item Fiche bilan
\end{enumerate}
\end{multicols}
\end{sommaire}

\begin{objectifs}
\begin{itemize}
\item À la fin de cette leçon, je sais échanger des informations sur la visite de sites touristiques camerounais en utilisant le vocabulaire approprié.
\item À la fin de cette leçon, je sais employer correctement used to, didn't use to et to be used to pour décrire des habitudes passées et présentes.
\item À la fin de cette leçon, je sais comprendre des textes audio sur la plantation d'arbres, le jardinage et la protection des habitats naturels.
\item À la fin de cette leçon, je sais former et utiliser les tag questions pour confirmer une information.
\item À la fin de cette leçon, je sais lire et exploiter des textes sur la prévention des maladies transmissibles.
\item À la fin de cette leçon, je sais former les noms composés et leurs pluriels.
\item À la fin de cette leçon, je sais rapporter les paroles de quelqu'un au discours indirect (déclarations, questions, ordres).
\item À la fin de cette leçon, je sais rédiger une composition pour communiquer sur les questions de santé et la prévention des maladies transmissibles.
\end{itemize}
\end{objectifs}

\begin{prerequis}
\begin{itemize}
\item Les temps de base de l'anglais (simple present, simple past, present continuous) vus en classe de Seconde
\item Le vocabulaire élémentaire de la nature et du corps humain
\item La formation des questions en anglais (auxiliaires do/does/did, inversion)
\item Les pronoms personnels, possessifs et les déterminants de base
\item La rédaction d'un paragraphe simple en anglais (niveau Seconde)
\end{itemize}
\end{prerequis}

\newpage

\coursec{Visiting tourist sites: speaking and grammar (used to)}

\textbf{Situation d'accroche.} Pendant les grandes vacances, Amina, élève de Première technique à Bafoussam, a visité le lac de Bambalang avec sa famille. De retour en classe, elle raconte à ses camarades comment le paysage a changé : \emph{``People used to plant trees here, didn't they?''} Mais la discussion dérive vite : un camarade revient de l'hôpital où son frère est soigné pour le paludisme, une maladie évitable. Comment parler des sites touristiques du Cameroun, des actions de protection de l'environnement et de la prévention des maladies transmissibles en anglais correct ? Cette leçon donne les outils linguistiques pour échanger sur le bien-être environnemental et la santé.

\textbf{Problématique.} How can we exchange information about tourist sites in Cameroon and talk about how places and habits have changed over time? In this first part of the lesson, you will learn the vocabulary of tourism, discover six major Cameroonian tourist sites, practise speaking in real situations, and master the grammar of \cle{used to} and \cle{be used to}.

\courssub{Vocabulary: words and expressions for visiting tourist sites}

Before you can speak about tourism, you need the right words. Read the list carefully, repeat each word aloud, and note the pronunciation of the underlined syllables.

\begin{center}
\begin{tabularx}{\linewidth}{|l|X|}\hline
\rowcolor{popblueL} \textbf{English word} & \textbf{Meaning / French gloss} \\ \hline
tourist attraction & un site touristique, un lieu qui attire les visiteurs \\ \hline
sightseeing & le tourisme, la visite de monuments et de sites (\emph{to go sightseeing}) \\ \hline
tour guide & un guide touristique (personne qui accompagne et explique) \\ \hline
entrance fee & le droit d'entrée, le prix du billet \\ \hline
landscape / scenery & le paysage, le décor naturel \\ \hline
wildlife reserve & une réserve de faune, une zone protégée pour les animaux \\ \hline
craft market & un marché artisanal (objets faits à la main) \\ \hline
booking / reservation & une réservation (\emph{to make a booking}) \\ \hline
itinerary & un itinéraire, le programme détaillé d'un voyage \\ \hline
opening hours & les heures d'ouverture \\ \hline
souvenir & un souvenir, un objet rapporté d'un voyage \\ \hline
accommodation & l'hébergement (hôtel, auberge, camping) \\ \hline
\end{tabularx}
\end{center}

\begin{exemplebox}[Using the vocabulary in context]
\emph{Last holidays, my family planned an \textbf{itinerary} in the West Region. Our \textbf{tour guide} took us to the Foumban \textbf{craft market}. The \textbf{entrance fee} to the Royal Palace was low, and the \textbf{landscape} around the palace was wonderful. We made a \textbf{booking} at a small hotel for two nights. The \textbf{accommodation} was basic but clean.}
\end{exemplebox}

\courssub{Major tourist sites in Cameroon}

Cameroon is often called \emph{``Africa in miniature''} because it has almost all the landscapes of the continent: beaches, mountains, savannah, rainforest and lakes. Here are six major sites you must know for your exams and general culture.

\begin{itemize}
\item \textbf{Waza National Park} (Far North Region): a famous \cle{wildlife reserve} with elephants, lions, giraffes and antelopes. It receives thousands of visitors each year.
\item \textbf{Mount Cameroon} (South West Region): an active volcano, about \num{4095} m high, near Buea. It is the highest peak in West and Central Africa. Many tourists climb it every year.
\item \textbf{Kribi beaches} (South Region): white sand beaches on the Atlantic Ocean, and the famous Lobé Falls where the river flows directly into the sea.
\item \textbf{Foumban Royal Palace} (West Region): the palace of the Bamoun sultanate, with a museum and a large craft market famous for bronze casting and embroidery.
\item \textbf{Dja Faunal Reserve} (South / East Regions): a huge rainforest reserve, a UNESCO World Heritage site since 1987, rich in biodiversity (gorillas, chimpanzees, forest elephants).
\item \textbf{Lake Bambalang} (North West Region, Ndop Plain): a large artificial lake surrounded by hills and rice fields, ideal for fishing and boat rides.
\end{itemize}

\begin{popfigure}
\begin{center}
\begin{tikzpicture}[scale=1.0, every node/.style={font=\small}]
% Contour très simplifié du Cameroun (polygone approximatif)
\draw[thick, popdark, fill=popcream] (2.2,6.4) -- (3.4,7.6) -- (4.6,6.8) -- (4.4,5.6) -- (5.2,4.6) -- (5.0,3.2) -- (4.2,2.2) -- (4.4,0.8) -- (3.0,0.2) -- (1.6,0.8) -- (0.6,1.8) -- (0.4,3.4) -- (1.2,4.6) -- (1.4,5.6) -- (2.2,6.4) -- cycle;
% Flèche du nord
\draw[-{Stealth}, thick, popblue] (5.8,6.6) -- (5.8,7.6);
\node[popblue] at (5.8,7.95) {N};
% Sites (points orange + labels)
\fill[poporange] (3.2,6.6) circle (2.2pt); \node[anchor=west, popdark] at (3.4,6.6) {1. Waza National Park};
\fill[poporange] (0.9,2.6) circle (2.2pt); \node[anchor=west, popdark] at (1.05,2.6) {2. Mount Cameroon};
\fill[poporange] (1.6,0.7) circle (2.2pt); \node[anchor=west, popdark] at (1.75,0.7) {3. Kribi beaches};
\fill[poporange] (1.7,3.7) circle (2.2pt); \node[anchor=west, popdark] at (1.85,3.7) {4. Foumban Royal Palace};
\fill[poporange] (3.6,1.2) circle (2.2pt); \node[anchor=west, popdark] at (3.75,1.2) {5. Dja Faunal Reserve};
\fill[poporange] (1.9,4.7) circle (2.2pt); \node[anchor=west, popdark] at (2.05,4.7) {6. Lake Bambalang};
\end{tikzpicture}
\end{center}
\legende{Carte schématique du Cameroun : localisation approximative de six sites touristiques majeurs (croquis simplifié, sans échelle précise ; la flèche indique le nord).}
\end{popfigure}

\begin{savaistu}[Mount Cameroon, the highest peak of West and Central Africa]
Mount Cameroon rises to about \num{4095} metres. It is the highest point in West Africa and Central Africa. It is still an active volcano: its last eruptions happened in 1999 and 2000. Every year, athletes run the famous \emph{Race of Hope} from Buea to the summit and back. In the past, fewer tourists used to climb the mountain because there were no good paths; today, trained guides lead visitors safely to the top.
\end{savaistu}

\courssub{Speaking: exchanging information on visiting tourist sites}

When you visit a tourist site, you need to ask practical questions and describe your experience. Here are useful expressions classified by function.

\begin{center}
\begin{tabularx}{\linewidth}{|l|X|}\hline
\rowcolor{popblueL} \textbf{Function} & \textbf{Useful expressions} \\ \hline
Asking for directions & \emph{Excuse me, how can I get to the craft market? / Is it far from here? / Could you show me the way to the palace?} \\ \hline
Asking about prices & \emph{How much is the entrance fee? / Is there a discount for students? / What is the price for a guided tour?} \\ \hline
Asking about opening hours & \emph{What time does the museum open? / When does it close? / Is it open on Sundays?} \\ \hline
Describing a past visit & \emph{We visited Waza last year. The wildlife was amazing! / We used to go there every holiday.} \\ \hline
Giving an opinion & \emph{In my opinion, Kribi is the best beach in Cameroon. / I really enjoyed the landscape. / I found the trip exhausting but rewarding.} \\ \hline
\end{tabularx}
\end{center}

\begin{methode}[How to exchange information about a tourist site]
\begin{enumerate}
\item \textbf{Greet} the person politely: \emph{Good morning, Sir / Madam.}
\item \textbf{Ask your question} clearly: directions, price or opening hours.
\item \textbf{Listen} to the answer and \textbf{react}: \emph{Thank you very much! / That is interesting! / I see.}
\item \textbf{Share your own experience}: say where you went, what you saw, and what you liked (use \emph{used to} for past habits).
\item \textbf{Close} the exchange politely: \emph{Have a nice day! / I hope to visit again. / Goodbye!}
\end{enumerate}
\end{methode}

\begin{exemplebox}[Guided dialogue: a tourist at the Foumban craft market]
\textbf{Tourist:} Good morning, Madam. I am visiting Foumban for the first time. How can I get to the Royal Palace from here?\par
\textbf{Craftswoman:} Good morning, Sir. It is very easy. Go straight along this road, then turn left at the mosque. The palace is just behind it.\par
\textbf{Tourist:} Thank you. And how much is the entrance fee?\par
\textbf{Craftswoman:} It is about \num{1000} francs CFA for visitors. The museum opens at 9 a.m. and closes at 5 p.m.\par
\textbf{Tourist:} That is good to know. Your market is wonderful! Did people use to make these bronze statues a long time ago?\par
\textbf{Craftswoman:} Yes, they did. Our artisans used to work only for the Sultan's palace, but today they sell to everyone. We are used to receiving many tourists now.\par
\textbf{Tourist:} Fantastic! I will buy a souvenir before I leave. Have a nice day!\par
\textbf{Craftswoman:} You too, Sir. Enjoy your visit!
\end{exemplebox}

\courssub{Grammar focus: used to + base verb}

\textbf{Intuition first.} Look at Amina's sentence: \emph{``People used to plant trees here.''} She is not talking about today. She is talking about a \textbf{past habit} that has stopped: people planted trees regularly in the past, but they do not do it anymore. English has a special structure for this idea: \cle{used to} + base verb.

\begin{definition}[Used to + base verb]
\cle{Used to} + base verb (infinitive without \emph{to}) expresses a \textbf{past habit} or a \textbf{past state} that is \textbf{no longer true} today (une habitude ou un état révolu du passé).\par
\textbf{Affirmative:} Subject + used to + base verb.\par
\emph{People \textbf{used to travel} by train to Ngaoundéré.} (Ils avaient l'habitude de voyager en train ; aujourd'hui ce n'est plus le cas.)\par
\emph{There \textbf{used to be} a dense forest around the lake.} (État passé révolu.)\par
\textbf{Negative:} Subject + didn't use to + base verb.\par
\emph{We \textbf{didn't use to have} mobile phones at school.}\par
\textbf{Question:} Did + subject + use to + base verb?\par
\emph{\textbf{Did} tourists \textbf{use to visit} Waza in large numbers?}
\end{definition}

\begin{attention}[Spelling trap: used to vs use to]
The most frequent error is to write \emph{use to} without the \textbf{-d} in affirmative sentences. Remember:\par
\begin{itemize}
\item Affirmative: always \textbf{used to} with \textbf{-d}: \emph{She \textbf{used to} live in Buea.} (never \emph{she use to live}).
\item Negative and question: after \emph{did / didn't}, the \textbf{-d disappears}: \emph{He \textbf{didn't use to} smoke.} (never \emph{didn't used to}); \emph{\textbf{Did} you \textbf{use to} play football?}
\end{itemize}
Why? Because the past tense is already carried by \emph{did}, exactly like \emph{He didn't go} (never \emph{He didn't went}).
\end{attention}

\begin{attention}[Pronunciation tip]
\begin{itemize}
\item \textbf{Used to} (habit): pronounced \textit{/juːstuː/} (the \emph{s} sounds like \emph{s} in \emph{see}, strong form) or \textit{/juːstə/} (weak form in fast speech).
\item \textbf{Use to} (after did/didn't): pronounced \textit{/juːstə/} (the \emph{s} sounds like \emph{z} in \emph{easy}).
\end{itemize}
Practise: \emph{They used to come} \textit{/juːstuː kʌm/} vs \emph{Did they use to come?} \textit{/dɪd ðeɪ juːstə kʌm/}.
\end{attention}

\begin{exemplebox}[Waza National Park: a changing site]
Waza National Park receives thousands of visitors each year. But the park has changed a lot. In the 1980s, it \textbf{used to host} even more elephants than today. Poaching and drought reduced the herds. Tourists also \textbf{used to sleep} inside the park in simple camps; now most of them stay in lodges outside. Rangers \textbf{didn't use to carry} radios, but today they communicate with modern equipment to protect the animals.
\end{exemplebox}

\courssub{Grammar focus: to be used to + noun / -ing}

\textbf{Intuition first.} Compare these two sentences:\par
\emph{My grandfather \textbf{used to live} in Maroua.} (He lived there in the past; now he does not.)\par
\emph{My grandfather \textbf{is used to} the hot climate.} (The heat is normal for him; he is accustomed to it.)\par
The structures look similar but the meanings are completely different. \cle{To be used to} + noun or verb-\emph{ing} means \emph{to be accustomed to} (être habitué à).

\begin{definition}[To be used to + noun / -ing]
\cle{To be used to} + noun or verb-\emph{ing} expresses that a situation is \textbf{familiar or normal} for someone (accoutumance présente ou passée selon le temps de \emph{be}).\par
\textbf{Structure:} Subject + be (conjugated) + used to + noun / verb-\emph{ing}.\par
\emph{Tourists \textbf{are used to} the hot climate of the Far North.}\par
\emph{She \textbf{is used to walking} long distances.} (Elle a l'habitude de marcher longtemps.)\par
\emph{We \textbf{are not used to eating} spicy food.}\par
\emph{He \textbf{was used to} the noise.} (Passé : il avait l'habitude.)
\end{definition}

\begin{methode}[How to distinguish ``used to do'' and ``be used to doing'']
\begin{enumerate}
\item \textbf{Look at the verb \emph{be}.} Is there a form of \emph{be} (am, is, are, was, were) before \emph{used to}? If yes, it is the structure of \textbf{accoutumance} (\emph{be used to}).
\item \textbf{Look at what follows.} A base verb (infinitive) follows \emph{used to do}; a noun or an \emph{-ing} form follows \emph{be used to}.
\item \textbf{Check the meaning.} \emph{Used to do} = a past habit that has stopped (révolu). \emph{Be used to doing} = something normal and familiar now (habitude présente).
\item \textbf{Translate mentally.} If you can say \emph{avant, on faisait...}, use \emph{used to do}. If you can say \emph{on est habitué à...}, use \emph{be used to doing}.
\end{enumerate}
\end{methode}

\begin{center}
\begin{tabularx}{\linewidth}{|X|X|}\hline
\rowcolor{popblueL} \textbf{used to + base verb (past habit/state, now finished)} & \textbf{be used to + noun/-ing (accustomed, familiar)} \\ \hline
\emph{People \textbf{used to travel} by train to Ngaoundéré.} (They do not anymore.) & \emph{Travellers \textbf{are used to} long journeys in the North.} (Long journeys are normal for them.) \\ \hline
\emph{We \textbf{used to swim} in Lake Bambalang every holiday.} & \emph{The fishermen \textbf{are used to swimming} in the lake.} \\ \hline
\emph{Foumban artisans \textbf{used to work} only for the Sultan.} & \emph{Foumban artisans \textbf{are used to} selling to tourists.} \\ \hline
\emph{I \textbf{didn't use to like} spicy food.} & \emph{I \textbf{am not used to} spicy food yet.} \\ \hline
\emph{There \textbf{used to be} a forest here.} & \emph{We \textbf{are used to} the noise of the city.} \\ \hline
\end{tabularx}
\end{center}

\begin{attention}[Common confusion -- do not mix!]
\begin{itemize}
\item Wrong: \emph{Tourists used to the heat.} (missing \emph{are}) $\rightarrow$ Right: \emph{Tourists \textbf{are used to} the heat.}
\item Wrong: \emph{I am used to walk there.} $\rightarrow$ Right: \emph{I am used to \textbf{walking} there.} (after \emph{be used to}, use the \emph{-ing} form).
\item Wrong: \emph{I used to walking there.} $\rightarrow$ Right: \emph{I \textbf{used to walk} there.} (after \emph{used to}, use the base verb).
\item Wrong: \emph{Did you used to go?} $\rightarrow$ Right: \emph{\textbf{Did} you \textbf{use} to go?} (no \emph{-d} after \emph{did}).
\end{itemize}
\end{attention}

\begin{exoresolu}[Choosing the correct structure]
\textbf{Énoncé.} Complete each sentence with \emph{used to} + base verb or \emph{be used to} + \emph{-ing}, using the verb in brackets.\par
1. When I was a child, my family \dots\ (visit) Kribi every December.\par
2. The guides at Waza \dots\ (work) under the sun; it is normal for them.\par
3. Tourists in the 1980s \dots\ (not / take) many photos because cameras were rare.\par
4. My grandmother \dots\ (cook) on a wood fire; she has done it all her life.
\tcblower
\textbf{Solution détaillée.}\par
1. \emph{used to visit} — a repeated action in childhood, finished today: past habit $\rightarrow$ \emph{used to} + base verb.\par
2. \emph{are used to working} — working under the sun is normal for them: accoutumance $\rightarrow$ \emph{be used to} + \emph{-ing}.\par
3. \emph{didn't use to take} — negative past habit; after \emph{didn't}, write \emph{use to} without \emph{-d}.\par
4. \emph{is used to cooking} — a lifelong habit that is still true and familiar: accoutumance $\rightarrow$ \emph{be used to} + \emph{-ing}.
\end{exoresolu}

\begin{exercice}[Practice: talking about change]
\textbf{Instructions:} Complete the sentences with the correct form (\emph{used to} + base verb, \emph{didn't use to}, or \emph{be used to} + \emph{-ing} / noun).\par
1. Fifty years ago, many people \dots\ (travel) to Douala by boat.\par
2. The rangers at Dja Reserve \dots\ (walk) in the forest; they do it every day.\par
3. \dots\ you \dots\ (go) sightseeing with your parents when you were younger?\par
4. We \dots\ (not / have) a tour guide in the past; we visited alone.\par
5. Visitors \dots\ (the entrance fee) at Foumban; they pay it without complaining.\par
6. Rewrite with \emph{used to}: \emph{In the past, people planted many trees around Lake Bambalang.}\par
7. Complete with the correct form of \emph{be used to}: My little brother \dots\ (not) \dots\ (sleep) in a tent. He finds it scary.
\end{exercice}

\begin{corrige}[Exercice: talking about change]
1. \emph{used to travel} (past habit, finished).\par
2. \emph{are used to walking} (accoutumance, present).\par
3. \emph{Did you use to go} (question about a past habit; no \emph{-d} after \emph{did}).\par
4. \emph{didn't use to have} (negative past habit).\par
5. \emph{are used to the entrance fee} (familiar situation; \emph{be used to} + noun).\par
6. \emph{People used to plant many trees around Lake Bambalang.}\par
7. \emph{is not used to sleeping} (accoutumance négative présente).
\end{corrige}

\begin{aretenir}[Key points of this section]
\begin{itemize}
\item Tourism vocabulary: \emph{tourist attraction, sightseeing, tour guide, entrance fee, landscape/scenery, wildlife reserve, craft market, booking, itinerary, opening hours, souvenir, accommodation}.
\item Six major Cameroonian sites: Waza National Park (Far North), Mount Cameroon (\num{4095} m, South West), Kribi beaches (South), Foumban Royal Palace (West), Dja Faunal Reserve (South/East, UNESCO), Lake Bambalang (North West).
\item \cle{Used to} + base verb = past habit or state, no longer true: \emph{People used to travel by train.} Negative: \emph{didn't use to}. Question: \emph{Did ... use to?}
\item \cle{Be used to} + noun/\emph{-ing} = to be accustomed to (familiar): \emph{Tourists are used to the hot climate.}
\item Pronunciation: \emph{used to} \textit{/juːstuː/} (affirmative) vs \emph{use to} \textit{/juːstə/} (after did/didn't).
\end{itemize}
\end{aretenir}

\coursec{Protecting the environment: listening, vocabulary and tag questions}

\courssub{Transition from the previous section}
In the previous section we practised talking about \cle{visiting tourist sites} and we learned how to use \cle{used to} for past habits (\emph{We used to hike in the mountains}).
Now we turn our attention to \cle{protecting the environment}. You will hear short recordings about school green clubs and community reforestation campaigns, learn the key vocabulary for planting trees and gardening, and master \cle{tag questions} to check information politely.

\courssub{Vocabulary: planting trees, gardening and protecting natural habitats}

\begin{definition}[Key vocabulary]
The following words and expressions are frequently used when we talk about tree-planting, gardening and habitat protection. Learn the English term, its French gloss and a short example sentence.
\end{definition}

\begin{center}
\begin{tabularx}{\linewidth}{|l|X|X|}
\hline
\rowcolor{popblueL}
\textbf{English term} & \textbf{French gloss} & \textbf{Example sentence} \\ \hline
seedling & jeune plant / semis & \emph{The children planted a \cle{seedling} in the school garden.} \\ \hline
nursery & pépinière & \emph{We visited the \cle{nursery} where young trees are grown.} \\ \hline
reforestation & reboisement & \emph{The \cle{reforestation} project aims to restore 500 hectares.} \\ \hline
pruning & élagage / taille & \emph{Regular \cle{pruning} helps trees grow stronger.} \\ \hline
compost & compost & \emph{We add \cle{compost} to enrich the soil.} \\ \hline
watering can & arrosoir & \emph{She filled the \cle{watering can} and watered the seedlings.} \\ \hline
habitat & habitat naturel & \emph{Deforestation destroys the natural \cle{habitat} of many animals.} \\ \hline
endangered species & espèce en danger & \emph{The gorilla is an \cle{endangered species} in this region.} \\ \hline
deforestation & déforestation & \emph{\cle{Deforestation} reduces biodiversity dramatically.} \\ \hline
biodiversity & biodiversité & \emph{Protecting forests preserves \cle{biodiversity}.} \\ \hline
\end{tabularx}
\end{center}

\begin{exemplebox}[Vocabulary in context]
Listen to the short dialogue and underline the vocabulary items you hear.

\textbf{Dialogue (audio script)}

\textit{Teacher:} ``Good morning, club members! Today we will visit the school \cle{nursery} to check the \cle{seedling}s.''

\textit{Student A:} ``Are we going to do some \cle{pruning} as well?''

\textit{Teacher:} ``Yes, and we will add \cle{compost} to the beds. Don't forget your \cle{watering can}s.''

\textit{Student B:} ``I read that \cle{deforestation} threatens the \cle{habitat} of \cle{endangered species}.''

\textit{Teacher:} ``Exactly. Our \cle{reforestation} work helps protect \cle{biodiversity}.''
\end{exemplebox}

\courssub{Listening: school green club and community reforestation}

\begin{methode}[How to succeed in a listening task]
\begin{enumerate}
\item \textbf{Read the questions first.} Identify the topic and predict the vocabulary you will hear.
\item \textbf{Highlight key words} in the questions (numbers, names, actions).
\item \textbf{During the first listening,} write only abbreviations or symbols (e.g. ``+50 trees'', ``2023'').
\item \textbf{Second listening:} complete full sentences and check your notes.
\item \textbf{After listening,} compare answers with a partner and discuss any doubts.
\end{enumerate}
\end{methode}

\begin{exemplebox}[Listening text 1 -- School Green Club]
\textbf{Title:} ``Our Green Club's Achievements''

\textbf{Script (read aloud by the teacher or played from audio):}

``Good afternoon everyone. I am Marie, the president of the Green Club at Lycée Technique de Yaoundé. Over the past five years our club has planted more than 300 trees on the school grounds. In 2019 we started with 40 seedlings. Each year we increased the number by about 20 percent. Last year, 2023, we planted 85 trees in a single campaign. We also organise weekly \cle{watering} sessions and monthly \cle{pruning} workshops. The club works closely with the local forestry office to choose native species that support \cle{biodiversity}. Thank you.''
\end{exemplebox}

\begin{exercice}[Comprehension questions -- Listening text 1]
\begin{enumerate}
\item How many trees were planted in 2019?
\item By what percentage did the club increase planting each year?
\item How many trees were planted in 2023?
\item Name two activities the club organises regularly.
\item With which institution does the club cooperate?
\end{enumerate}
\end{exercice}

\begin{corrige}[Exercice 1]
\begin{enumerate}
\item 40 seedlings.
\item About 20\% per year.
\item 85 trees.
\item Weekly watering sessions and monthly pruning workshops.
\item The local forestry office.
\end{enumerate}
\end{corrige}

\begin{exemplebox}[Listening text 2 -- Community reforestation campaign]
\textbf{Title:} ``Reforestation in the East Region''

\textbf{Script:}

``Hello, I am Paul, a community leader from Bertoua. Last month we launched a \cle{reforestation} campaign along the banks of the Lom River. Volunteers from three villages gathered 1200 \cle{seedling}s from the regional \cle{nursery}. In two days they planted all the seedlings, covering 15 hectares. The project aims to stop \cle{deforestation} and restore the \cle{habitat} of \cle{endangered species} such as the forest elephant. We will monitor the growth for the next three years. Every volunteer received a \cle{watering can} and a bag of \cle{compost}. Together we can make a difference.''
\end{exemplebox}

\begin{exercice}[Comprehension questions -- Listening text 2]
\begin{enumerate}
\item Where did the campaign take place?
\item How many seedlings were planted?
\item What area was covered?
\item Which endangered species is mentioned?
\item What did each volunteer receive?
\end{enumerate}
\end{exercice}

\begin{corrige}[Exercice 2]
\begin{enumerate}
\item Along the banks of the Lom River in the East Region.
\item 1200 seedlings.
\item 15 hectares.
\item The forest elephant.
\item A watering can and a bag of compost.
\end{enumerate}
\end{corrige}

\courssub{Grammar: Tag questions}

\begin{definition}[Tag question]
A \cle{tag question} is a short question added at the end of a statement to ask for confirmation or agreement. It consists of an auxiliary verb + a pronoun. The polarity of the tag is opposite to the polarity of the main clause.
\end{definition}

\begin{propriete}[Formation rules]
\begin{itemize}
\item Identify the auxiliary verb in the main clause (\cle{be}, \cle{do/does/did}, \cle{have}, modal).
\item Use the same auxiliary in the tag.
\item Invert the polarity: \textbf{positive statement $\rightarrow$ negative tag}, \textbf{negative statement $\rightarrow$ positive tag}.
\item The subject of the tag is always a pronoun that matches the subject of the main clause.
\end{itemize}
\end{propriete}

\begin{center}
\begin{tabularx}{\linewidth}{|l|X|X|}
\hline
\rowcolor{popblueL}
\textbf{Main clause auxiliary} & \textbf{Positive statement $\rightarrow$ Tag} & \textbf{Negative statement $\rightarrow$ Tag} \\ \hline
\cle{be} (am/is/are/was/were) & You are ready, \textbf{aren't you?} & She isn't here, \textbf{is she?} \\ \hline
\cle{do/does/did} & They plant trees, \textbf{don't they?} & He didn't water them, \textbf{did he?} \\ \hline
\cle{have} (perfect) & We have finished, \textbf{haven't we?} & It hasn't rained, \textbf{has it?} \\ \hline
Modals (can, could, will, would, should, must, might) & You can help, \textbf{can't you?} & They shouldn't cut, \textbf{should they?} \\ \hline
\end{tabularx}
\end{center}

\begin{attention}[Common mistakes]
\begin{itemize}
\item \textbf{Wrong polarity in tag:} ``You planted trees, did you?'' $\rightarrow$ correct: ``You planted trees, \textbf{didn't you?}''
\item \textbf{Forgetting polarity inversion:} ``We shouldn't cut down forests, shouldn't we?'' $\rightarrow$ correct: ``We shouldn't cut down forests, \textbf{should we?}''
\item \textbf{Using \emph{do} with \emph{be}:} ``She is late, does she?'' $\rightarrow$ correct: ``She is late, \textbf{isn't she?}''
\end{itemize}
\end{attention}

\begin{exemplebox}[Special cases]
\begin{itemize}
\item \textbf{Let's} $\rightarrow$ \textbf{shall we?} \quad ``Let's plant trees, \textbf{shall we?}''
\item \textbf{I am} $\rightarrow$ \textbf{aren't I?} \quad ``I am late, \textbf{aren't I?}'' (not \emph{amn't I})
\item \textbf{Imperative} $\rightarrow$ \textbf{will you? / won't you? / would you?} \quad ``Water the seedlings, \textbf{will you?}''
\item \textbf{There is/are} $\rightarrow$ \textbf{isn't there? / aren't there?} \quad ``There are many birds, \textbf{aren't there?}''
\end{itemize}
\end{exemplebox}

\begin{exemplebox}[Intonation]
\begin{itemize}
\item \textbf{Falling tone} -- the speaker expects agreement or confirmation. \quad ``Trees absorb carbon dioxide, \textbf{don't they?}''
\item \textbf{Rising tone} -- a genuine question; the speaker is unsure. \quad ``You haven't seen my watering can, \textbf{have you?}''
\end{itemize}
\end{exemplebox}

\begin{exoresolu}[Tag question practice]
\textbf{Task:} Add the correct tag question to each sentence.

1. The seedlings need water every day, \_\_\_\_\_\_?

2. We didn't use chemical fertiliser, \_\_\_\_\_\_?

3. She has already planted ten trees, \_\_\_\_\_\_?

4. Let's organise a cleaning day, \_\_\_\_\_\_?

5. I am responsible for the compost, \_\_\_\_\_\_?
\tcblower
\textbf{Solution:}

1. \textbf{don't they?} (positive \emph{need} $\rightarrow$ negative tag with \emph{do})

2. \textbf{did we?} (negative \emph{didn't use} $\rightarrow$ positive tag with \emph{did})

3. \textbf{hasn't she?} (positive \emph{has planted} $\rightarrow$ negative tag with \emph{has})

4. \textbf{shall we?} (special case \emph{Let's})

5. \textbf{aren't I?} (special case \emph{I am})
\end{exoresolu}

\courssub{Pronunciation and intonation practice}

\begin{methode}[Pair work -- Interview about gardening experience]
\begin{enumerate}
\item Sit face-to-face with a partner.
\item Student A asks Student B about his/her tree-planting or gardening experience using tag questions. Example: ``You have planted a tree before, \textbf{haven't you?}''
\item Student B answers briefly and adds a tag question back. Example: ``Yes, I have. You enjoy working outdoors, \textbf{don't you?}''
\item Switch roles after five questions.
\item Use both falling and rising intonation. Note which tags sound like confirmation and which sound like real questions.
\end{enumerate}
\end{methode}

\begin{exercice}[Written practice -- Tag questions]
Complete each sentence with the correct tag question.
\begin{enumerate}
\item The community planted 500 trees last year, \_\_\_\_\_\_?
\item Deforestation destroys natural habitats, \_\_\_\_\_\_?
\item We shouldn't waste water, \_\_\_\_\_\_?
\item There is a nursery near the school, \_\_\_\_\_\_?
\item Let's start a compost bin, \_\_\_\_\_\_?
\end{enumerate}
\end{exercice}

\begin{corrige}[Exercice 3]
\begin{enumerate}
\item \textbf{didn't they?}
\item \textbf{doesn't it?}
\item \textbf{should we?}
\item \textbf{isn't there?}
\item \textbf{shall we?}
\end{enumerate}
\end{corrige}

\courssub{Illustration: Tree anatomy vocabulary}

\begin{popfigure}
\begin{tikzpicture}[scale=0.9, every node/.style={font=\small}]
% Roots
\draw[popgreen!70!black, thick] (0,0) .. controls (-0.8,-0.6) and (-1.2,-1.2) .. (-1.5,-1.8);
\draw[popgreen!70!black, thick] (0,0) .. controls (0.8,-0.6) and (1.2,-1.2) .. (1.5,-1.8);
% Trunk
\draw[poporange!80!black, line width=4pt] (0,0) -- (0,3);
% Branches
\draw[poporange!80!black, line width=2pt] (0,2.5) .. controls (-1,3) and (-1.5,3.5) .. (-2,4);
\draw[poporange!80!black, line width=2pt] (0,2.5) .. controls (1,3) and (1.5,3.5) .. (2,4);
\draw[poporange!80!black, line width=2pt] (0,1.5) .. controls (-0.8,2) and (-1.2,2.3) .. (-1.5,2.6);
\draw[poporange!80!black, line width=2pt] (0,1.5) .. controls (0.8,2) and (1.2,2.3) .. (1.5,2.6);
% Leaves (clusters)
\foreach \x/\y in {-2/4, -1.5/3.8, -1/4.2, 2/4, 1.5/3.8, 1/4.2, -1.5/2.6, -1/2.3, 1.5/2.6, 1/2.3} {
  \fill[popgreen!60!black] (\x,\y) circle (0.15);
}
% Labels
\node[anchor=east, popgreen!70!black] at (-2.2,-1.5) {\cle{roots}};
\node[anchor=west, poporange!80!black] at (2.2,1.5) {\cle{trunk}};
\node[anchor=south west, poporange!80!black] at (-2.2,4.2) {\cle{branches}};
\node[anchor=south east, popgreen!60!black] at (2.2,4.2) {\cle{leaves}};
\node[anchor=north, popblue] at (0,-2.3) {\cle{seedling} (young tree)};
\end{tikzpicture}
\legende{Figure 1 -- Main parts of a tree with key vocabulary labels.}
\end{popfigure}

\courssub{Illustration: Evolution of trees planted by the school green club (2019--2023)}

\begin{popfigure}
\begin{tikzpicture}
\begin{axis}[
  ybar,
  bar width=12pt,
  width=\linewidth,
  height=8cm,
  ylabel={Number of trees planted},
  xlabel={Year},
  symbolic x coords={2019,2020,2021,2022,2023},
  xtick=data,
  ymin=0,
  ymax=100,
  ytick={0,20,40,60,80,100},
  tick label style={font=\small},
  label style={font=\small},
  nodes near coords,
  nodes near coords align={vertical},
  every node near coord/.append style={font=\tiny, popdark},
  every axis plot/.append style={fill=popgreen!70!black},
  axis lines=left,
  enlarge x limits=0.2,
]
\addplot coordinates {(2019,40) (2020,48) (2021,58) (2022,70) (2023,85)};
\end{axis}
\end{tikzpicture}
\legende{Figure 2 -- Bar chart showing the fictitious but realistic increase in trees planted by the Lycée Technique Green Club from 2019 to 2023 (about +20\% per year).}
\end{popfigure}

\courssub{Illustration: Tag question formation table (summary)}

\begin{popfigure}
\begin{center}
\begin{tabular}{|c|c|c|c|}
\hline
\rowcolor{popblueL}
\textbf{Main clause} & \textbf{Auxiliary} & \textbf{Positive $\rightarrow$ Tag} & \textbf{Negative $\rightarrow$ Tag} \\ \hline
She \emph{is} a gardener. & be & isn't she? & is she? \\ \hline
They \emph{plant} trees. & do & don't they? & do they? \\ \hline
We \emph{have} finished. & have & haven't we? & have we? \\ \hline
You \emph{can} help. & modal (can) & can't you? & can you? \\ \hline
Let's \emph{go}. & --- & shall we? & --- \\ \hline
I \emph{am} ready. & be & aren't I? & am I? \\ \hline
\end{tabular}
\end{center}
\legende{Figure 3 -- Quick reference for forming tag questions with the most common auxiliaries and special cases.}
\end{popfigure}

\courssub{Did you know? -- Cameroon's green heritage}

\begin{savaistu}[Cameroon and the Congo Basin]
Cameroon hosts part of the \cle{Congo Basin}, the world's second largest tropical forest after the Amazon. This vast forest ecosystem stores enormous quantities of carbon and shelters exceptional \cle{biodiversity} (gorillas, forest elephants, endemic birds). The \cle{reforestation} actions carried out by school green clubs and local communities contribute directly to the preservation of this world heritage.

\emph{Le Cameroun abrite une partie du bassin du Congo, deuxième forêt tropicale du monde après l'Amazonie. Les actions de reboisement des clubs verts scolaires et des communautés locales contribuent à préserver ce patrimoine mondial.}
\end{savaistu}

\courssub{Quantitative example -- Impact of planting trees}

\begin{exemplebox}[Carbon absorption estimate]
A single mature tree can absorb around \SI{20}{kg} of CO$_2$ per year.
If the school club plants \textbf{50} trees and they all reach maturity, the annual absorption would be:
\[
50 \times \SI{20}{kg} = \SI{1000}{kg} = \SI{1}{tonne} \text{ of CO}_2 \text{ per year.}
\]
\textbf{Tag question practice:} ``Planting 50 trees makes a difference, \textbf{doesn't it?}''
\end{exemplebox}

\courssub{Integration activity -- Mini-project}

\begin{methode}[Mini-project: Design a school awareness poster]
\begin{enumerate}
\item Work in groups of three.
\item Choose one environmental action (tree planting, composting, water saving, protecting a local habitat).
\item Create a poster (A3 size) that includes:
  \begin{itemize}
  \item A catchy title in English.
  \item At least five vocabulary items from this section (e.g. \cle{seedling}, \cle{compost}, \cle{biodiversity}).
  \item Two tag questions to engage the reader (e.g. ``You want a greener school, \textbf{don't you?}'').
  \item A simple bar chart or drawing showing a target (e.g. ``Goal: 100 trees in 2025'').
  \end{itemize}
\item Present the poster to the class in 2 minutes, using tag questions to involve the audience.
\end{enumerate}
\end{methode}

\courssub{Self-evaluation checklist}

\begin{exercice}[Check your progress]
Tick the boxes you can do confidently.
\begin{itemize}
\item[$\square$] I can list and spell 10 vocabulary items related to tree planting and habitat protection.
\item[$\square$] I can understand the main ideas of a short listening text about a reforestation campaign.
\item[$\square$] I can take notes using abbreviations while listening.
\item[$\square$] I can form correct tag questions with \emph{be}, \emph{do/does/did}, \emph{have}, modals and the special cases.
\item[$\square$] I can distinguish falling and rising intonation in tag questions.
\item[$\square$] I can interview a classmate about gardening experience using tag questions.
\item[$\square$] I can write a short paragraph (50--80 words) describing a local environmental action, including at least two tag questions.
\end{itemize}
\end{exercice}

\begin{corrige}[Self-evaluation]
Discuss your ticks with a partner. For any unchecked item, review the corresponding box (vocabulary list, listening scripts, tag-question table, pronunciation practice) and redo the related exercise.
\end{corrige}

\coursec{Preventing communicable diseases: reading and compound nouns}

In the previous section, we listened to texts about planting trees and protecting natural habitats, and we practised \emph{tag questions}. A clean environment is the first step towards good health. But even in a clean environment, diseases can spread from person to person. In this section, we move from protecting nature to protecting our bodies: we will read texts about \cle{preventing communicable diseases} (les maladies transmissibles), learn the vocabulary of health, and study a very useful grammar point: \cle{compound nouns} (les noms composés).

\courssub{What is a communicable disease?}

\begin{definition}[Communicable disease]
A \cle{communicable disease} (maladie transmissible) is an illness that can pass from one person, animal or environment to another person, through air, water, food, insects or direct contact. Examples: malaria, cholera, typhoid, tuberculosis and COVID-19.
\end{definition}

To understand prevention, we must first understand \cle{transmission} (la transmission). A disease always travels along a \emph{chain}: it leaves a \textbf{source} (la source), moves through a \textbf{vector or agent} (le vecteur ou l'agent), and reaches a \textbf{host} (l'hôte, la personne infectée). Prevention means \emph{breaking the chain} at one or more points.

\begin{popfigure}
\begin{center}
\begin{tikzpicture}[font=\small, >=Stealth, node distance=1.6cm]
\node[draw, rounded corners, fill=poporangeL, minimum width=2.8cm, minimum height=1.1cm] (source) {\textbf{Source}};
\node[draw, rounded corners, fill=poppurpleL, minimum width=2.8cm, minimum height=1.1cm, right=2.4cm of source] (vector) {\textbf{Vector / Agent}};
\node[draw, rounded corners, fill=popblueL, minimum width=2.8cm, minimum height=1.1cm, right=2.4cm of vector] (host) {\textbf{Host}};
\node[below=0.15cm of source, text width=3cm, align=center, color=popmuted] {sick person, dirty water, contaminated food};
\node[below=0.15cm of vector, text width=3cm, align=center, color=popmuted] {mosquito, air, hands, water};
\node[below=0.15cm of host, text width=3cm, align=center, color=popmuted] {a healthy person};
\draw[->, very thick, popdark] (source) -- (vector);
\draw[->, very thick, popdark] (vector) -- (host);
\node[draw=popgreen, very thick, fill=popgreenL, rounded corners, below=1.9cm of vector, text width=5cm, align=center] (prev) {\textbf{Prevention breaks the chain}\\ handwashing, mosquito nets, vaccines, clean water, quarantine};
\draw[->, very thick, popgreen, dashed] (prev.north west) -- ($(source)!0.5!(vector)$);
\draw[->, very thick, popgreen, dashed] (prev.north east) -- ($(vector)!0.5!(host)$);
\end{tikzpicture}
\end{center}
\legende{The chain of transmission of a communicable disease: the germ leaves a source, travels through a vector or agent, and reaches a host. Prevention measures (green) interrupt the chain.}
\end{popfigure}

\begin{savaistu}[Did you know?]
Cholera is transmitted mainly through \textbf{contaminated water} (eau contaminée). One of the most effective prevention measures is extremely simple: \textbf{washing your hands with soap}. During cholera outbreaks in Cameroonian towns such as Douala, Yaoundé or Bafoussam, health districts always start with handwashing campaigns before anything else.
\end{savaistu}

\courssub{Vocabulary: health issues and prevention}

Here are the key words you need to read and speak about disease prevention. Learn them with their French glosses. Pay attention to the stress in compound nouns (usually on the first element: \emph{MOSquito net, HEALTH centre, HANDwashing}).

\begin{center}
\begin{tabularx}{\linewidth}{|l|X|l|}\hline
\rowcolor{popblueL} \textbf{English} & \textbf{Meaning / example} & \textbf{Français} \\ \hline
symptom & Fever is a symptom of malaria. & symptôme \\ \hline
transmission & The transmission of cholera is through dirty water. & transmission \\ \hline
vaccine & The COVID-19 vaccine protects the body. & vaccin \\ \hline
hygiene & Good hygiene prevents many diseases. & hygiène \\ \hline
mosquito net & Sleep under a treated mosquito net. & moustiquaire \\ \hline
sanitation & Poor sanitation causes typhoid. & assainissement \\ \hline
outbreak & There was a cholera outbreak in Douala. & épidémie \\ \hline
handwashing & Handwashing with soap saves lives. & lavage des mains \\ \hline
quarantine & The patient stayed in quarantine for ten days. & quarantaine \\ \hline
immune system & Vaccines help the immune system fight germs. & système immunitaire \\ \hline
health centre & Go to the health centre for a test. & centre de santé \\ \hline
blood test & A blood test confirms the diagnosis. & analyse de sang \\ \hline
hand sanitizer & Use hand sanitizer when water is not available. & gel hydroalcoolique \\ \hline
\end{tabularx}
\end{center}

\begin{exemplebox}[A figure to remember]
Sleeping under a \textbf{treated mosquito net} (moustiquaire imprégnée) reduces the risk of malaria infection significantly. In Cameroon, treated nets distributed by health districts protect \textbf{millions of children} every year, especially during the rainy season when mosquitoes multiply.
\end{exemplebox}

\courssub{Reading strategies: how to read a health text}

Before we read, let us learn three strategies that good readers use. These strategies are essential for the Probatoire and Baccalauréat technique reading comprehension papers.

\begin{methode}[Three reading strategies]
\begin{enumerate}
\item \textbf{Skimming} (lecture globale): read the text quickly, only to find the \emph{general idea}. Look at the title, the first sentence of each paragraph and any pictures. Do not read every word.
\item \textbf{Scanning} (recherche ciblée): move your eyes fast through the text to find \emph{one specific piece of information} (a date, a number, a symptom, a place). Do not read every word.
\item \textbf{Guessing meaning from context}: when you meet an unknown word, do not stop. Use the surrounding words to guess its meaning.
\end{enumerate}
\end{methode}

\begin{methode}[Guessing an unknown word from context]
\begin{enumerate}
\item Read the whole sentence and the sentences before and after the word.
\item Look for words of the \textbf{same family}: \emph{infect}, \emph{infection}, \emph{infectious}.
\item If it is a \textbf{compound noun}, break it into parts: \emph{handwashing} = \emph{hand} + \emph{washing} = washing your hands.
\item Check if the text itself gives a definition or an example (\emph{such as}, \emph{that is}, \emph{for example}, \emph{meaning}).
\end{enumerate}
\end{methode}

\courssub{Guided reading: a public health poster}

Let us apply these strategies to an authentic document: a poster from a Cameroonian health district.

\begin{exemplebox}[Reading text: health district poster]
\textbf{BAFOUSSAM HEALTH DISTRICT --- STOP CHOLERA!}

The rainy season is here, and with it comes the danger of a cholera \textbf{outbreak}. Cholera is a communicable disease transmitted through contaminated water and food. Its main \textbf{symptoms} are violent diarrhoea, vomiting and rapid dehydration.

\textbf{Protect your family!}
\begin{itemize}
\item Wash your hands with soap before eating and after using the toilet. \textbf{Handwashing} saves lives!
\item Drink only treated or boiled water.
\item Wash fruits and vegetables with clean water.
\item Cook food well and eat it hot.
\item Keep your environment clean: good \textbf{sanitation} stops transmission.
\item If someone has symptoms, take them to the nearest \textbf{health centre} immediately. A rapid \textbf{blood test} and stool test will confirm the disease.
\end{itemize}

Together, we can break the chain of transmission!
\end{exemplebox}

\textbf{Comprehension questions (skimming and scanning):}
\begin{enumerate}
\item (Skimming) What is the general purpose of this poster?
\item (Scanning) What are the three main symptoms of cholera mentioned?
\item (Scanning) Where must a sick person be taken?
\item (Context) What does \emph{dehydration} mean? Use the symptoms to guess.
\item (Scanning) Name two prevention measures mentioned in the poster.
\end{enumerate}

\begin{corrige}[Reading questions]
\begin{enumerate}
\item The poster wants to inform the population and prevent a cholera outbreak during the rainy season.
\item Violent diarrhoea, vomiting and rapid dehydration.
\item To the nearest health centre, immediately.
\item The symptoms are diarrhoea and vomiting: the body loses a lot of water. So \emph{dehydration} means ``losing too much water from the body'' (déshydratation).
\item Wash hands with soap / drink treated or boiled water / wash fruits and vegetables / cook food well / keep environment clean (any two).
\end{enumerate}
\end{corrige}

\courssub{Grammar: compound nouns}

Look again at the poster. It contains \emph{handwashing}, \emph{health centre}, \emph{blood test}, \emph{mosquito net}, \emph{hand sanitizer}. These are \cle{compound nouns}.

\begin{definition}[Compound noun]
A \cle{compound noun} (nom composé) is a noun made of \textbf{two or more words} that function together as \textbf{one unit of meaning}: \emph{health centre} (centre de santé), \emph{mosquito net} (moustiquaire), \emph{toothbrush} (brosse à dents).
\end{definition}

\textbf{Formation of compound nouns.} There are four main patterns:

\begin{center}
\begin{tabularx}{\linewidth}{|l|X|X|}\hline
\rowcolor{popblueL} \textbf{Pattern} & \textbf{Examples} & \textbf{Plural} \\ \hline
noun + noun & health centre, blood test, mosquito net, toothbrush, hand sanitizer, waiting room, outbreak & health centres, blood tests, mosquito nets, toothbrushes, hand sanitizers, waiting rooms, outbreaks \\ \hline
adjective + noun & waiting room, public health, grown-up, emergency room & waiting rooms, grown-ups, emergency rooms \\ \hline
verb + noun & washing machine, swimming pool, drinking water & washing machines, swimming pools \\ \hline
prepositional compound & mother-in-law, passer-by, disease outbreak, check-up & mothers-in-law, passers-by, disease outbreaks, check-ups \\ \hline
\end{tabularx}
\end{center}

\begin{propriete}[Plurals of compound nouns]
The plural is generally formed on the \textbf{main word} (le mot principal) of the compound:
\begin{itemize}
\item \emph{mother-in-law} $\rightarrow$ \emph{mothers-in-law} (the main word is \emph{mother})
\item \emph{passer-by} $\rightarrow$ \emph{passers-by}
\item \emph{health centre} $\rightarrow$ \emph{health centres} (main word: \emph{centre})
\item \emph{hand sanitizer} $\rightarrow$ \emph{hand sanitizers} (main word: \emph{sanitizer})
\end{itemize}
When there is \textbf{no main word} (e.g. compounds formed from adjective/participle + adverb/preposition), the plural goes on the \textbf{last element}: \emph{grown-up} $\rightarrow$ \emph{grown-ups}, \emph{check-up} $\rightarrow$ \emph{check-ups}.
\end{propriete}

\begin{attention}[Spelling trap]
Do not write compound nouns joined or separated at random! \emph{Toothbrush} and \emph{handwashing} are written as \textbf{one word} (soudé), but \emph{health centre}, \emph{mosquito net} and \emph{waiting room} are written as \textbf{two separate words} (séparé). Some compounds use a hyphen: \emph{mother-in-law}, \emph{passer-by}, \emph{check-up}. When in doubt, check a dictionary. \textbf{Note:} In compound nouns, the stress usually falls on the first element (\emph{MOSquito net, HEALTH centre, TOOTHbrush}).
\end{attention}

\begin{exoresolu}[Form the plurals]
Give the plural of: (a) mosquito net, (b) mother-in-law, (c) grown-up, (d) blood test, (e) passer-by, (f) hand sanitizer, (g) check-up.
\tcblower
\textbf{Solution:}
\begin{enumerate}
\item[(a)] \emph{mosquito nets} --- the main word is \emph{net}, so it takes the \emph{-s}.
\item[(b)] \emph{mothers-in-law} --- the main word is \emph{mother}; \emph{in-law} never changes.
\item[(c)] \emph{grown-ups} --- there is no main noun, so the \emph{-s} goes on the last element.
\item[(d)] \emph{blood tests} --- the main word is \emph{test}.
\item[(e)] \emph{passers-by} --- the main word is \emph{passer}.
\item[(f)] \emph{hand sanitizers} --- the main word is \emph{sanitizer}.
\item[(g)] \emph{check-ups} --- no main noun, plural on last element.
\end{enumerate}
\end{exoresolu}

\courssub{Class discussion: malaria or cholera?}

Use the target vocabulary to compare the prevention of two major diseases in Cameroon. Work in pairs. One student plays a community health worker, the other plays a parent asking questions.

\begin{exemplebox}[Useful sentences for the discussion]
\begin{itemize}
\item Malaria is transmitted by \textbf{mosquitoes}, but cholera is transmitted by \textbf{contaminated water}.
\item To prevent malaria, we sleep under \textbf{treated mosquito nets}; to prevent cholera, we practise \textbf{handwashing} and drink clean water.
\item Both diseases show \textbf{symptoms} like fever or diarrhoea, and both need rapid treatment at a \textbf{health centre}.
\item Good \textbf{sanitation} helps to prevent both diseases.
\item A \textbf{blood test} confirms malaria; a \textbf{stool test} helps confirm cholera.
\item During an \textbf{outbreak}, the \textbf{immune system} of children is most at risk.
\end{itemize}
\end{exemplebox}

\begin{exercice}[Practice: compound nouns and vocabulary]
\begin{enumerate}
\item Complete with a compound noun from the lesson: (a) Before a consultation, patients sit in the \dots\ (b) A \dots\ confirms if you have malaria parasites in your blood. (c) During the COVID-19 outbreak, people used \dots\ to clean their hands when water was not available. (d) The \dots\ will examine you and prescribe treatment.
\item Give the plural of: toothbrush, mother-in-law, waiting room, grown-up, hand sanitizer, check-up, passer-by.
\item Guess from context: ``After the floods, poor \emph{sanitation} caused many cases of typhoid.'' What does \emph{sanitation} mean? Explain your reasoning.
\item Match each disease to its main vector/agent: (a) Malaria --- (b) Cholera --- (c) Typhoid --- (d) Tuberculosis. Options: contaminated water, mosquito, air, contaminated food.
\end{enumerate}
\end{exercice}

\begin{corrige}[Practice exercise]
\begin{enumerate}
\item (a) waiting room; (b) blood test; (c) hand sanitizer; (d) health centre (or doctor/nurse).
\item toothbrushes; mothers-in-law; waiting rooms; grown-ups; hand sanitizers; check-ups; passers-by.
\item The context mentions floods and disease: \emph{sanitation} means the systems that keep water and the environment clean (assainissement). Reasoning: floods damage sanitation systems $\rightarrow$ dirty water $\rightarrow$ typhoid.
\item (a) Malaria $\rightarrow$ mosquito; (b) Cholera $\rightarrow$ contaminated water; (c) Typhoid $\rightarrow$ contaminated food/water; (d) Tuberculosis $\rightarrow$ air.
\end{enumerate}
\end{corrige}

\begin{aretenir}[Key points]
\begin{itemize}
\item A communicable disease spreads along a chain: \textbf{source} $\rightarrow$ \textbf{vector/agent} $\rightarrow$ \textbf{host}. Prevention breaks the chain.
\item Master the health vocabulary: symptom, transmission, vaccine, hygiene, mosquito net, sanitation, outbreak, handwashing, quarantine, immune system, health centre, blood test, hand sanitizer.
\item Read with three strategies: \textbf{skimming} (general idea), \textbf{scanning} (specific information), \textbf{guessing from context}.
\item A \textbf{compound noun} is two or more words with one meaning. The plural goes on the \textbf{main word} (\emph{mothers-in-law}, \emph{health centres}) or on the \textbf{last element} when there is no main word (\emph{grown-ups}, \emph{check-ups}).
\item Spelling: some compounds are one word (\emph{toothbrush}), some two words (\emph{health centre}), some hyphenated (\emph{mother-in-law}). Stress is usually on the first element.
\end{itemize}
\end{aretenir}

In the next section, we will learn how to \textbf{report what people say} --- for example, what a doctor or a health officer tells the population --- using direct and indirect speech.

\coursec{Reporting what people say: direct and indirect speech (Part 1)}

In the previous section we read texts about preventing communicable diseases and learned how to form \cle{compound nouns} and their plurals. Now we move from reading to speaking and writing: we will learn how to report what health workers, teachers or radio announcers say. This skill is essential when you summarise a health talk, write a report for your school magazine or share a sensitisation message with your community.

\courssub{Direct speech vs. indirect speech}

\begin{definition}[Direct and indirect speech]
\textbf{Direct speech} reproduces the exact words of a speaker, enclosed in quotation marks.  
\textbf{Indirect speech} (also called \emph{reported speech}) reformulates the speaker's words without quotation marks, introduced by a \cle{reporting verb} such as \emph{say}, \emph{tell}, \emph{ask}, etc.  
\textit{Français : le discours direct cite les mots exacts (guillemets) ; le discours indirect les rapporte sans guillemets, avec un verbe introducteur.}
\end{definition}

\begin{exemplebox}[Health‑talk example]
Direct: The nurse said, \emph{``I vaccinate 30 children every week.''}  
Indirect: The nurse said (that) she vaccinated 30 children every week.  
\textit{Français : L'infirmière a dit : « Je vaccine 30 enfants chaque semaine. » → L'infirmière a dit qu'elle vaccinait 30 enfants chaque semaine.}
\end{exemplebox}

\courssub{Reporting statements: \emph{say} and \emph{tell}}

When we report a statement we can use two common reporting verbs:

\begin{itemize}
  \item \textbf{say (that) + clause} — no indirect object.  
  \item \textbf{tell + object (that) + clause} — the person who receives the information must be mentioned.
\end{itemize}

\begin{exemplebox}[Say vs. tell]
The doctor \textbf{said} (that) the vaccine was safe.  
The doctor \textbf{told us} (that) the vaccine was safe.  
\textit{Français : Le médecin a dit que le vaccin était sûr. / Le médecin nous a dit que le vaccin était sûr.}
\end{exemplebox}

\begin{attention}[Common mistake: confusing \emph{say} and \emph{tell}]
Never say *\emph{The doctor told that…}* or *\emph{The doctor said us that…}*.  
Correct patterns: \emph{say (that) …}  |  \emph{tell somebody (that) …}.  
\textit{Français : Ne dites pas « told that » ni « said us that ».}
\end{attention}

\begin{aretenir}[The word \emph{that}]
In indirect speech, the conjunction \emph{that} is optional after \emph{said} and \emph{told}. Both forms are correct:  
\emph{The nurse said she was busy.} = \emph{The nurse said that she was busy.}  
\textit{Français : « that » est facultatif après said/told.}
\end{aretenir}

\courssub{Backshift of tenses}

\begin{propriete}[Backshift rule]
When the reporting verb is in the past tense (\emph{said}, \emph{told}, \emph{asked}…), the verb in the reported clause usually moves one step back in time. This is called \emph{backshift}.
\end{propriete}

\begin{popfigure}
\begin{tikzpicture}[node distance=0.5cm, every node/.style={font=\small}]
% Table using matrix
\matrix (tbl) [matrix of nodes, nodes in empty cells,
  column sep=0.2cm, row sep=0.15cm,
  nodes={anchor=center, minimum width=2.8cm, minimum height=0.7cm},
  column 1/.style={nodes={minimum width=2.2cm, fill=popblueL, text=popdark, font=\bfseries}},
  column 2/.style={nodes={fill=popcream}},
  column 3/.style={nodes={fill=popgreenL}},
  row 1/.style={nodes={fill=popblue, text=white, font=\bfseries}}]
{
  \textbf{Direct tense} & \textbf{Indirect tense (backshift)} & \textbf{Example (health / environment)} \\
  Present simple & Past simple & ``I \textbf{plant} trees'' $\rightarrow$ He said he \textbf{planted} trees. \\
  Present continuous & Past continuous & ``We \textbf{are cleaning} the river'' $\rightarrow$ They said they \textbf{were cleaning} the river. \\
  Past simple & Past perfect & ``The nurse \textbf{vaccinated} 120 children'' $\rightarrow$ She said she \textbf{had vaccinated} 120 children. \\
  Present perfect & Past perfect & ``I \textbf{have finished} the garden'' $\rightarrow$ He said he \textbf{had finished} the garden. \\
  Will & Would & ``We \textbf{will protect} the forest'' $\rightarrow$ They said they \textbf{would protect} the forest. \\
  Can & Could & ``You \textbf{can prevent} malaria'' $\rightarrow$ He said you \textbf{could prevent} malaria. \\
};
\draw[popline, thick] (tbl-1-1.north west) rectangle (tbl-7-3.south east);
\foreach \r in {2,...,7} {
  \draw[popline] (tbl-\r-1.north west) -- (tbl-\r-3.north east);
}
\end{tikzpicture}
\legende{Tableau de correspondance des temps (direct $\rightarrow$ indirect) avec exemples santé/environnement.}
\end{popfigure}

\courssub{Changes of time and place expressions}

When we shift from direct to indirect speech, words that indicate \emph{when} and \emph{where} the original words were spoken must also change.

\begin{center}
\begin{tabularx}{\linewidth}{|l|X|X|}\hline
\textbf{Direct} & \textbf{Indirect} & \textbf{Example} \\\hline
now & then & ``I am busy \textbf{now}'' $\rightarrow$ He said he was busy \textbf{then}. \\\hline
today & that day & ``The talk is \textbf{today}'' $\rightarrow$ She said the talk was \textbf{that day}. \\\hline
yesterday & the day before / the previous day & ``We planted trees \textbf{yesterday}'' $\rightarrow$ They said they had planted trees \textbf{the day before}. \\\hline
tomorrow & the next day / the following day & ``The campaign starts \textbf{tomorrow}'' $\rightarrow$ He said the campaign would start \textbf{the next day}. \\\hline
here & there & ``Meet me \textbf{here}'' $\rightarrow$ She told me to meet her \textbf{there}. \\\hline
this / these & that / those & ``\textbf{This} vaccine works'' $\rightarrow$ He said \textbf{that} vaccine worked. \\\hline
\end{tabularx}
\end{center}

\courssub{Changes of pronouns and possessives}

Pronouns adapt to the point of view of the reporter.

\begin{exemplebox}[Pronoun shift]
Direct: ``\textbf{I} will check \textbf{my} blood pressure,'' the nurse said.  
Indirect: The nurse said (that) \textbf{she} would check \textbf{her} blood pressure.  
\textit{Français : « Je vais vérifier ma tension », a dit l'infirmière → L'infirmière a dit qu'elle vérifierait sa tension.}
\end{exemplebox}

\begin{aretenir}[Pronoun reference]
Always ask: \emph{Who is speaking? Who is reporting?}  
\emph{I} $\rightarrow$ he/she (the speaker) | \emph{you} $\rightarrow$ I/we/they (depending on context) | \emph{we} $\rightarrow$ they | \emph{my} $\rightarrow$ his/her | \emph{our} $\rightarrow$ their.
\end{aretenir}

\courssub{Cases with no backshift}

If the reported statement expresses a \cle{general truth} or a fact that is still true at the moment of reporting, the tense does \emph{not} change.

\begin{savaistu}[Scientific facts stay in the present]
The teacher said that \textbf{trees produce} oxygen.  
The doctor explained that \textbf{malaria is transmitted} by mosquitoes.  
\textit{Français : Les vérités scientifiques (généralités) gardent le présent au discours indirect.}
\end{savaistu}

\begin{attention}[Modal verbs that do not backshift]
\emph{should}, \emph{ought to}, \emph{might}, \emph{could}, \emph{would} — these modals already have a past form and do not change.  
Example: ``You \textbf{should wash} your hands.'' $\rightarrow$ He said we \textbf{should wash} our hands.
\end{attention}

\courssub{Method: transforming direct speech into indirect speech (4 steps)}

\begin{methode}[Four‑step transformation]
\begin{enumerate}
  \item \textbf{Identify the reporting verb} (usually \emph{said} or \emph{told}) and note its tense.
  \item \textbf{Adapt pronouns and possessives} to match the reporter's perspective.
  \item \textbf{Apply backshift} to the verb(s) in the reported clause (unless it is a general truth).
  \item \textbf{Change time/place markers} (now $\rightarrow$ then, today $\rightarrow$ that day, etc.).
\end{enumerate}
\end{methode}

\begin{exoresolu}[Reporting a doctor's health talk]
\textbf{Direct speech (doctor at school assembly):} \\
``\emph{Today I will explain how to avoid cholera. You must wash your hands before eating. The water here is not safe. We can all help by boiling water.}''

\textbf{Step‑by‑step transformation:}
\begin{enumerate}
  \item Reporting verb: \emph{The doctor said (that)} — past tense.
  \item Pronouns: \emph{I} $\rightarrow$ \emph{he}; \emph{you} $\rightarrow$ \emph{the students / they}; \emph{we} $\rightarrow$ \emph{they}; \emph{our} $\rightarrow$ \emph{their}.
  \item Backshift: \emph{will explain} $\rightarrow$ \emph{would explain}; \emph{must wash} $\rightarrow$ \emph{had to wash}; \emph{is not safe} $\rightarrow$ \emph{was not safe} (reported as past); \emph{can help} $\rightarrow$ \emph{could help}.
  \item Time/place: \emph{today} $\rightarrow$ \emph{that day}; \emph{here} $\rightarrow$ \emph{there}.
\end{enumerate}

\textbf{Indirect speech:} \\
The doctor said (that) \emph{he would explain how to avoid cholera that day. He added that the students had to wash their hands before eating. He warned that the water there was not safe. He concluded that they could all help by boiling water.}
\end{exoresolu}

\courssub{Guided practice: reporting a school health talk}

\begin{exercice}[Guided practice]
Report the following statements made by the school nurse during a health talk. Use \emph{said} or \emph{told us} and apply all necessary changes.
\begin{enumerate}
  \item ``I \textbf{am checking} the vaccination records \textbf{now}.''
  \item ``You \textbf{should drink} clean water \textbf{every day}.''
  \item ``Last week we \textbf{distributed} mosquito nets \textbf{here}.''
  \item ``The campaign \textbf{will continue} \textbf{tomorrow}.''
  \item ``Malaria \textbf{is caused} by a parasite.''
\end{enumerate}
\end{exercice}

\begin{corrige}[Exercice n°1]
\begin{enumerate}
  \item The nurse said (that) she \textbf{was checking} the vaccination records \textbf{then}.
  \item The nurse told us (that) we \textbf{should drink} clean water \textbf{every day}. (no backshift for \emph{should}; time marker unchanged)
  \item The nurse said (that) they \textbf{had distributed} mosquito nets \textbf{there} \textbf{the week before}.
  \item The nurse said (that) the campaign \textbf{would continue} \textbf{the next day}.
  \item The nurse said (that) malaria \textbf{is caused} by a parasite. (general truth — no backshift)
\end{enumerate}
\end{corrige}

\courssub{Mini‑project: reporting a radio sensitisation message on tree planting}

\begin{exercice}[Mini‑project]
Listen to (or read) the following five sentences from a radio message about tree planting. Write a short paragraph reporting the message in indirect speech. Use a variety of reporting verbs (\emph{said}, \emph{announced}, \emph{explained}, \emph{urged}, \emph{reminded}).

\begin{enumerate}
  \item ``We \textbf{are planting} 500 trees \textbf{this month}.''
  \item ``Every citizen \textbf{must protect} the young trees.''
  \item ``The forest \textbf{provides} clean air and water.''
  \item ``Next week we \textbf{will organise} a community planting day.''
  \item ``You \textbf{can help} by watering the seedlings \textbf{every evening}.''
\end{enumerate}
\end{exercice}

\begin{corrige}[Mini‑project]
The radio announcer \textbf{announced} that they \textbf{were planting} 500 trees \textbf{that month}. He \textbf{explained} that every citizen \textbf{had to protect} the young trees. He \textbf{added} that the forest \textbf{provides} clean air and water (general truth). He \textbf{said} that they \textbf{would organise} a community planting day \textbf{the following week}. Finally, he \textbf{urged} listeners to help by watering the seedlings \textbf{every evening}.
\end{corrige}

\begin{popfigure}
\begin{tikzpicture}[>=Stealth, node distance=1.2cm, every node/.style={font=\small}]
% Nodes for direct sentence
\node[draw, rounded corners, fill=poppinkL, minimum width=5cm, minimum height=1.2cm, align=center] (direct){
  \textbf{Direct speech}\\
  ``I vaccinated 120 children yesterday,'' the nurse said.
};
% Arrow 1: reporting verb
\node[draw, rounded corners, fill=poporangeL, right=1.5cm of direct, minimum width=3.5cm, minimum height=1cm] (verb) {Reporting verb: \emph{said} (past)};
% Arrow 2: pronoun change
\node[draw, rounded corners, fill=popgreenL, below=1.2cm of verb, minimum width=3.5cm, minimum height=1cm] (pronoun) {Pronoun: \emph{I} $\rightarrow$ \emph{she}};
% Arrow 3: backshift
\node[draw, rounded corners, fill=popblueL, below=1.2cm of pronoun, minimum width=3.5cm, minimum height=1cm] (tense) {Backshift: \emph{vaccinated} $\rightarrow$ \emph{had vaccinated}};
% Arrow 4: time marker
\node[draw, rounded corners, fill=popgoldL, below=1.2cm of tense, minimum width=3.5cm, minimum height=1cm] (time) {Time: \emph{yesterday} $\rightarrow$ \emph{the day before}};
% Result node
\node[draw, rounded corners, fill=popcream, right=1.5cm of tense, minimum width=5.5cm, minimum height=1.2cm, align=center] (indirect){
  \textbf{Indirect speech}\\
  The nurse said (that) she \textbf{had vaccinated} 120 children \textbf{the day before}.
};

% Arrows
\draw[->, thick, popdark] (direct.east) -- (verb.west) node[midway, above] {1};
\draw[->, thick, popdark] (verb.south) -- (pronoun.north) node[midway, right] {2};
\draw[->, thick, popdark] (pronoun.south) -- (tense.north) node[midway, right] {3};
\draw[->, thick, popdark] (tense.south) -- (time.north) node[midway, right] {4};
\draw[->, thick, popdark] (tense.east) -- (indirect.west) node[midway, above] {Result};
\end{tikzpicture}
\legende{Schéma de transformation d'une phrase directe en phrase rapportée (pronoms, temps, marqueurs spatio‑temporels).}
\end{popfigure}

\courssub{Summary checklist}

\begin{itemize}
  \item Identify the reporting verb and its tense.
  \item Change pronouns/possessives to the reporter's perspective.
  \item Apply backshift (except for general truths and modals like \emph{should}, \emph{might}, \emph{could}, \emph{would}).
  \item Adjust time and place expressions.
  \item Choose \emph{say} (no object) or \emph{tell} + object.
  \item Remember \emph{that} is optional.
\end{itemize}

\begin{aretenir}[Exam tip — Probatoire / Baccalauréat technique]
In the exam, you may be asked to:
\begin{itemize}
  \item Rewrite a short dialogue or health talk in reported speech.
  \item Fill in blanks with the correct backshifted verb form.
  \item Choose between \emph{say} and \emph{tell} in a given sentence.
  \item Identify errors in reported speech sentences.
\end{itemize}
Practise with authentic health/environment messages (radio, posters, school talks) to build confidence.
\end{aretenir}

With these tools you can now report health talks, environmental campaigns or any spoken message accurately — a key skill for the Probatoire and for real‑life communication in your future technical career.

\coursec{Reporting questions and commands: indirect speech (Part 2)}

In the previous section \emph{Reporting what people say: direct and indirect speech (Part 1)}, we mastered the art of reporting \textbf{statements} (declarative sentences). We learned how to shift tenses back (\cle{backshift}), change pronouns and adjust time/place references. Now, we extend these skills to \textbf{questions} and \textbf{commands/requests}. These are essential when you report interviews, medical advice, safety instructions or conversations with tourists and health workers — typical situations in your technical training.

\courssub{Reporting Yes/No Questions}

When we report a closed question (a question expecting \emph{Yes} or \emph{No}), we do not keep the auxiliary verb at the beginning. We transform the question into a \textbf{statement word order} (Subject + Verb) and introduce it with \textbf{if} or \textbf{whether}.

\begin{definition}[Reported Yes/No Question]
A \cle{reported yes/no question} is an indirect question introduced by \textbf{ask + if/whether} followed by \textbf{statement word order} (Subject + Verb) and \textbf{no question mark}. The tense usually shifts back (backshift).
\end{definition}

\begin{propriete}[Structure: Yes/No Questions]
\textbf{Direct:} \emph{``Did you visit the botanical garden?''} \\
\textbf{Reported:} \emph{He asked \textbf{if/whether} I \textbf{had visited} the botanical garden.}
\begin{itemize}
    \item \textbf{Introductory verb:} \cle{ask} (most common), \cle{inquire}, \cle{wonder}.
    \item \textbf{Linker:} \cle{if} or \cle{whether} (\emph{whether} is slightly more formal).
    \item \textbf{Word order:} \textbf{Subject + Verb} (NO inversion: \emph{not} ``asked if had I visited'').
    \item \textbf{Punctuation:} Full stop (.), \textbf{no question mark}.
    \item \textbf{Backshift:} Present Simple $\to$ Past Simple; Past Simple/Present Perfect $\to$ Past Perfect; \emph{will} $\to$ \emph{would}; \emph{can} $\to$ \emph{could}.
\end{itemize}
\end{propriete}

\begin{exemplebox}[Tourist Site Context]
\textbf{Direct:} The guide asked, ``\emph{Do you have a map of the park?}'' \\
\textbf{Reported:} The guide asked \textbf{if} I \textbf{had} a map of the park. \\[4pt]
\textbf{Direct:} The tourist asked, ``\emph{Will the bus leave at noon?}'' \\
\textbf{Reported:} The tourist asked \textbf{whether} the bus \textbf{would leave} at noon.
\end{exemplebox}

\begin{attention}[Common Error: Keeping Inversion]
\textbf{Wrong:} He asked \emph{if did I see} the animals. \\
\textbf{Right:} He asked \emph{if I saw/had seen} the animals. \\
\textbf{Reason:} In reported questions, the subject \textbf{always} comes before the verb. The auxiliary \emph{do/does/did} disappears (except in negative questions: \emph{if I didn't see}).
\end{attention}

\begin{exoresolu}[Reporting a Health Survey Question]
\textbf{Task:} Report the nurse's question: ``\emph{Did you sleep under a mosquito net last night?}'' \\[4pt]
\textbf{Solution:}
\begin{enumerate}
    \item Identify type: Yes/No question (auxiliary \emph{Did}).
    \item Choose linker: \textbf{if} (or \emph{whether}).
    \item Apply backshift: \emph{Did sleep} (Past Simple) $\to$ \textbf{had slept} (Past Perfect).
    \item Adjust time reference: \emph{last night} $\to$ \textbf{the night before} (or \emph{the previous night}).
    \item Adjust pronoun: \emph{you} $\to$ \textbf{I} (reporter is the patient).
    \item Word order: Subject (\emph{I}) + Verb (\emph{had slept}).
\end{enumerate}
\textbf{Answer:} The nurse asked \textbf{if I had slept under a mosquito net the night before}.
\end{exoresolu}

\courssub{Reporting Wh-Questions (Open Questions)}

Open questions start with a \textbf{question word} (\emph{who, what, where, when, why, how, how much/many}). The good news: you \textbf{keep the question word}. The rule for word order remains the same: \textbf{Statement order (Subject + Verb)}.

\begin{propriete}[Structure: Wh-Questions]
\textbf{Direct:} \emph{``Where is the park entrance?''} \\
\textbf{Reported:} \emph{The tourist asked \textbf{where} the park entrance \textbf{was}.}
\begin{itemize}
    \item \textbf{Linker:} The original \cle{wh-word} (\emph{where, when, how, what, etc.}).
    \item \textbf{Word order:} \textbf{Wh-word + Subject + Verb} (NO inversion: \emph{not} ``where was the entrance'').
    \item \textbf{Backshift:} Applies normally (\emph{is} $\to$ \emph{was}; \emph{are} $\to$ \emph{were}; \emph{have} $\to$ \emph{had}).
    \item \textbf{Subject questions (Who/What + Verb):} No change in structure, just backshift. \emph{``Who called you?''} $\to$ He asked who \textbf{had called} me.
\end{itemize}
\end{propriete}

\begin{exemplebox}[Environment Context]
\textbf{Direct:} The activist asked, ``\emph{How do you plant these seedlings?}'' \\
\textbf{Reported:} The activist asked \textbf{how} we \textbf{planted/had planted} those seedlings. \\[4pt]
\textbf{Direct:} The ranger asked, ``\emph{Why are you cutting that tree?}'' \\
\textbf{Reported:} The ranger asked \textbf{why} I \textbf{was cutting} that tree.
\end{exemplebox}

\begin{attention}[Trap: Inversion in Wh-Questions]
\textbf{Wrong:} She asked \emph{where was the clinic}. \\
\textbf{Right:} She asked \emph{where the clinic was}. \\
\textbf{Tip:} Say the reported clause as a statement: \emph{``The clinic was there''} $\to$ \emph{where the clinic was}.
\end{attention}

\begin{exoresolu}[Reporting a Gardening Question]
\textbf{Task:} Report: ``\emph{When will the trees be ready for harvest?}'' asked the farmer. \\[4pt]
\textbf{Solution:}
\begin{enumerate}
    \item Wh-word: \textbf{When}.
    \item Subject: \textbf{the trees}.
    \item Verb backshift: \emph{will be} $\to$ \textbf{would be}.
    \item Object/Complement: \emph{ready for harvest}.
\end{enumerate}
\textbf{Answer:} The farmer asked \textbf{when the trees would be ready for harvest}.
\end{exoresolu}

\courssub{Reporting Commands, Requests and Advice}

Commands (imperatives) are reported using an \textbf{infinitive structure}. We choose the introductory verb based on the \textbf{force} or \textbf{intention} of the speaker.

\begin{methode}[Steps to Report a Command/Request]
\begin{enumerate}
    \item \textbf{Identify the imperative verb} (base form: \emph{Go, Stop, Plant, Don't drink}).
    \item \textbf{Choose the reporting verb}: \cle{tell} (strong command), \cle{ask} (request), \cle{order} (authority), \cle{advise} (suggestion for good), \cle{warn} (danger), \cle{encourage} (support), \cle{invite} (social), \cle{beg} (desperate).
    \item \textbf{Add the object} (person receiving the command): \emph{me, him, her, us, them, the tourists, the workers}.
    \item \textbf{Use \textbf{to-infinitive}} (positive) or \textbf{not to-infinitive} (negative).
\end{enumerate}
\end{methode}

\begin{definition}[Reported Command / Request]
A \cle{reported command} uses the structure: \textbf{Reporting Verb + Object + (not) + to-infinitive}. It has no question mark and no subject in the reported clause (the object becomes the subject of the infinitive).
\end{definition}

\begin{exemplebox}[Health \& Environment Contexts]
\begin{tabularx}{\linewidth}{@{}lX@{}}
\textbf{Positive Command:} & \emph{``Plant more trees,'' said the minister.} $\to$ The minister \textbf{told/encouraged} the farmers \textbf{to plant} more trees. \\
\textbf{Request:} & \emph{``Please show me your ticket,'' said the guard.} $\to$ The guard \textbf{asked} the visitor \textbf{to show} his ticket. \\
\textbf{Medical Advice:} & \emph{``Drink boiled water,'' said the doctor.} $\to$ The doctor \textbf{advised} the patient \textbf{to drink} boiled water. \\
\textbf{Negative Command:} & \emph{``Don't swim in the river,'' said the guide.} $\to$ The guide \textbf{warned/told} us \textbf{not to swim} in the river. \\
\textbf{Negative Request:} & \emph{``Please don't litter,'' said the ranger.} $\to$ The ranger \textbf{asked} the tourists \textbf{not to litter}.
\end{tabularx}
\end{exemplebox}

\begin{attention}[Negative Commands: \textbf{not to} vs \textbf{to not}]
\textbf{Standard:} \emph{He told me \textbf{not to go}.} \\
\textbf{Avoid:} \emph{He told me \textbf{to not go}.} (Split infinitive, less formal, often considered an error in exams). \\
\textbf{Verb Choice Matters:} \emph{``Don't touch the chemical.''} $\to$ He \textbf{warned} me not to touch it (danger). $\neq$ He \textbf{advised} me not to touch it (good practice).
\end{attention}

\begin{exoresolu}[Reporting Safety Instructions]
\textbf{Task:} Report the following instructions given by a site supervisor in Kribi. \\
1. ``\emph{Wear your helmets at all times.}'' (to the workers) \\
2. ``\emph{Don't remove your gloves.}'' (to the workers) \\
3. ``\emph{Please report any accident immediately.}'' (to the team leader) \\[4pt]
\textbf{Solution:}
\begin{enumerate}
    \item The supervisor \textbf{told/ordered} the workers \textbf{to wear} their helmets at all times.
    \item The supervisor \textbf{told/warned} the workers \textbf{not to remove} their gloves.
    \item The supervisor \textbf{asked/advised} the team leader \textbf{to report} any accident immediately.
\end{enumerate}
\end{exoresolu}

\courssub{Choosing the Right Reporting Verb: Nuance and Precision}

At Probatoire level, using \emph{said} or \emph{told} for everything limits your score. Precise verbs show \cle{pragmatic competence} — you understand the \textbf{speaker's intention}.

\begin{propriete}[Reporting Verbs \& Intentions]
\begin{center}
\begin{tabularx}{\linewidth}{|l|X|l|}\hline
\textbf{Verb} & \textbf{Intention / Context} & \textbf{Structure} \\ \hline
\textbf{Advise} & Recommending what is good/healthy & advise + obj + to-inf \\ \hline
\textbf{Warn} & Alerting about danger/risk & warn + obj + (not) to-inf \\ \hline
\textbf{Encourage} & Giving confidence/support & encourage + obj + to-inf \\ \hline
\textbf{Invite} & Asking someone socially & invite + obj + to-inf \\ \hline
\textbf{Beg} & Desperate/emotional request & beg + obj + to-inf \\ \hline
\textbf{Promise} & Committing to future action & promise + (obj) + to-inf \\ \hline
\textbf{Order} & Authority, no choice & order + obj + to-inf \\ \hline
\textbf{Ask} & Neutral request / question & ask + obj + to-inf / if/wh \\ \hline
\textbf{Tell} & Neutral command / statement & tell + obj + to-inf / that-clause \\ \hline
\end{tabularx}
\end{center}
\end{propriete}

\begin{popfigure}
\begin{tikzpicture}[
    node distance=1.2cm and 1.5cm,
    start/.style={rectangle, rounded corners, draw=popblue, thick, fill=popblueL, text width=3.5cm, align=center, font=\bfseries},
    decision/.style={diamond, aspect=2, draw=poporange, thick, fill=poporangeL, text width=3cm, align=center, font=\small},
    action/.style={rectangle, rounded corners, draw=popgreen, thick, fill=popgreenL, text width=4cm, align=center, font=\small},
    end/.style={ellipse, draw=poppurple, thick, fill=poppurpleL, text width=3.5cm, align=center, font=\bfseries\small},
    arrow/.style={->, thick, draw=popdark}
]
\node[start] (start) {Direct Speech Act};
\node[decision, below of=start] (type) {What is the sentence type?};
\node[action, below left=1.5cm and 2cm of type] (stmt) {Statement $\to$ \textbf{say/tell} + that-clause};
\node[action, below of=type] (quest) {Question};
\node[action, below right=1.5cm and 2cm of type] (cmd) {Command / Request};

\node[decision, below=1.5cm of quest] (qtype) {Yes/No or Wh-?};
\node[action, below left=1.2cm and 1.5cm of qtype] (yn) {Yes/No $\to$ \textbf{ask if/whether} + S+V};
\node[action, below right=1.2cm and 1.5cm of qtype] (wh) {Wh- $\to$ \textbf{ask wh-word} + S+V};

\node[decision, below=1.5cm of cmd] (intent) {What is the speaker's intention?};
\node[action, below left=1.2cm and 1.8cm of intent] (warn) {Danger/Risk $\to$ \textbf{warn} + obj + not to};
\node[action, below left=0.5cm and 0.5cm of intent] (adv) {Health/Best practice $\to$ \textbf{advise} + obj + to};
\node[action, below=0.2cm of intent] (enc) {Support/Motivate $\to$ \textbf{encourage} + obj + to};
\node[action, below right=0.5cm and 0.5cm of intent] (ord) {Authority $\to$ \textbf{order/tell} + obj + to};
\node[action, below right=1.2cm and 1.8cm of intent] (req) {Polite/Social $\to$ \textbf{ask/invite} + obj + to};

\draw[arrow] (start) -- (type);
\draw[arrow] (type) -| node[near start, above, sloped, font=\tiny] {Statement} (stmt);
\draw[arrow] (type) -- node[left, font=\tiny] {Question} (quest);
\draw[arrow] (type) -| node[near start, above, sloped, font=\tiny] {Imperative} (cmd);
\draw[arrow] (quest) -- (qtype);
\draw[arrow] (qtype) -| node[near start, above, font=\tiny] {Yes/No} (yn);
\draw[arrow] (qtype) -| node[near start, above, font=\tiny] {Wh-} (wh);
\draw[arrow] (cmd) -- (intent);
\draw[arrow] (intent) -| node[near start, above, font=\tiny] {Danger} (warn);
\draw[arrow] (intent) -| node[near start, above, font=\tiny] {Advice} (adv);
\draw[arrow] (intent) -- node[left, font=\tiny] {Support} (enc);
\draw[arrow] (intent) -| node[near start, above, font=\tiny] {Authority} (ord);
\draw[arrow] (intent) -| node[near start, above, font=\tiny] {Polite} (req);
\end{tikzpicture}
\legende{Decision tree: Choosing the correct reporting structure and verb based on sentence type and speaker intention.}
\end{popfigure}

\begin{savaistu}[Exam Tip: Pragmatic Variety]
In the \cle{Probatoire} and \cle{Baccalauréat technique} writing tasks (compositions, reports), examiners reward \textbf{lexical variety}. Instead of writing ``He said I must...'' or ``He told me to...'' five times, use: \emph{The doctor \textbf{advised} me to...}; \emph{The guide \textbf{warned} us not to...}; \emph{The activist \textbf{encouraged} the youth to...}; \emph{The minister \textbf{promised} to...}. This demonstrates B2-level control of functional language.
\end{savaistu}

\courssub{Consolidation: Mixed Reporting Practice}

Real conversations mix statements, questions and commands. You must switch structures fluidly.

\begin{exercice}[Journalist Interview in Douala]
\textbf{Context:} A journalist interviews \textbf{Mr. Ngu}, an environmental activist in Douala, about a mangrove replanting project. Report the \textbf{entire interview} in indirect speech. Use a variety of reporting verbs (\emph{said, told, asked, advised, warned, encouraged, promised}).

\textbf{Interview Transcript:}
\begin{enumerate}
    \item \textbf{Journalist:} ``Good morning, Mr. Ngu. \textbf{What is the main goal of this project?}''
    \item \textbf{Mr. Ngu:} ``\textbf{We want} to restore 50 hectares of mangrove. \textbf{It protects} the coast from erosion.''
    \item \textbf{Journalist:} ``\textbf{Did you receive} government funding?'' 
    \item \textbf{Mr. Ngu:} ``\textbf{Yes, we did.} The Ministry \textbf{gave} us a grant last year.''
    \item \textbf{Journalist:} ``\textbf{How can} local people help?''
    \item \textbf{Mr. Ngu:} ``\textbf{Plant} seedlings with us next Saturday. \textbf{Don't throw} plastic waste in the water.''
    \item \textbf{Journalist:} ``\textbf{Will you organize} more events?''
    \item \textbf{Mr. Ngu:} ``\textbf{Yes, I will.} I \textbf{promise} to organize monthly clean-ups.''
\end{enumerate}
\end{exercice}

\begin{corrige}[Exercice n°1 : Journalist Interview in Douala]
\textbf{Reported Version:}
\begin{enumerate}
    \item The journalist \textbf{greeted} Mr. Ngu and \textbf{asked} \textbf{what the main goal of the project was}.
    \item Mr. Ngu \textbf{replied/stated} \textbf{that they wanted} to restore 50 hectares of mangrove and \textbf{added} \textbf{that it protected} the coast from erosion.
    \item The journalist \textbf{asked} \textbf{if/whether they had received} government funding.
    \item Mr. Ngu \textbf{confirmed} \textbf{that they had} and \textbf{explained} \textbf{that the Ministry had given} them a grant the previous year.
    \item The journalist \textbf{asked} \textbf{how local people could help}.
    \item Mr. Ngu \textbf{encouraged/invited} them \textbf{to plant} seedlings with him the following Saturday and \textbf{warned/advised} them \textbf{not to throw} plastic waste in the water.
    \item The journalist \textbf{asked} \textbf{if/whether he would organize} more events.
    \item Mr. Ngu \textbf{confirmed} \textbf{that he would} and \textbf{promised} \textbf{to organize} monthly clean-ups.
\end{enumerate}
\textbf{Key Checks:} Backshift (want $\to$ wanted, protect $\to$ protected, give $\to$ had given, will $\to$ would). Pronouns (we $\to$ they, us $\to$ them, I $\to$ he). Time/Place (last year $\to$ previous year, next Saturday $\to$ following Saturday). Word order in Qs (What is $\to$ what ... was; How can $\to$ how ... could). Neg command $\to$ not to throw. Promise + to-inf.
\end{corrige}

\courssub{Role Play: Journalist \& Environmental Activist}

\begin{methode}[Role Play Procedure]
\begin{enumerate}
    \item \textbf{Pair Work:} Student A = Journalist (prepares 5-6 questions: Yes/No, Wh-, requests for advice). Student B = Activist (prepares facts: goals, funding, needs, future plans, warnings).
    \item \textbf{Perform:} Conduct a 2-minute interview. \textbf{Record} or take notes.
    \item \textbf{Switch Roles:} Repeat.
    \item \textbf{Write Report:} Individually, write a 100-word news report in \textbf{indirect speech} summarizing the interview. Use at least \textbf{4 different reporting verbs}.
    \item \textbf{Peer Assessment:} Exchange reports. Use the \textbf{Correction Grid} below.
\end{enumerate}
\end{methode}

\begin{definition}[Peer Assessment Grid for Reported Speech]
\textbf{Instructions:} Tick (✓) the box that best describes your partner's work. Total: 10 marks.
\begin{center}
\begin{tabularx}{\linewidth}{|l|X|c|c|c|}\hline
\textbf{Criteria} & \textbf{Descriptors} & \textbf{2 pts} & \textbf{1 pt} & \textbf{0 pt} \\ \hline
\textbf{Backshift Accuracy} & Tenses correctly shifted (Pres$\to$Past, Past$\to$Past Perf, Will$\to$Would) & All correct & 1-2 errors & >2 errors \\ \hline
\textbf{Question Structures} & Yes/No: ask if/whether + S+V; Wh-: ask wh-word + S+V (no inversion) & Perfect & Minor inversion error & Major errors \\ \hline
\textbf{Command Structures} & Verb + Obj + (not) to-inf; Correct negative form (not to) & Perfect & 1 error (split inf / missing obj) & Wrong structure \\ \hline
\textbf{Reporting Verb Variety} & Uses $\ge$ 4 distinct verbs (advise, warn, encourage, promise, etc.) appropriately & $\ge$ 4 varied & 2-3 varied (mostly say/tell) & Only say/tell \\ \hline
\textbf{Pronouns / Time Ref.} & Pronouns (I$\to$he, we$\to$they) \& Time (now$\to$then, today$\to$that day) correct & All correct & 1-2 slips & Confusing \\ \hline
\end{tabularx}
\end{center}
\textbf{Feedback:} Write one \emph{strength} and one \emph{target for improvement} for your partner.
\end{definition}

\courssub{Summary: Three Types of Reported Speech at a Glance}

\begin{popfigure}
\begin{tikzpicture}
\node[inner sep=0pt, align=center] (table){
\begin{tabularx}{\linewidth}{|>{\bfseries}l|X|X|X|}\hline
\rowcolor{popblueL}
\textbf{Type} & \textbf{Direct Speech Example} & \textbf{Reported Structure} & \textbf{Reported Example} \\ \hline
\textbf{1. Statement} & \emph{``We protect the forest.''} & \textbf{say/tell/explain/add} + \textbf{that} + S + V (backshifted) & He \textbf{said (that) they protected} the forest. \\ \hline
\textbf{2. Yes/No Question} & \emph{``Do you recycle?''} & \textbf{ask/inquire} + \textbf{if/whether} + S + V (backshifted) & She \textbf{asked if I recycled/had recycled}. \\ \hline
\textbf{3. Wh-Question} & \emph{``Where is the nursery?''} & \textbf{ask} + \textbf{wh-word} + S + V (backshifted) & He \textbf{asked where the nursery was}. \\ \hline
\textbf{4. Command} & \emph{``Plant this tree.''} & \textbf{tell/order/advise/encourage/warn} + Obj + \textbf{to} + V\textsubscript{base} & The ranger \textbf{told us to plant} that tree. \\ \hline
\textbf{5. Negative Command} & \emph{``Don't cut the trees.''} & \textbf{tell/warn/advise} + Obj + \textbf{not to} + V\textsubscript{base} & He \textbf{warned us not to cut} the trees. \\ \hline
\textbf{6. Request} & \emph{``Please help us.''} & \textbf{ask/invite/beg} + Obj + \textbf{to} + V\textsubscript{base} & She \textbf{asked me to help} them. \\ \hline
\end{tabularx}
};
\legende{Summary table of Reported Speech structures covered in Part 1 (Statements) and Part 2 (Questions, Commands, Requests).}
\end{tikzpicture}
\end{popfigure}

\begin{aretenir}[Key Takeaways for Exams]
\begin{itemize}
    \item \textbf{Questions:} \textbf{No inversion, no question mark, no do/does/did}. Use \textbf{if/whether} (Yes/No) or \textbf{Wh-word} (Open). Backshift the tense.
    \item \textbf{Commands/Requests:} \textbf{Verb + Object + (not) to-infinitive}. Choose verb for \textbf{meaning}: \emph{warn} (danger), \emph{advise} (health), \emph{encourage} (support), \emph{promise} (commitment).
    \item \textbf{Mixed Texts:} Identify each speech act (Statement? Question? Command?) and apply the correct formula instantly.
    \item \textbf{Probatoire Writing:} In compositions (health, environment topics), use reported speech to cite experts, witnesses or authorities: \emph{The doctor \textbf{advised} the population to vaccinate their children. The expert \textbf{warned} that pollution would increase.}
\end{itemize}
\end{aretenir}

\coursec{Writing and integration: communicating on health issues}

In the previous section, you learnt how to report questions and commands in indirect speech: \emph{The nurse asked me if I washed my hands regularly} or \emph{The doctor told me to drink clean water}. Now it is time to use everything you have studied in this lesson — \cle{used to}, \cle{tag questions}, \cle{compound nouns} and \cle{reported speech} — in real communication. In this final section, you will learn how to \textbf{write clear, well-organised compositions} about health issues and the prevention of communicable diseases, and you will combine all your skills in an integration project: a \emph{School Health and Environment Day}.

\courssub{What is a composition?}

Imagine the situation: cholera has appeared in a quarter of Yaoundé. The principal of your school asks you to write a text for the students. What do you write? A poem? A shopping list? No — you write a \cle{composition}: an organised text with a clear message for a clear audience.

\begin{definition}[Composition]
A \cle{composition} is an organised piece of writing made of three parts: an \textbf{introduction} (which presents the issue), a \textbf{body} or development (which gives causes, consequences and prevention measures), and a \textbf{conclusion} (which often ends with a call to action). A good composition is always adapted to its \textbf{reader} (\emph{destinataire}) and to its \textbf{communicative objective} (\emph{objectif de communication}).
\end{definition}

In this unit, you will meet four common text types:

\begin{center}
\begin{tabularx}{\linewidth}{|l|X|X|}
\hline
\rowcolor{popblueL} \textbf{Text type} & \textbf{Reader and objective} & \textbf{Style} \\
\hline
Sensitization poster & The whole school; to inform quickly and push to action & Short sentences, slogans, imperatives \\
\hline
Letter to a friend & One friend; to advise or share experience & Informal, warm, personal \\
\hline
Article for the school magazine & Students and teachers; to inform and explain & Semi-formal, factual, organised \\
\hline
Dialogue patient / health worker & A reader or an audience; to show good practice & Spoken style, questions and answers \\
\hline
\end{tabularx}
\end{center}

\courssub{The structure of a good health composition}

Whatever the text type, the skeleton is the same. Let us build it step by step, like a house: foundations (introduction), walls and rooms (body), roof (conclusion).

\textbf{1. The introduction} presents the issue and catches the reader's attention. You can start with a question, a fact or a short anecdote: \emph{Did you know that malaria still kills thousands of Cameroonians every year?}

\textbf{2. The body} develops two or three ideas, one paragraph per idea:
\begin{itemize}
\item \textbf{Causes}: \emph{Malaria is transmitted by the bite of an infected mosquito.}
\item \textbf{Consequences}: \emph{As a result, many pupils miss school during the rainy season.}
\item \textbf{Prevention measures}: \emph{Moreover, sleeping under a mosquito net is cheap and effective.}
\end{itemize}

\textbf{3. The conclusion} summarises and calls the reader to action: \emph{Therefore, let us all sleep under treated mosquito nets. Prevention is better than cure!}

\begin{popfigure}
\begin{center}
\begin{tikzpicture}[
box/.style={draw=popblue, thick, rounded corners, fill=popblueL, text width=3.4cm, align=center, minimum height=1.6cm, font=\small},
lnk/.style={draw=poporange, thick, rounded corners, fill=poporangeL, text width=4.6cm, align=left, font=\scriptsize},
arr/.style={-{Stealth[length=3mm]}, very thick, popdark}]
\node[box, align=center] (intro) at (0,4.4){\textbf{INTRODUCTION}\\ Present the issue\\ Catch attention};
\node[box, align=center] (body) at (0,1.8){\textbf{BODY}\\ Causes $\rightarrow$ consequences\\ $\rightarrow$ prevention};
\node[box, align=center] (concl) at (0,-0.8){\textbf{CONCLUSION}\\ Summary + call to action};
\node[lnk, align=center] (l1) at (5.6,4.4){\textbf{Openers:}\\ Did you know that...?\\ Nowadays, ...\\ It is true that..., isn't it?};
\node[lnk, align=center] (l2) at (5.6,1.8){\textbf{Linking words:}\\ First, ... / Moreover, ...\\ However, ... / As a result, ...};
\node[lnk, align=center] (l3) at (5.6,-0.8){\textbf{Concluders:}\\ Therefore, ...\\ In conclusion, ...\\ Let us all ...};
\draw[arr] (intro) -- (body);
\draw[arr] (body) -- (concl);
\draw[arr, poporange] (l1.west) -- (intro.east);
\draw[arr, poporange] (l2.west) -- (body.east);
\draw[arr, poporange] (l3.west) -- (concl.east);
\end{tikzpicture}
\end{center}
\legende{Structure of a health composition in three parts, with the linking words (\emph{mots de liaison}) associated to each part.}
\end{popfigure}

\begin{aretenir}[Key linking words for health writing]
\begin{itemize}
\item \cle{First,} / \cle{Moreover,} — to add ideas (\emph{pour ajouter une idée}).
\item \cle{However,} — to contrast ideas (\emph{pour opposer}): \emph{Vaccines are free; however, many people do not go to the health centre.}
\item \cle{Therefore,} / \cle{As a result,} — to show a consequence (\emph{pour montrer une conséquence}).
\item \cle{In conclusion,} — to close the text (\emph{pour conclure}).
\end{itemize}
\end{aretenir}

\courssub{Reusing the lesson's grammar in your writing}

A Probatoire composition must show the examiner that you master the grammar of the unit. Here is how to integrate each point naturally.

\textbf{Used to — for past habits.} Compare the past and the present: \emph{People in our village \textbf{used to drink} water directly from the stream, but now they boil it first.} / \emph{We \textbf{didn't use to wash} our hands before meals.}

\textbf{Tag questions — to engage the reader.} A tag question involves your audience: \emph{Prevention is better than cure, \textbf{isn't it}?} / \emph{Nobody wants to be sick, \textbf{do they}?}

\textbf{Compound nouns — for precision.} Precise vocabulary makes your text credible: \emph{mosquito net}, \emph{health centre}, \emph{hand washing}, \emph{clean water supply}, \emph{bed nets}. Remember the plural rule: usually only the second noun takes the plural \emph{s}: \emph{health centre}s, \emph{mosquito net}s.

\textbf{Reported speech — to quote authorities.} Quoting a health worker makes your argument stronger: \emph{The nurse said that vaccination \textbf{was} free for all children.} / \emph{The doctor advised us \textbf{to sleep} under mosquito nets.}

\begin{exemplebox}[A model paragraph using all four points]
\emph{In the past, the pupils of our school \textbf{used to throw} rubbish everywhere, and many of them fell sick. Last month, the head nurse of the district hospital visited us. She \textbf{said that} dirty environments \textbf{caused} many communicable diseases, and she \textbf{advised us to keep} our classrooms clean. Since then, we clean the \textbf{school grounds} every Friday and we wash our hands at the new \textbf{water point}. Our school is much healthier now, \textbf{isn't it}?}
\end{exemplebox}

\begin{attention}[Do not translate word for word from French!]
The most frequent error is to think in French and translate literally: \emph{*I have paludism} (for \emph{j'ai le paludisme}) instead of \emph{I have malaria}; \emph{*the maladies diarrhéiques} instead of \emph{diarrhoeal diseases}. Think directly in \textbf{simple, correct English structures}: short sentences, subject + verb + complement. A simple correct sentence always scores more than a long incorrect one.
\end{attention}

\courssub{The writing process: five steps}

Good writers never write their final text immediately. They follow a process.

\begin{methode}[Writing a composition in five steps]
\begin{enumerate}
\item \textbf{Brainstorm} (\emph{remue-méninges}): list all your ideas, vocabulary and facts about the topic, in any order.
\item \textbf{Plan}: select the best ideas and organise them — introduction, two or three body paragraphs, conclusion.
\item \textbf{Draft} (\emph{brouillon}): write a first version without worrying too much about mistakes; use your linking words.
\item \textbf{Peer review} (\emph{correction par un camarade}): exchange your draft with a classmate; check grammar, vocabulary and organisation with the checklist below.
\item \textbf{Final copy}: rewrite neatly, correcting all the errors found.
\end{enumerate}
\end{methode}

\begin{exoresolu}[Writing a short article for the school magazine]
\textbf{Task:} Your school magazine asks you to write an article (150–200 words) entitled \emph{``Stop cholera in our school!''}. Present the issue, its causes and consequences, prevention measures, and end with a call to action.
\tcblower
\textbf{Solution — model answer (about 175 words):}\par
\emph{Last year, cholera appeared in our town, and everybody was afraid. This dangerous disease can kill in a few days, can't it? We must act now to protect our school.}\par
\emph{Cholera is caused by dirty water and dirty food. People used to drink water from the river without boiling it, and they didn't use to wash their hands after using the toilet. As a result, the disease spread very quickly from family to family.}\par
\emph{However, prevention is simple and cheap. First, we must drink only clean, boiled or treated water. Moreover, hand washing with soap before meals is essential. The health worker who visited our school said that most cases of cholera could be avoided with these simple gestures, and she advised us to build clean latrines far from the water points.}\par
\emph{Therefore, dear schoolmates, let us keep our school grounds clean, use the dustbins and wash our hands every day. In conclusion, remember: a clean school is a healthy school!}
\end{exoresolu}

\begin{savaistu}[What really counts at the Probatoire technique]
At the Probatoire technique, the clarity of your \textbf{organisation} and the \textbf{grammatical accuracy} of your sentences count as much as the richness of your vocabulary. A composition of 150 to 200 words, clearly organised in three parts, with correct \emph{used to}, \emph{tag questions}, \emph{compound nouns} and \emph{reported speech}, will score higher than a long text full of beautiful but wrongly used words. Examiners read hundreds of papers: make yours easy to read and correct!
\end{savaistu}

\courssub{Integration activity: a School Health and Environment Day}

Now let us combine all your skills — speaking, listening, reading and writing — in one group project. This is the \cle{integration activity} of the lesson.

\begin{experience}[Group project: organising a School Health and Environment Day]
\textbf{Material:} cardboard, markers, the texts studied in this lesson, a notebook for the draft.\par
\textbf{Protocol (work in groups of four or five):}
\begin{enumerate}
\item \textbf{Listening}: listen to the teacher's text about planting trees and protecting natural habitats; note five key words.
\item \textbf{Reading}: read a short text about preventing a communicable disease (malaria, cholera or typhoid); extract the causes and the prevention measures.
\item \textbf{Speaking}: in your group, exchange information about a tourist or natural site in your region and explain why it must be protected; use at least two tag questions.
\item \textbf{Writing}: produce together one sensitization poster and one short article (150–200 words) for the school magazine, following the five-step writing process.
\item \textbf{Presentation}: each group presents its poster to the class in two minutes; the class asks questions.
\end{enumerate}
\textbf{Observations and interpretation:} each group checks with the self-assessment grid below whether the four skills were correctly used. The class votes for the clearest and most convincing poster.
\end{experience}

\courssub{Evaluation and remediation}

Use this grid to assess your own composition (self-assessment) or your classmate's (peer review). Give 0, 1 or 2 points per criterion: 0 = absent or wrong, 1 = present but imperfect, 2 = fully achieved.

\begin{center}
\begin{tabularx}{\linewidth}{|l|X|c|c|}
\hline
\rowcolor{popblueL} \textbf{Criterion} & \textbf{What the examiner checks} & \textbf{Max} & \textbf{Score} \\
\hline
Content & The issue, causes, consequences and prevention measures are all treated & 4 & \\
\hline
Vocabulary & At least six words of the unit (health, environment) used correctly & 4 & \\
\hline
Grammar & \emph{Used to}, one tag question, compound nouns, one reported sentence — all correct & 6 & \\
\hline
Organisation & Clear introduction / body / conclusion, linking words, 150–200 words & 4 & \\
\hline
Presentation & Clean handwriting, correct spelling and punctuation & 2 & \\
\hline
\rowcolor{popgreenL} \textbf{Total} & & \textbf{20} & \\
\hline
\end{tabularx}
\end{center}

\begin{exercice}[Targeted remediation on weak grammar points]
Rewrite or correct the following sentences.
\begin{enumerate}
\item Correct: \emph{Before, people drinked water from the stream, but now they boil it.} (use \emph{used to})
\item Add the correct tag question: \emph{Mosquito nets protect us from malaria, \ldots?}
\item Write the plural of: \emph{health centre}, \emph{mosquito net}, \emph{water point}.
\item Report the sentence: The nurse said: ``Vaccination is free for all children.''
\item Report the command: The doctor said to the pupils: ``Wash your hands before meals!''
\end{enumerate}
\end{exercice}

\begin{corrige}[Exercice : remediation]
\begin{enumerate}
\item \emph{People \textbf{used to drink} water from the stream, but now they boil it.} (past habit $\rightarrow$ \emph{used to} + base verb)
\item \emph{Mosquito nets protect us from malaria, \textbf{don't they}?} (affirmative sentence $\rightarrow$ negative tag; plural subject $\rightarrow$ \emph{they})
\item \emph{health centre\textbf{s}}, \emph{mosquito net\textbf{s}}, \emph{water point\textbf{s}} (only the second noun takes the plural \emph{s}).
\item \emph{The nurse said (that) vaccination \textbf{was} free for all children.} (present simple $\rightarrow$ past simple in reported speech)
\item \emph{The doctor \textbf{told/advised the pupils to wash} their hands before meals.} (reported command: \emph{tell/advise + object + to}-infinitive)
\end{enumerate}
\end{corrige}

\begin{exercice}[Final writing task]
Write a letter to a friend (150–200 words) telling him or her about the Health and Environment Day organised in your school. Say what your school used to be like before, describe the activities of the day, quote something a health worker said, and end with advice and a tag question. Follow the five-step writing process and use the evaluation grid to check your work.
\end{exercice}

\begin{aretenir}[To remember]
A good health composition is \textbf{organised} (introduction, body, conclusion), \textbf{precise} (compound nouns, unit vocabulary), \textbf{grammatical} (\emph{used to}, tag questions, reported speech) and \textbf{convincing} (linking words, call to action). Write simply, write correctly, and always think of your reader. Prevention is better than cure — and a clear composition is better than a long one, isn't it?
\end{aretenir}

\newpage

\coursec{Activité d'intégration}

\courssub{Objectifs de l'activité}

By the end of this integration activity, you will be able to:
\begin{itemize}
\item use vocabulary related to \cle{tourist sites}, \cle{gardening}, \cle{natural habitats} and \cle{communicable diseases} (\emph{maladies transmissibles}) in an authentic communication task;
\item report what people said using \cle{indirect speech} (statements, questions and commands);
\item use \cle{used to / didn't use to / to be used to} to speak about past habits and present habits;
\item build and pluralise \cle{compound nouns} (\emph{noms composés}) correctly;
\item write \cle{tag questions} to create catchy slogans (\emph{slogans accrocheurs});
\item write a short newspaper article and design a sensitization poster (\emph{affiche de sensibilisation});
\item speak in front of the class and react to classmates' presentations.
\end{itemize}

\courssub{Situation}

The \textbf{Ecology Club} of your technical high school in Maroua is organising a \textbf{School Health and Environment Day} next month. The Principal has asked every class of Première to prepare a \textbf{sensitization campaign} (\emph{campagne de sensibilisation}) for the whole school. Your class will produce two things: a short \textbf{article for the school magazine} and a \textbf{poster} with slogans.

The club gives your group of four students the following documents:

\textbf{Document 1 -- Poster from the District Health Service (Maroua):} ``Cholera is a dangerous communicable disease. Wash your hands with clean water and soap. Drink only treated or boiled water. Cook food well.''

\textbf{Document 2 -- Photographs of Waza National Park} (Far North Region), with elephants, giraffes and lions, and a short note: ``Waza National Park covers about \SI{170000}{hectares}. It receives about \num{10000} visitors per year.''

\textbf{Document 3 -- Figures (statistiques):}
\begin{itemize}
\item Malaria (\emph{paludisme}) causes about \num{4000} deaths per year in Cameroon, mostly children under five.
\item About \num{70} \% of Cameroonians live in areas where malaria is common.
\item The school planted \num{120} trees last year; only \num{80} survived. This year the club wants to plant \num{200} trees and protect them all.
\end{itemize}

\textbf{Document 4 -- Recorded interviews.} During a preparatory visit, the club interviewed two people:

\begin{itemize}
\item \textbf{Nurse Amina} (health centre) said: ``Malaria is the most common disease in our area. Sleep under mosquito nets and keep your environment clean. Many years ago, people did not use mosquito nets, but now they are used to sleeping under them.''
\item \textbf{Mr Bello} (tourist guide at Waza) said: ``Tourists used to visit the park only in December, but now they come all year. Do not throw plastic bags in the park. Protect the animals' natural habitat.''
\end{itemize}

\textbf{Your mission:} with your group, prepare the campaign and present it to the class. The best campaign will be displayed on the school notice board (\emph{panneau d'affichage}) during the Health and Environment Day.

\courssub{Consignes}

Work in groups of four. Follow the tasks in order.

\begin{enumerate}
\item \textbf{Read and select.} Read the four documents carefully. Underline ten useful words or expressions (for example: \emph{mosquito net, hand washing, natural habitat, tourist guide}). Classify them into two columns: \emph{Health} and \emph{Environment}.

\item \textbf{Compound nouns hunt.} Find at least \textbf{five compound nouns} in the documents (or build your own from the lesson vocabulary, e.g. \emph{mosquito net, health centre, tourist guide, washing powder, school magazine}). Write their plurals correctly.

\item \textbf{Report the interviews.} Write a short article (about 120--150 words) for the school magazine. Report what Nurse Amina and Mr Bello said using \textbf{indirect speech}: at least two reported statements, one reported question and one reported command. Do not forget the tense changes (backshift).

\item \textbf{Past and present habits.} In your article, include at least \textbf{two structures with \emph{used to}} (one affirmative or negative, and one with \emph{to be used to + noun/-ing}) to compare the past and the present.

\item \textbf{Design the poster.} Create a poster with a title, one drawing or symbol, and \textbf{three slogans containing tag questions} (example: \emph{We want a clean school, don't we?}). Check that each tag is correct (auxiliary, subject pronoun, positive/negative).

\item \textbf{Present your campaign.} Each group presents its article and poster to the class in three minutes. Speak clearly and loudly. The other students listen, then react and ask at least one question per group (use tag questions if possible: \emph{You planted the trees behind the classrooms, didn't you?}).

\item \textbf{Evaluate and improve.} Fill in the evaluation grid for one other group. Correct your own work according to the remarks (remediation).
\end{enumerate}

\courssub{Ressources à mobiliser}

\begin{aretenir}[Grammar reminders]
\textbf{Used to:} \emph{subject + used to + base verb} (past habit, finished). Negative: \emph{didn't use to}. Habit now: \emph{to be used to + noun/-ing} (être habitué à).

\textbf{Tag questions:} positive sentence $\rightarrow$ negative tag; negative sentence $\rightarrow$ positive tag. Use the same auxiliary and a pronoun: \emph{You wash your hands, don't you?}

\textbf{Indirect speech:} present simple $\rightarrow$ past simple; present continuous $\rightarrow$ past continuous; will $\rightarrow$ would; imperative $\rightarrow$ \emph{told/asked + object + (not) to + verb}. Questions: \emph{asked + if/whether} or \emph{wh-word}, affirmative word order, no question mark.

\textbf{Compound nouns:} the plural usually goes on the main (last) word: \emph{mosquito nets, health centres, tourist guides}.
\end{aretenir}

\begin{center}
\begin{tabularx}{\linewidth}{|l|X|X|}\hline
\rowcolor{popblueL} \textbf{Direct speech} & \textbf{Indirect speech} & \textbf{Change} \\ \hline
``Malaria is common.'' & She said (that) malaria was common. & is $\rightarrow$ was \\ \hline
``Sleep under nets.'' & She told us to sleep under nets. & imperative $\rightarrow$ to + verb \\ \hline
``Do you wash your hands?'' & She asked if we washed our hands. & do $\rightarrow$ if + past \\ \hline
``Don't throw plastic bags.'' & He told us not to throw plastic bags. & don't $\rightarrow$ not to \\ \hline
\end{tabularx}
\end{center}

\courssub{Démarche et corrigé commenté}

\begin{exoresolu}[Model answer: the campaign of Group 2]
Énoncé : produire l'article du journal scolaire et l'affiche du groupe, en appliquant toutes les consignes.

\tcblower

\textbf{Step 1 -- Vocabulary selection (Task 1).}

Health: \emph{mosquito net, hand washing, clean water, communicable disease, health centre}. Environment: \emph{natural habitat, tree planting, national park, plastic bags, tourist guide}.

\textbf{Step 2 -- Compound nouns and plurals (Task 2).}

\emph{mosquito net} $\rightarrow$ \emph{mosquito nets}; \emph{health centre} $\rightarrow$ \emph{health centres}; \emph{tourist guide} $\rightarrow$ \emph{tourist guides}; \emph{washing powder} $\rightarrow$ \emph{washing powders}; \emph{school magazine} $\rightarrow$ \emph{school magazines}. The plural -s goes on the last word, which carries the main meaning.

\textbf{Step 3 -- Article with indirect speech and \emph{used to} (Tasks 3 and 4).}

\emph{Health and Environment Day is coming! Last week, the Ecology Club interviewed Nurse Amina and Mr Bello, a tourist guide at Waza National Park.}

\emph{Nurse Amina said that malaria was the most common disease in our area. She told us to sleep under mosquito nets and to keep our environment clean. She explained that many years ago people didn't use to sleep under mosquito nets, but that now they were used to sleeping under them.}

\emph{Mr Bello said that tourists used to visit Waza only in December, but that they now came all year. He asked us if we liked animals, and he told us not to throw plastic bags in the park. He added that we had to protect the animals' natural habitat.}

\emph{Remember: malaria kills about \num{4000} people every year in Cameroon. Wash your hands, sleep under a net, and plant a tree!}

Justifications: ``is'' $\rightarrow$ ``was'' (backshift); the imperative ``Sleep'' $\rightarrow$ ``told us to sleep''; the negative imperative ``Do not throw'' $\rightarrow$ ``told us not to throw''; the yes/no question ``Do you like\ldots'' $\rightarrow$ ``asked us if we liked'' (affirmative order, no question mark). Two \emph{used to} structures are included: ``didn't use to sleep'' and ``used to visit'', plus ``were used to sleeping''.

\textbf{Step 4 -- Poster slogans with tag questions (Task 5).}

\begin{itemize}
\item We want a clean school, \textbf{don't we}?
\item Nobody likes cholera, \textbf{do they}?
\item Let's plant more trees, \textbf{shall we}?
\end{itemize}

Check: ``we want'' is positive $\rightarrow$ negative tag ``don't we''; ``nobody likes'' is negative in meaning $\rightarrow$ positive tag ``do they''; after ``Let's'', the tag is always ``shall we''.

\textbf{Step 5 -- Oral presentation (Task 6).} Each member of the group speaks: one presents the article, one reads the reported speech, one presents the poster, one answers questions.

\textbf{Step 6 -- Evaluation grid (Task 7).}

\begin{center}
\begin{tabularx}{\linewidth}{|l|X|c|}\hline
\rowcolor{popgreenL} \textbf{Criterion} & \textbf{What is expected} & \textbf{Mark /5} \\ \hline
Content & All documents used, figures correct & \\ \hline
Grammar & Indirect speech, used to, tag questions correct & \\ \hline
Vocabulary & At least 5 compound nouns, rich lexis & \\ \hline
Presentation & Clear voice, poster neat and attractive & \\ \hline
\end{tabularx}
\end{center}
\end{exoresolu}

\courssub{Bilan de l'activité}

This activity shows that the grammar and vocabulary of the lesson are not isolated rules: they are \textbf{tools for real communication}. Reporting interviews with \cle{indirect speech} is what journalists do every day; \cle{tag questions} make slogans lively and invite the audience to agree; \cle{used to} helps us compare past and present habits, for example in health education; \cle{compound nouns} make our English precise and natural. By writing an article, designing a poster and speaking in public, you have practised the four skills -- reading, writing, listening and speaking -- in one meaningful project, exactly like a real sensitization campaign organised by health services and ecology clubs in Cameroon.

\newpage

\coursec{Méthodes et exercices résolus}

\courssub{Méthodes et exercices résolus}

\begin{methode}[Parler de sites touristiques avec \emph{used to}]
\textbf{Objectif :} Décrire des habitudes passées ou des états anciens liés aux visites touristiques.
\medskip
\textbf{Étapes :}
\begin{enumerate}
    \item \textbf{Identifier le contexte :} S'agit-il d'une habitude passée (action répétée) ou d'un état passé ? $\rightarrow$ \emph{used to + BV}.
    \item \textbf{Former l'affirmation :} Sujet + \emph{used to} + base verbale. \textit{Ex: We \textbf{used to visit} the botanical garden every year.}
    \item \textbf{Former la négation :} Sujet + \emph{didn't use to} + base verbale (pas de \emph{d} à \emph{use}). \textit{Ex: They \textbf{didn't use to allow} cameras inside.}
    \item \textbf{Former la question :} \emph{Did} + sujet + \emph{use to} + base verbale ? \textit{Ex: \textbf{Did you use to go} there by bus?}
    \item \textbf{Ne pas confondre avec :} \emph{be used to + V-ing / nom} (avoir l'habitude de, présent). \textit{Ex: I \textbf{am used to walking} long distances.}
\end{enumerate}
\textbf{Vocabulaire clé :} \cle{tourist site} (site touristique), \cle{guided tour} (visite guidée), \cle{entrance fee} (droit d'entrée), \cle{souvenir} (souvenir), \cle{landscape} (paysage), \cle{heritage} (patrimoine).
\end{methode}

\begin{exoresolu}[Expression orale : souvenirs de visites]
\textbf{Énoncé :} Vous guidez un groupe de touristes au Parc National de Waza. En utilisant \emph{used to / didn't use to}, complétez le dialogue ci-dessous. Puis, à l'oral, posez 3 questions à un camarade sur ses anciennes visites scolaires.

\begin{center}
\begin{tabularx}{\linewidth}{|X|X|}\hline
\textbf{Guide :} Welcome to Waza Park! This area \_\_\_\_\_\_\_\_ (be) a hunting reserve. & \textbf{Tourist :} Really? \_\_\_\_\_\_\_\_ people \_\_\_\_\_\_\_\_ (hunt) lions here? \\ \hline
\textbf{Guide :} Yes, they \_\_\_\_\_\_\_\_ (hunt) them legally. But they \_\_\_\_\_\_\_\_ (protect) them. & \textbf{Tourist :} When \_\_\_\_\_\_\_\_ it \_\_\_\_\_\_\_\_ (become) a national park? \\ \hline
\textbf{Guide :} In 1968. Since then, rangers \_\_\_\_\_\_\_\_ (patrol) daily. & \textbf{Tourist :} \_\_\_\_\_\_\_\_ you \_\_\_\_\_\_\_\_ (work) here before? \\ \hline
\end{tabularx}
\end{center}

\tcblower
\textbf{Solution détaillée :}
\begin{enumerate}
    \item \textbf{used to be} — État passé (était une réserve de chasse). \textit{Français : « Cette zone était une réserve de chasse. »}
    \item \textbf{Did ... use to hunt} — Question sur habitude passée. \textit{« Est-ce que les gens chassaient des lions ici ? »}
    \item \textbf{used to hunt} — Habitude passée affirmative. \textit{« Oui, ils les chassaient légalement. »}
    \item \textbf{didn't use to protect} — Négation d'habitude passée (attention : pas de \emph{d} à \emph{use}). \textit{« Mais ils ne les protégeaient pas. »}
    \item \textbf{did ... become} — \textbf{Piège !} \emph{Become} est ici un événement ponctuel daté, pas une habitude. On utilise le \emph{preterite} : \textbf{did ... become}. \textit{« Quand est-elle devenue un parc national ? »}
    \item \textbf{have patrolled} ou \emph{have been patrolling} — Depuis 1968 (marqueur \emph{since}) $\rightarrow$ \emph{present perfect}. \textit{« Depuis, les gardes patrouillent chaque jour. »}
    \item \textbf{Did ... use to work} — Question sur habitude passée. \textit{« Est-ce que vous travailliez ici avant ? »}
\end{enumerate}
\textbf{Dialogue corrigé :}
\begin{quote}
\textbf{Guide :} Welcome to Waza Park! This area \textbf{used to be} a hunting reserve. \\
\textbf{Tourist :} Really? \textbf{Did} people \textbf{use to hunt} lions here? \\
\textbf{Guide :} Yes, they \textbf{used to hunt} them legally. But they \textbf{didn't use to protect} them. \\
\textbf{Tourist :} When \textbf{did} it \textbf{become} a national park? \\
\textbf{Guide :} In 1968. Since then, rangers \textbf{have patrolled} daily. \\
\textbf{Tourist :} \textbf{Did} you \textbf{use to work} here before?
\end{quote}
\textbf{Conclusion :} \emph{Used to} ne s'emploie que pour des habitudes/états \textbf{révolus}. Pour un événement daté unique, on utilise le \emph{preterite}. Pour une action commencée dans le passé et qui continue, on utilise le \emph{present perfect}.
\end{exoresolu}

\begin{methode}[Former et utiliser les \emph{tag questions}]
\textbf{Objectif :} Vérifier une information ou demander confirmation en fin de phrase.
\medskip
\textbf{Règle de base :} \emph{Affirmative statement} + \textbf{negative tag} \quad|\quad \emph{Negative statement} + \textbf{affirmative tag}
\medskip
\textbf{Étapes :}
\begin{enumerate}
    \item \textbf{Repérer l'auxiliaire/modal de la phrase principale} (be, have, do, can, will, must, etc.).
    \item \textbf{Inverser la polarité :} Positif $\rightarrow$ tag négatif (\emph{n't}) ; Négatif $\rightarrow$ tag positif.
    \item \textbf{Reprendre le sujet sous forme de pronom personnel} (I, you, he, she, it, we, they).
    \item \textbf{Accorder le temps :} Le tag est au même temps que la phrase principale.
    \item \textbf{Cas particuliers :}
    \begin{itemize}
        \item \emph{I am} $\rightarrow$ \textbf{aren't I?} (pas \emph{amn't I})
        \item \emph{Let's} $\rightarrow$ \textbf{shall we?}
        \item \emph{Imperatif} $\rightarrow$ \textbf{will you?} / \textbf{would you?}
        \item \emph{There is/are} $\rightarrow$ \textbf{isn't/aren't there?}
        \item \emph{Everybody/Somebody/No one} $\rightarrow$ \textbf{they?}
    \end{itemize}
    \item \textbf{Intonation :} Descendante = confirmation attendue ; Montante = vraie question.
\end{enumerate}
\textbf{Vocabulaire environnement :} \cle{deforestation} (déforestation), \cle{seedling} (jeune plant), \cle{ecosystem} (écosystème), \cle{endangered species} (espèce menacée), \cle{reforestation} (reboisement), \cle{habitat loss} (perte d'habitat).
\end{methode}

\begin{exoresolu}[Tag questions sur la protection de l'environnement]
\textbf{Énoncé :} Complétez chaque phrase par la \emph{tag question} correcte. Justifiez chaque choix en une phrase.

\begin{enumerate}
    \item Planting trees reduces soil erosion, \_\_\_\_\_\_\_\_ ?
    \item The villagers didn't cut down the old baobab, \_\_\_\_\_\_\_\_ ?
    \item There are many endangered species in this forest, \_\_\_\_\_\_\_\_ ?
    \item Let's organize a cleaning campaign next Saturday, \_\_\_\_\_\_\_\_ ?
    \item Nobody knows the exact number of gorillas left, \_\_\_\_\_\_\_\_ ?
    \item The river has been polluted by chemical waste, \_\_\_\_\_\_\_\_ ?
    \item You'll help us water the seedlings tomorrow, \_\_\_\_\_\_\_\_ ?
\end{enumerate}

\tcblower
\textbf{Solution détaillée :}
\begin{enumerate}
    \item \textbf{doesn't it?} — Phrase affirmative (\emph{reduces} = présent simple, 3\textsuperscript{e} pers. sg.) $\rightarrow$ tag négatif. Auxiliaire \emph{does} + \emph{n't} + \emph{it} (sujet \emph{Planting trees} = \emph{it}).
    \item \textbf{did they?} — Phrase négative (\emph{didn't cut}) $\rightarrow$ tag positif. Auxiliaire \emph{did} + \emph{they} (sujet \emph{The villagers} = \emph{they}).
    \item \textbf{aren't there?} — Structure \emph{There are} (affirmatif) $\rightarrow$ tag négatif \emph{aren't there?}.
    \item \textbf{shall we?} — Cas particulier \emph{Let's} (proposition) $\rightarrow$ toujours \emph{shall we?}.
    \item \textbf{do they?} — \emph{Nobody} (négatif de sens) $\rightarrow$ tag positif. Sujet indéfini $\rightarrow$ \emph{they}. Verbe \emph{knows} (présent simple) $\rightarrow$ auxiliaire \emph{do}.
    \item \textbf{hasn't it?} — \emph{Present perfect passive} (\emph{has been polluted}) affirmatif $\rightarrow$ tag négatif. Auxiliaire \emph{has} + \emph{n't} + \emph{it} (sujet \emph{The river} = \emph{it}).
    \item \textbf{won't you?} — \emph{Future simple} (\emph{will help}) affirmatif $\rightarrow$ tag négatif. \emph{Will} + \emph{n't} = \emph{won't} + \emph{you}.
\end{enumerate}
\textbf{Conclusion :} Identifier l'auxiliaire et la polarité de la principale est la clé. Attention aux sujets \emph{there}, \emph{nobody/everybody}, \emph{let's} et à l'auxiliaire \emph{do} au présent/prétérit simple.
\end{exoresolu}

\begin{methode}[Former les noms composés et leurs pluriels]
\textbf{Objectif :} Construire et accorder correctement les \emph{compound nouns} (noms composés) du vocabulaire santé/environnement.
\medskip
\textbf{Types principaux et règles de pluriel :}
\begin{enumerate}
    \item \textbf{Nom + Nom} : Le pluriel se marque sur le \textbf{deuxième} nom. \textit{Ex: \textbf{health centre} $\rightarrow$ \textbf{health centres} ; \textbf{blood test} $\rightarrow$ \textbf{blood tests}.}
    \item \textbf{Verbe + Particule (phrasal noun)} : Pluriel sur le \textbf{nom/verbe}. \textit{Ex: \textbf{check-up} $\rightarrow$ \textbf{check-ups} ; \textbf{breakdown} $\rightarrow$ \textbf{breakdowns}.}
    \item \textbf{Adjectif + Nom} : Pluriel sur le \textbf{nom}. \textit{Ex: \textbf{communicable disease} $\rightarrow$ \textbf{communicable diseases}.}
    \item \textbf{Nom + Préposition + Nom} : Pluriel sur le \textbf{premier} nom. \textit{Ex: \textbf{mother-in-law} $\rightarrow$ \textbf{mothers-in-law} ; \textbf{outbreak of cholera} $\rightarrow$ \textbf{outbreaks of cholera}.}
\end{enumerate}
\textbf{Hyphenation :} Vérifier au dictionnaire (ex: \emph{health-care} ou \emph{healthcare}, \emph{well-being} ou \emph{wellbeing}). Au Bac, les deux formes acceptées si cohérentes.
\textbf{Vocabulaire clé :} \cle{waterborne disease} (maladie hydrique), \cle{vector-borne disease} (maladie vectorielle), \cle{handwashing station} (station de lavage des mains), \cle{mosquito net} (moustiquaire), \cle{vaccination campaign} (campagne de vaccination), \cle{health worker} (agent de santé).
\end{methode}

\begin{exoresolu}[Noms composés : prévention des maladies]
\textbf{Énoncé :} Mettez les noms composés entre parenthèses au pluriel si nécessaire. Attention : certains restent au singulier (non-dénombrables ou concepts).

\begin{enumerate}
    \item The Ministry of Health launched three new \_\_\_\_\_\_\_\_ (vaccination campaign) last year.
    \item \_\_\_\_\_\_\_\_ (Malaria) is a dangerous \_\_\_\_\_\_\_\_ (vector-borne disease) transmitted by mosquitoes.
    \item Every \_\_\_\_\_\_\_\_ (health centre) in the district received five \_\_\_\_\_\_\_\_ (handwashing station).
    \item We need more \_\_\_\_\_\_\_\_ (community health worker) in rural areas.
    \item The \_\_\_\_\_\_\_\_ (outbreak of cholera) in the North region was controlled quickly.
    \item Sleeping under \_\_\_\_\_\_\_\_ (mosquito net) prevents malaria.
    \item Regular \_\_\_\_\_\_\_\_ (check-up) can detect diseases early.
\end{enumerate}

\tcblower
\textbf{Solution détaillée :}
\begin{enumerate}
    \item \textbf{vaccination campaigns} — \emph{Nom + Nom} : pluriel sur le 2\textsuperscript{e} nom (\emph{campaigns}). \emph{Three} impose le pluriel.
    \item \textbf{Malaria} (invariable, maladie) — \textbf{vector-borne disease} — \emph{Adjectif composé + Nom} : singulier après \emph{a} (une maladie). Pas de \emph{s} à \emph{disease}.
    \item \textbf{health centre} (singulier après \emph{every}) — \textbf{handwashing stations} — \emph{Nom (gérondif) + Nom} : pluriel sur \emph{stations} (\emph{five} impose le pluriel).
    \item \textbf{community health workers} — \emph{Nom + Nom + Nom} (trois éléments) : pluriel sur le dernier (\emph{workers}).
    \item \textbf{outbreak of cholera} — \emph{Nom + Préposition + Nom} : ici singulier (\emph{The ... was}). Au pluriel : \emph{outbreaks of cholera} (pluriel sur le 1\textsuperscript{er} nom). \emph{Cholera} invariable.
    \item \textbf{mosquito nets} — \emph{Nom + Nom} : pluriel sur \emph{nets}. (Généralité $\rightarrow$ pluriel).
    \item \textbf{check-ups} — \emph{Verbe + Particule} (phrasal noun) : pluriel sur le nom verbal (\emph{check-ups}). Trait d'union conservé.
\end{enumerate}
\textbf{Conclusion :} La règle d'or : identifier le \textbf{nom principal} (head noun) dans le composé. C'est lui qui prend la marque du pluriel. Pour les composés à trois mots ou avec préposition, c'est le \textbf{premier} nom.
\end{exoresolu}

\begin{methode}[Discours direct $\rightarrow$ indirect (Déclarations) — Partie 1]
\textbf{Objectif :} Reporter des affirmations en changeant les personnes, le temps, les marqueurs spatio-temporels.
\medskip
\textbf{Étapes (méthode en 5 temps) :}
\begin{enumerate}
    \item \textbf{Verbe introducteur au passé} (\emph{said, told, explained, added}) $\rightarrow$ \textbf{reculer d'un cran les temps} (backshifting).
    \item \textbf{Tableau de correspondance des temps (Backshifting) :}
    \begin{center}
    \begin{tabular}{|l|l|}\hline
    \textbf{Direct} & \textbf{Indirect} \\ \hline
    Present Simple & Past Simple \\
    Present Continuous & Past Continuous \\
    Present Perfect & Past Perfect \\
    Past Simple & Past Perfect \\
    Will & Would \\
    Can & Could \\
    Must & Had to \\ \hline
    \end{tabular}
    \end{center}
    \item \textbf{Changer les pronoms et possessifs} selon le nouveau locuteur (\emph{I $\rightarrow$ he/she, my $\rightarrow$ his/her, we $\rightarrow$ they, our $\rightarrow$ their}).
    \item \textbf{Adapter les marqueurs de temps/lieu :}
    \begin{itemize}
        \item \emph{now $\rightarrow$ then / at that moment}
        \item \emph{today $\rightarrow$ that day}
        \item \emph{yesterday $\rightarrow$ the day before / the previous day}
        \item \emph{tomorrow $\rightarrow$ the next day / the following day}
        \item \emph{here $\rightarrow$ there}
        \item \emph{this/these $\rightarrow$ that/those}
    \end{itemize}
    \item \textbf{Structure :} Sujet + verbe introducteur + \emph{(that)} + nouvelle phrase (ordre affirmatif, pas d'inversion, pas de point d'interrogation).
\end{enumerate}
\textbf{Exception :} Si le verbe introducteur est au \textbf{présent} (\emph{says, tells}) $\rightarrow$ \textbf{pas de backshifting}. Vérités générales $\rightarrow$ temps inchangé.
\end{methode}

\begin{exoresolu}[Discours indirect : déclarations sur la santé]
\textbf{Énoncé :} Transformez les phrases suivantes au discours indirect. Le verbe introducteur est au passé.

\begin{enumerate}
    \item The nurse said: "I \textbf{am checking} the patients' temperatures \textbf{now}."
    \item Dr. Atanga told us: "The vaccination campaign \textbf{started} \textbf{yesterday}."
    \item "We \textbf{have distributed} 500 mosquito nets \textbf{this month}," the coordinator explained.
    \item The minister announced: "We \textbf{will build} a new hospital \textbf{here} \textbf{next year}."
    \item "You \textbf{must boil} drinking water," the health worker warned the villagers.
\end{enumerate}

\tcblower
\textbf{Solution détaillée :}
\begin{enumerate}
    \item \textbf{The nurse said (that) she was checking the patients' temperatures then.}
    \begin{itemize}
        \item \emph{am checking} (present continuous) $\rightarrow$ \emph{was checking} (past continuous).
        \item \emph{I} $\rightarrow$ \emph{she} (infirmière = 3\textsuperscript{e} pers. fém.).
        \item \emph{now} $\rightarrow$ \emph{then}.
    \end{itemize}
    \item \textbf{Dr. Atanga told us (that) the vaccination campaign had started the day before.}
    \begin{itemize}
        \item \emph{started} (past simple) $\rightarrow$ \emph{had started} (past perfect).
        \item \emph{yesterday} $\rightarrow$ \emph{the day before}.
        \item \emph{told us} conserve \emph{us} (complément d'objet indirect).
    \end{itemize}
    \item \textbf{The coordinator explained (that) they had distributed 500 mosquito nets that month.}
    \begin{itemize}
        \item \emph{have distributed} (present perfect) $\rightarrow$ \emph{had distributed} (past perfect).
        \item \emph{We} $\rightarrow$ \emph{they} (l'équipe).
        \item \emph{this month} $\rightarrow$ \emph{that month}.
    \end{itemize}
    \item \textbf{The minister announced (that) they would build a new hospital there the following year.}
    \begin{itemize}
        \item \emph{will build} $\rightarrow$ \emph{would build}.
        \item \emph{We} $\rightarrow$ \emph{they} (le gouvernement).
        \item \emph{here} $\rightarrow$ \emph{there}.
        \item \emph{next year} $\rightarrow$ \emph{the following year} (ou \emph{the next year}).
    \end{itemize}
    \item \textbf{The health worker warned the villagers (that) they had to boil drinking water.}
    \begin{itemize}
        \item \emph{must} (obligation) $\rightarrow$ \emph{had to} (pas \emph{musted}).
        \item \emph{You} $\rightarrow$ \emph{they} (les villageois).
        \item \emph{warned} + objet + \emph{that} (verbe de communication forte).
    \end{itemize}
\end{enumerate}
\textbf{Conclusion :} Appliquer systématiquement les 4 changements : \textbf{temps, pronoms, marqueurs, structure}. Vérifier que la phrase rapportée est à la forme affirmative (sujet + verbe).
\end{exoresolu}

\begin{methode}[Discours indirect : Questions et Ordres — Partie 2]
\textbf{Objectif :} Reporter des questions (fermées/ouvertes) et des impératifs (ordres/conseils).
\medskip
\textbf{A. Questions (Interrogatives) :}
\begin{enumerate}
    \item \textbf{Question fermée (Yes/No)} $\rightarrow$ \emph{asked if / whether} + sujet + verbe (ordre affirmatif, \textbf{pas d'auxiliaire do/does/did}). Temps reculé.
    \item \textbf{Question ouverte (Wh-)} $\rightarrow$ \emph{asked} + \textbf{mot interrogatif} (who, what, where, when, why, how) + sujet + verbe (ordre affirmatif). Temps reculé.
    \item \textbf{Jamais} \emph{do/does/did} dans la phrase indirecte. Pas de point d'interrogation final.
\end{enumerate}
\textbf{B. Ordres / Conseils (Impératifs) :}
\begin{enumerate}
    \item \textbf{Affirmatif :} \emph{told / ordered / advised / warned} + objet + \textbf{to} + base verbale.
    \item \textbf{Négatif :} \emph{told / ordered / advised / warned} + objet + \textbf{not to} + base verbale.
    \item \textbf{Let's} $\rightarrow$ \emph{suggested} + \emph{V-ing} / \emph{that} + \emph{should} + BV.
\end{enumerate}
\textbf{Verbes introducteurs courants :} \emph{ask, wonder, want to know} (questions) ; \emph{tell, order, ask, advise, warn, beg, forbid} (ordres).
\end{methode}

\begin{exoresolu}[Questions et ordres rapportés : contexte sanitaire]
\textbf{Énoncé :} Mettez au discours indirect.

\begin{enumerate}
    \item The journalist asked: "How many cases of cholera \textbf{have been reported} this week?"
    \item " \textbf{Do you wash} your hands before eating?" the teacher asked the pupils.
    \item The doctor told the patient: " \textbf{Take} this medicine three times a day."
    \item " \textbf{Don't drink} water from the river," the chief warned the children.
    \item " \textbf{Let's clean} the school compound together," the prefect proposed.
\end{enumerate}

\tcblower
\textbf{Solution détaillée :}
\begin{enumerate}
    \item \textbf{The journalist asked how many cases of cholera had been reported that week.}
    \begin{itemize}
        \item Question ouverte (\emph{How many}) $\rightarrow$ conserve \emph{how many}.
        \item \emph{have been reported} (present perfect passive) $\rightarrow$ \emph{had been reported} (past perfect passive).
        \item \emph{this week} $\rightarrow$ \emph{that week}.
        \item Ordre affirmatif : \emph{cases ... had been reported} (pas \emph{had cases been reported}).
    \end{itemize}
    \item \textbf{The teacher asked the pupils if / whether they washed their hands before eating.}
    \begin{itemize}
        \item Question fermée $\rightarrow$ \emph{if / whether}.
        \item \emph{Do you wash} (present simple) $\rightarrow$ \emph{washed} (past simple). \textbf{Pas de \emph{did}}.
        \item \emph{you} $\rightarrow$ \emph{they} (les élèves). \emph{your} $\rightarrow$ \emph{their}.
    \end{itemize}
    \item \textbf{The doctor told the patient to take that medicine three times a day.}
    \begin{itemize}
        \item Impératif affirmatif \emph{Take} $\rightarrow$ \emph{to take}.
        \item \emph{told} + objet (\emph{the patient}) + \emph{to} + BV.
        \item \emph{this} $\rightarrow$ \emph{that}.
    \end{itemize}
    \item \textbf{The chief warned the children not to drink water from the river.}
    \begin{itemize}
        \item Impératif négatif \emph{Don't drink} $\rightarrow$ \emph{not to drink}.
        \item \emph{warned} + objet + \emph{not to} + BV.
    \end{itemize}
    \item \textbf{The prefect suggested cleaning the school compound together.} \quad(\textbf{ou} \textbf{suggested that they should clean}...)
    \begin{itemize}
        \item \emph{Let's} $\rightarrow$ \emph{suggested} + \emph{V-ing} (forme la plus courante au Bac).
        \item \emph{Let's} implique \emph{we} ; avec \emph{V-ing}, le sujet n'est pas exprimé.
    \end{itemize}
\end{enumerate}
\textbf{Conclusion :} Questions $\rightarrow$ \emph{asked if/whether} ou \emph{asked [Wh-]} + ordre affirmatif. Ordres $\rightarrow$ \emph{told/advised/warned} + objet + \emph{(not) to} + BV. \emph{Let's} $\rightarrow$ \emph{suggested} + \emph{V-ing}.
\end{exoresolu}

\begin{methode}[Rédiger une composition sur la santé / prévention]
\textbf{Objectif :} Produire un texte cohérent (150--200 mots) informant, sensibilisant ou argumentant sur un problème de santé.
\medskip
\textbf{Structure type (4 paragraphes) :}
\begin{enumerate}
    \item \textbf{Introduction (1 phrase d'accroche + thèse) :} Présenter le problème, son importance, annoncer le plan. \textit{Ex: "Cholera outbreaks remain a major threat in our region. This essay explains the causes, consequences and prevention measures."}
    \item \textbf{Développement 1 -- Causes / Transmission :} \emph{Link words : Firstly, The main cause is, Due to, Because of.} Vocabulaire : \cle{contaminated water}, \cle{poor sanitation}, \cle{open defecation}, \cle{lack of hygiene}.
    \item \textbf{Développement 2 -- Solutions / Prévention :} \emph{Link words : Secondly, To prevent this, We must, It is essential to.} Structures : \emph{People should / must / have to + BV}, \emph{The government needs to + BV}, \emph{Communities ought to + BV}. Vocabulaire : \cle{oral rehydration}, \cle{vaccination}, \cle{handwashing}, \cle{latrine construction}.
    \item \textbf{Conclusion (Appel à l'action) :} \emph{In conclusion, To sum up, Let us all...} Résumer + exhorter. \textit{Ex: "Everyone has a role to play. Let's wash hands, use latrines and report symptoms early."}
\end{enumerate}
\textbf{Grammar check-list :} \emph{Used to} (habitudes passées), \emph{Tag questions} (rhétoriques), \emph{Compound nouns} (vocabulaire précis), \emph{Reported speech} (citer autorités/experts), \emph{Modals} (obligation/conseil).
\end{methode}

\begin{exoresolu}[Rédaction : "Preventing communicable diseases in our community"]
\textbf{Énoncé :} Sujet de type Bac : \emph{"Write a composition of 150--200 words on the following topic: 'How can we prevent communicable diseases in our community?' Use at least 5 compound nouns, 2 tag questions, 1 reported speech sentence and 1 sentence with \emph{used to}."}

\tcblower
\textbf{Composition modèle (environ 180 mots) :}

\textbf{Preventing Communicable Diseases in Our Community}

Communicable diseases such as cholera, typhoid and malaria still affect many families in our district. \textbf{Didn't we use to lose} many children to diarrhoea every rainy season? \textbf{Health workers} say that things can change if we act together.

\textbf{Firstly}, the main causes are \textbf{contaminated water}, \textbf{poor sanitation} and the lack of \textbf{handwashing stations} in public places. Many households \textbf{used to drink} water directly from streams without boiling it. \textbf{Secondly}, prevention is possible. The District Medical Officer \textbf{stated that} every family must use \textbf{mosquito nets} and build improved \textbf{pit latrines}. \textbf{Isn't it time} we all took responsibility?

\textbf{Moreover}, \textbf{vaccination campaigns} against cholera and typhoid should reach every village. \textbf{Community health workers} must teach \textbf{oral rehydration} techniques to mothers. Schools need \textbf{handwashing facilities} with soap.

\textbf{In conclusion}, preventing diseases is cheaper than curing them. Let us boil water, wash hands, use latrines and sleep under treated nets. A healthy community is a wealthy community.

\medskip
\textbf{Analyse des exigences respectées :}
\begin{itemize}
    \item \textbf{5+ noms composés :} \emph{health workers, contaminated water, poor sanitation, handwashing stations, mosquito nets, pit latrines, vaccination campaigns, community health workers, oral rehydration, handwashing facilities} (10).
    \item \textbf{2 tag questions :} \emph{Didn't we use to lose...? ; Isn't it time...?}
    \item \textbf{1 reported speech :} \emph{The District Medical Officer stated that every family must use...} (discours indirect ; \emph{must} conservé car obligation toujours valable).
    \item \textbf{1 used to :} \emph{Many households used to drink... ; Didn't we use to lose...}
    \item \textbf{Connecteurs :} Firstly, Secondly, Moreover, In conclusion.
    \item \textbf{Modaux :} must, should, can, need to.
\end{itemize}
\textbf{Conclusion :} Une composition réussie intègre \textbf{explicitement} les points de grammaire demandés dans un texte fluide et argumenté, pas sous forme d'exercices isolés.
\end{exoresolu}

\begin{methode}[Activité d'intégration : Réinvestir l'unité complète]
\textbf{Objectif :} Mobiliser l'ensemble des savoirs (speaking, grammar, vocabulary, listening, reading, writing) dans une tâche complexe simulée.
\medskip
\textbf{Scénario type :} \emph{"Vous êtes membre du Club Santé-Environnement de votre lycée. Préparez une campagne de sensibilisation pour la Journée Mondiale de la Santé (7 avril)."}
\medskip
\textbf{Étapes de la tâche intégrée :}
\begin{enumerate}
    \item \textbf{Écoute/Prise de notes :} Écouter un court reportage (audio/vidéo) sur une campagne de reboisement / vaccination. Noter : problème, actions, résultats, citations.
    \item \textbf{Oral (Interaction) :} En binôme, jouer un dialogue : un journaliste interviewe le président du club. Utiliser \emph{used to} (passé), \emph{tag questions} (vérification), vocabulaire spécifique.
    \item \textbf{Grammaire en contexte :} Transformer 3 citations du reportage en discours indirect (déclarations + questions + ordres).
    \item \textbf{Lecture/Analyse :} Lire un texte sur une maladie hydrique. Identifier : noms composés, marqueurs de cause/conséquence, modaux de recommandation.
    \item \textbf{Écriture finale :} Produire une \textbf{affiche} (poster) + un \textbf{communiqué de presse} (100 mots) + un \textbf{discours d'ouverture} (2 min) pour le jour J.
\end{enumerate}
\textbf{Critères d'évaluation (grille type Bac) :}
\begin{itemize}
    \item Contenu / Respect du sujet (4 pts)
    \item Cohérence / Organisation (3 pts)
    \item Richesse lexicale (vocabulaire unité) (3 pts)
    \item Correction grammaticale (points ciblés) (5 pts)
    \item Orthographe / Ponctuation / Lisibilité (2 pts)
    \item Originalité / Impact (3 pts) — \textbf{Total : 20 pts}
\end{itemize}
\end{methode}

\begin{exoresolu}[Simulation d'intégration : Journée Mondiale de la Santé]
\textbf{Énoncé :} À partir du scénario ci-dessus, effectuez les sous-tâches 3 et 5 (grammaire en contexte + écriture finale).

\textbf{Données du reportage (simulé) :}
\begin{quote}
\textbf{Reporter :} "The reforestation project started three years ago. \textbf{We have planted} 10,000 trees. \textbf{Did the villagers participate} actively?" \\
\textbf{Chief :} "Yes, they \textbf{did}. \textbf{We used to cut} trees for firewood. Now \textbf{we protect} the forest. \textbf{Please, continue} supporting us."
\end{quote}

\textbf{Tâche 3 :} Mettez les 4 interventions au discours indirect.
\textbf{Tâche 5 :} Rédigez le \textbf{communiqué de presse} (100 mots) annonçant l'événement du 7 avril.

\tcblower
\textbf{Solution détaillée :}

\textbf{Tâche 3 -- Discours indirect :}
\begin{enumerate}
    \item The reporter said (that) the reforestation project \textbf{had started} three years \textbf{before}.
    \item The reporter said (that) they \textbf{had planted} 10,000 trees.
    \item The reporter \textbf{asked if / whether the villagers had participated} actively. (Question fermée $\rightarrow$ \emph{if/whether} + past perfect).
    \item The chief \textbf{confirmed/replied} (that) they \textbf{had}. He \textbf{explained} (that) they \textbf{used to cut} trees for firewood. (Habitude passée $\rightarrow$ \emph{used to} inchangé car déjà passé). He \textbf{added} (that) \textbf{they protected} the forest \textbf{then}. (Present simple $\rightarrow$ past simple ; \emph{now} $\rightarrow$ \emph{then}). He \textbf{asked/urged} them \textbf{to continue} supporting \textbf{them}. (Impératif \emph{Please continue} $\rightarrow$ \emph{asked/urged} + \emph{to continue} ; \emph{us} $\rightarrow$ \emph{them}).
\end{enumerate}

\textbf{Tâche 5 -- Communiqué de presse (modèle, environ 100 mots) :}

\begin{quote}
\textbf{PRESS RELEASE -- World Health Day}

\textbf{Yaoundé, April 5} -- The Health-Environment Club of \textbf{Biyem-Assi Technical High School} will celebrate World Health Day on \textbf{Monday, April 7}, under the theme \emph{"Healthy Forests, Healthy Communities"}.

The event starts at \textbf{8:00 AM} with a \textbf{tree-planting ceremony} on the school campus, followed by a \textbf{health walk} to the district medical centre. \textbf{Dr. Ngo Bayock}, District Medical Officer, will deliver a talk on \emph{"Preventing Waterborne Diseases through Environmental Protection"}.

\textbf{Free screening} for malaria, diabetes and hypertension will be offered. \textbf{Handwashing stations} and \textbf{mosquito nets} will be distributed to vulnerable families.

The Club President, \textbf{Mbarga Jean}, declared: "\emph{We used to ignore the link between deforestation and disease. Now we act.}" All students, parents and partners are invited.

\textbf{Contact :} Club Coordinator, Tel: 6XX XXX XXX.
\end{quote}

\textbf{Vérification :} Vocabulaire unité (\emph{tree-planting, health walk, waterborne diseases, handwashing stations, mosquito nets, deforestation}), \emph{used to}, discours indirect (\emph{declared that...}), structure formelle, connecteurs. \textbf{Conclusion :} L'intégration réussit quand chaque item grammatical sert la communication réelle, pas l'exercice artificiel.
\end{exoresolu}

\begin{attention}[Les 5 erreurs qui coûtent des points au Bac]
\begin{enumerate}
    \item \textbf{Confondre \emph{used to} (habitude passée) et \emph{be used to} (habitude présente).} \\
    \emph{Faux :} \textit{I used to wake up early now.} \quad \emph{Juste :} \textit{I \textbf{am used to waking} up early now.} / \textit{I \textbf{used to wake} up early \textbf{when I was young}.}
    \item \textbf{Oublier le \emph{backshifting} au discours indirect (verbe introducteur au passé).} \\
    \emph{Faux :} \textit{He said he \textbf{will} come.} \quad \emph{Juste :} \textit{He said he \textbf{would} come.} (Sauf vérité générale : \textit{The teacher said the Earth \textbf{revolves} around the Sun.})
    \item \textbf{Mettre l'auxiliaire \emph{do/does/did} dans la question indirecte.} \\
    \emph{Faux :} \textit{She asked if I \textbf{did} like it.} \quad \emph{Juste :} \textit{She asked if I \textbf{liked} it.}
    \item \textbf{Pluriel mal placé dans les noms composés.} \\
    \emph{Faux :} \textit{healths centres, mother-in-laws, check-ups} (si trait d'union oublié : \textit{checkups} accepté mais \textit{check-ups} préféré). \quad \emph{Juste :} \textit{health centres, mothers-in-law, check-ups.}
    \item \textbf{Tag question : mauvaise polarité ou mauvais pronom.} \\
    \emph{Faux :} \textit{Let's go, \textbf{will we?} / Nobody came, \textbf{didn't they?} / I'm right, \textbf{amn't I?}} \\
    \emph{Juste :} \textit{Let's go, \textbf{shall we?} / Nobody came, \textbf{did they?} / I'm right, \textbf{aren't I?}}
\end{enumerate}
\textbf{Astuce de révision :} Pour chaque point, écrivez 3 phrases \emph{Faux/Juste} et expliquez la règle à un camarade. C'est la meilleure façon d'ancrer la correction.
\end{attention}

\newpage

\coursec{Exercices}

\coursec{Environment, Wellbeing and Health}

\courssub{Je vérifie mes connaissances}

\begin{exercice}[QCM : Used to / Be used to]
Choose the correct option to complete each sentence. (\emph{Choisis la bonne option.})
\begin{enumerate}
    \item When I was a child, my family \_\_\_\_\_\_ visit the Waza National Park every Christmas.
    \begin{itemize}
        \item[a)] used to \item[b)] are used to \item[c)] use to \item[d)] didn't used to
    \end{itemize}
    \item The tour guide \_\_\_\_\_\_ walking long distances in the Mandara Mountains; he does it every week.
    \begin{itemize}
        \item[a)] used to \item[b)] is used to \item[c)] uses to \item[d)] didn't use to
    \end{itemize}
    \item Tourists \_\_\_\_\_\_ pay an entrance fee to visit the Foumban Royal Palace, but now it is free on Sundays.
    \begin{itemize}
        \item[a)] are used to \item[b)] use to \item[c)] used to \item[d)] were used to
    \end{itemize}
    \item \_\_\_\_\_\_ you \_\_\_\_\_\_ eat \textit{ndolé} (\textit{plat traditionnel}) when you lived in Douala?
    \begin{itemize}
        \item[a)] Did / used to \item[b)] Are / used to \item[c)] Do / use to \item[d)] Were / used to
    \end{itemize}
    \item We \_\_\_\_\_\_ the hot climate in Maroua; we have lived there for ten years.
    \begin{itemize}
        \item[a)] used to \item[b)] didn't use to \item[c)] are used to \item[d)] use to
    \end{itemize}
\end{enumerate}
\end{exercice}

\begin{exercice}[Vrai ou Faux justifié : Tag Questions \& Environnement]
Read the statements. Write \textbf{True} or \textbf{False} and justify your answer with the correct tag question or vocabulary precision. (\emph{Écris Vrai ou Faux et justifie.})
\begin{enumerate}
    \item "Let's plant trees around the school compound, \textbf{shall we}?" \rightarrow The tag question is correct for a suggestion with "Let's".
    \item "Nobody threw waste in the river, \textbf{did they}?" \rightarrow The tag is correct because "Nobody" is a negative subject, requiring a positive tag.
    \item "A \textbf{seedling} (\textit{jeune plant}) is a fully grown tree ready for cutting."
    \item "Deforestation (\textit{déforestation}) helps protect natural habitats (\textit{habitats naturels}) for wildlife (\textit{faune sauvage})."
    \item "The gardener waters the flowers every morning, \textbf{doesn't he}?" \rightarrow The tag question follows the rule: positive statement, negative tag.
\end{enumerate}
\end{exercice}

\begin{exercice}[Définitions et Vocabulaire : Santé et Noms Composés]
Match each term (1--5) with its correct definition (a--e). Write the letter in the box. Then write the \textbf{plural form} of each compound noun. (\emph{Associe et donne le pluriel.})
\begin{center}
\begin{tabular}{|c|l|l|}
\hline
\textbf{Term} & \textbf{Definition} & \textbf{Plural} \\
\hline
1. Mosquito net (\textit{moustiquaire}) & a. A place where sick people are treated. & \_\_\_\_\_\_\_\_\_\_\_\_\_\_\_\_ \\
2. Health centre (\textit{centre de santé}) & b. A net used to prevent mosquito bites. & \_\_\_\_\_\_\_\_\_\_\_\_\_\_\_\_ \\
3. First-aid kit (\textit{trousse de secours}) & c. A disease spread from person to person. & \_\_\_\_\_\_\_\_\_\_\_\_\_\_\_\_ \\
4. Communicable disease (\textit{maladie transmissible}) & d. A box containing medical supplies for emergencies. & \_\_\_\_\_\_\_\_\_\_\_\_\_\_\_\_ \\
5. Passer-by (\textit{passant}) & e. A person walking past a place. & \_\_\_\_\_\_\_\_\_\_\_\_\_\_\_\_ \\
\hline
\end{tabular}
\end{center}
\end{exercice}

\begin{exercice}[Discours Direct / Indirect : Transformations de base]
Rewrite the following sentences in \textbf{Indirect (Reported) Speech}. The reporting verb is in the past (\textit{said / told / asked}). (\emph{Transforme en discours indirect.})
\begin{enumerate}
    \item The nurse said: "I \textbf{am vaccinating} the children today."
    \item The doctor told me: "You \textbf{must} drink clean water."
    \item A tourist asked: "\textbf{Where} is the nearest hospital?"
    \item The guide said: "The museum \textbf{opens} at 9 a.m."
    \item The teacher ordered: "\textbf{Don't touch} the medical waste."
\end{enumerate}
\end{exercice}

\courssub{Je m'entraîne}

\begin{exercice}[Dialogue complet : Visite touristique à Foumban]
Complete the dialogue between a Tourist (T) and a Guide (G) at the Foumban Royal Palace. Use the correct form: \textbf{used to / didn't use to / be used to (present or past)} + verb in brackets. (\emph{Complète le dialogue avec la forme correcte.})
\begin{enumerate}
    \item G: Welcome! Kings \_\_\_\_\_\_\_\_\_\_\_\_\_\_ (live) in this palace for centuries.
    \item T: Did they \_\_\_\_\_\_\_\_\_\_\_\_\_\_ (receive) foreign visitors often?
    \item G: Yes, they \_\_\_\_\_\_\_\_\_\_\_\_\_\_ (host) many diplomats. But they \_\_\_\_\_\_\_\_\_\_\_\_\_\_ (allow) photos inside the throne room (now it is allowed).
    \item T: I \_\_\_\_\_\_\_\_\_\_\_\_\_\_ (not / walk) on such beautiful traditional carpets.
    \item G: You will get used to it. The artisans \_\_\_\_\_\_\_\_\_\_\_\_\_\_ (work) here every day.
    \item T: \_\_\_\_\_\_\_\_\_\_\_\_\_\_ they \_\_\_\_\_\_\_\_\_\_\_\_\_\_ (make) bronze sculptures in the past?
    \item G: Absolutely. They \_\_\_\_\_\_\_\_\_\_\_\_\_\_ (use) the lost-wax technique.
    \item T: I \_\_\_\_\_\_\_\_\_\_\_\_\_\_ (see) many crafts, but these are unique.
    \item G: Tourists \_\_\_\_\_\_\_\_\_\_\_\_\_\_ (buy) many souvenirs here.
    \item T: I \_\_\_\_\_\_\_\_\_\_\_\_\_\_ (not / spend) much money usually, but I will today!
\end{enumerate}
\end{exercice}

\begin{exercice}[Tag Questions : Environnement et Santé]
Add the correct \textbf{Tag Question} to each sentence. (\emph{Ajoute la question tag correcte.})
\begin{enumerate}
    \item We should protect the Korup National Park, \_\_\_\_\_\_\_\_\_\_\_\_\_\_ ?
    \item Cutting down trees causes erosion, \_\_\_\_\_\_\_\_\_\_\_\_\_\_ ?
    \item Nobody wants to drink polluted water, \_\_\_\_\_\_\_\_\_\_\_\_\_\_ ?
    \item Let's organize a cleaning campaign this Saturday, \_\_\_\_\_\_\_\_\_\_\_\_\_\_ ?
    \item The nurses have vaccinated the villagers, \_\_\_\_\_\_\_\_\_\_\_\_\_\_ ?
    \item Malaria is transmitted by mosquitoes, \_\_\_\_\_\_\_\_\_\_\_\_\_\_ ?
    \item You rarely sleep without a mosquito net, \_\_\_\_\_\_\_\_\_\_\_\_\_\_ ?
    \item There isn't any medicine left in the kit, \_\_\_\_\_\_\_\_\_\_\_\_\_\_ ?
    \item Everyone must wash their hands before eating, \_\_\_\_\_\_\_\_\_\_\_\_\_\_ ?
    \item I am right about the prevention measures, \_\_\_\_\_\_\_\_\_\_\_\_\_\_ ?
\end{enumerate}
\end{exercice}

\begin{exercice}[Pluriels des noms composés : Spécialité Santé]
Write the \textbf{plural form} of the following compound nouns. Be careful with the rules (head noun plural, both parts plural, or fixed forms). (\emph{Écris le pluriel.})
\begin{enumerate}
    \item toothbrush (\textit{brosse à dents}) \hfill 6. mother-in-law (\textit{belle-mère})
    \item health centre (\textit{centre de santé}) \hfill 7. grown-up (\textit{adulte})
    \item passer-by (\textit{passant}) \hfill 8. blood test (\textit{analyse de sang})
    \item first-aid kit (\textit{trousse de secours}) \hfill 9. x-ray (\textit{radio / radiographie})
    \item mosquito net (\textit{moustiquaire}) \hfill 10. medical check-up (\textit{bilan médical})
\end{enumerate}
\end{exercice}

\begin{exercice}[Discours Indirect : Déclarations, Questions, Ordres (Médecin-chef)]
Report the following statements made by the \textbf{District Medical Officer} during a health campaign. Use reporting verbs: \textit{said, told, asked, ordered, advised, warned}. (\emph{Rapport le discours du médecin-chef.})
\begin{enumerate}
    \item "Cholera spreads through contaminated water."
    \item "Do not drink water from the river without boiling it."
    \item "How many cases of measles have you recorded this month?"
    \item "The vaccination campaign will start next Monday."
    \item "Where is the nearest health centre located?"
    \item "Parents must bring their children for routine immunization."
    \item "Did you distribute the mosquito nets to the pregnant women?"
    \item "Wash your hands with soap before every meal."
\end{enumerate}
\end{exercice}

\courssub{Je me prépare au Bac}

\begin{exercice}[Compréhension de lecture : La lutte contre le paludisme au Cameroun]
Read the text carefully and answer the questions that follow. (\emph{Lis le texte et réponds aux questions.})

\textbf{Malaria Prevention in Cameroon: A Collective Effort}

Malaria remains a major public health challenge in Cameroon. According to the Ministry of Public Health, the disease accounts for approximately 30\% of consultations in health facilities and is a leading cause of morbidity and mortality, especially among children under five and pregnant women. The primary vector is the female Anopheles mosquito, which breeds in stagnant water (\textit{eau stagnante}) found in puddles, blocked gutters (\textit{caniveaux bouchés}), and uncovered containers.

To combat this scourge (\textit{fléau}), the government, with support from partners like the Global Fund and WHO, has implemented several strategies. The mass distribution of Long-Lasting Insecticidal Nets (LLINs / \textit{Moustiquaires Imprégnées à Longue Durée d'Action - MILD}) is the cornerstone (\textit{pierre angulaire}) of prevention. Since 2011, nationwide campaigns have distributed millions of nets free of charge to households. "Sleeping under a treated net every night is the single most effective way to prevent malaria," stated the Minister of Health during the last launch ceremony in Yaoundé.

Indoor Residual Spraying (IRS / \textit{Pulvérisation Intradosicielle à Effet Rémanent - PIER}) is used in targeted high-transmission areas in the North and Far North regions. Furthermore, Intermittent Preventive Treatment in pregnancy (IPTp) with Sulphadoxine-Pyrimethamine is offered free during antenatal consultations (\textit{consultations prénatales}) to protect mothers and unborn babies.

Community engagement is vital. Community Health Workers (\textit{Agents de Santé Communautaires}) educate families on environmental hygiene: clearing bushes (\textit{débroussailler}), draining stagnant water, and covering water storage containers. Early diagnosis using Rapid Diagnostic Tests (RDTs / \textit{Tests de Diagnostic Rapide - TDR}) and prompt treatment with Artemisinin-based Combination Therapies (ACTs / \textit{Thérapies de Combinaison à base d'Artémisinine - TCA}) are available in all public health centres.

However, challenges persist. Misuse of nets (for fishing or gardening), resistance to insecticides, and self-medication with incomplete doses hinder progress. The National Malaria Control Programme (NMCP / \textit{Programme National de Lutte contre le Paludisme - PNLP}) aims for a "Malaria-Free Cameroon" by 2035, aligning with the WHO Global Technical Strategy.

\vspace{0.5em}
\textbf{Questions:}
\begin{enumerate}
    \item \textbf{True or False?} Justify by quoting the text. (\emph{Vrai ou Faux ? Justifie par une citation.})
    \begin{itemize}
        \item[a)] Malaria is the main reason for consultation in Cameroonian health facilities.
        \item[b)] LLINs are sold at a subsidized price to households.
        \item[c)] IRS is applied uniformly across all ten regions of Cameroon.
        \item[d)] Community Health Workers play a role in environmental hygiene education.
    \end{itemize}
    \item \textbf{Open questions:} (\emph{Questions ouvertes.})
    \begin{itemize}
        \item[e)] Name two specific groups most vulnerable to malaria according to the text.
        \item[f)] What does IPTp stand for and who benefits from it?
        \item[g)] Identify two challenges hindering malaria elimination mentioned in the last paragraph.
    \end{itemize}
    \item \textbf{Vocabulary in context:} (\emph{Vocabulaire en contexte.}) Find words/phrases in the text that mean:
    \begin{itemize}
        \item[h)] "A disease that causes death" (paragraph 1) $\rightarrow$ \_\_\_\_\_\_\_\_\_\_\_\_\_\_
        \item[i)] "The most important part of something" (paragraph 2) $\rightarrow$ \_\_\_\_\_\_\_\_\_\_\_\_\_\_
        \item[j)] "Giving medicine to oneself without a doctor's advice" (paragraph 5) $\rightarrow$ \_\_\_\_\_\_\_\_\_\_\_\_\_\_
    \end{itemize}
\end{enumerate}
\end{exercice}

\begin{exercice}[Rédaction type Probatoire : Article pour le magazine scolaire]
\textbf{Task:} Your school is organizing a campaign against communicable diseases. Write an article of \textbf{150--200 words} for the school magazine. (\emph{Ta école organise une campagne contre les maladies transmissibles. Écris un article de 150--200 mots pour le magazine scolaire.})

\textbf{Your article must include:}
\begin{itemize}
    \item A catchy title.
    \item Presentation of 2--3 common communicable diseases (causes, symptoms, transmission).
    \item Prevention measures (hygiene, vaccination, environment).
    \item At least \textbf{two direct quotations} from the school nurse (use Direct Speech) and \textbf{one indirect quotation} (Indirect Speech).
    \item A call to action for students.
\end{itemize}
\textbf{Useful vocabulary:} outbreak (\textit{épidémie}), contagious (\textit{contagieux}), immunization (\textit{vaccination}), sanitation (\textit{assainissement}), handwashing (\textit{lavage des mains}), "Prevention is better than cure".
\end{exercice}

\begin{exercice}[Compréhension orale simulée : Campagne de reboisement scolaire]
\textbf{Instructions:} Listen to the text read by your teacher (or read the transcript below if practicing alone) and complete the \textbf{Information Sheet}. (\emph{Écoute le texte et complète la fiche d'informations.})

\textbf{Transcript (Teacher's copy -- do not read to students initially):}
"Good morning everyone. Last Saturday, June 15th, our school, Government Technical High School Buea, launched its annual 'Green School' reforestation campaign. The event started at 8:00 a.m. sharp. Over 350 students from Form 1 to Upper Sixth participated, supervised by 15 teachers and the school administration. We received a donation of 2,000 seedlings from the Ministry of Forestry and Wildlife (MINFOF). These included 1,200 Prunus africana (\textit{Prunier d'Afrique}), 500 Mahogany (\textit{Acajou}), and 300 Eucalyptus trees. The students were divided into 20 groups. Each group planted 100 trees along the school boundaries and near the water catchment area (\textit{prise d'eau}) to protect the source. The Principal, Mr. Ngome, declared: 'We are not just planting trees; we are planting our future.' The campaign ended at 1:00 p.m. with a refreshment. Next year, we aim to plant 3,000 trees."

\vspace{0.5em}
\textbf{INFORMATION SHEET: SCHOOL REFORESTATION CAMPAIGN}
\begin{center}
\begin{tabular}{|l|l|}
\hline
\textbf{Information Required} & \textbf{Answers} \\
\hline
Name of School & \_\_\_\_\_\_\_\_\_\_\_\_\_\_\_\_\_\_\_\_\_\_\_\_\_\_\_\_\_\_\_\_\_\_ \\
Date of Campaign & \_\_\_\_\_\_\_\_\_\_\_\_\_\_\_\_\_\_\_\_\_\_\_\_\_\_\_\_\_\_\_\_\_\_ \\
Start Time / End Time & \_\_\_\_\_\_\_\_\_\_\_\_\_\_\_\_\_\_\_\_\_\_\_\_\_\_\_\_\_\_\_\_\_\_ \\
Number of Student Participants & \_\_\_\_\_\_\_\_\_\_\_\_\_\_\_\_\_\_\_\_\_\_\_\_\_\_\_\_\_\_\_\_\_\_ \\
Number of Teachers Involved & \_\_\_\_\_\_\_\_\_\_\_\_\_\_\_\_\_\_\_\_\_\_\_\_\_\_\_\_\_\_\_\_\_\_ \\
Donor of Seedlings & \_\_\_\_\_\_\_\_\_\_\_\_\_\_\_\_\_\_\_\_\_\_\_\_\_\_\_\_\_\_\_\_\_\_ \\
Total Number of Seedlings & \_\_\_\_\_\_\_\_\_\_\_\_\_\_\_\_\_\_\_\_\_\_\_\_\_\_\_\_\_\_\_\_\_\_ \\
Breakdown of Species & 1. \_\_\_\_\_\_\_\_\_\_\_\_ (\_\_\_\_\_\_) \quad 2. \_\_\_\_\_\_\_\_\_\_\_\_ (\_\_\_\_\_\_) \quad 3. \_\_\_\_\_\_\_\_\_\_\_\_ (\_\_\_\_\_\_) \\
Planting Locations (2) & 1. \_\_\_\_\_\_\_\_\_\_\_\_\_\_\_\_\_\_ \quad 2. \_\_\_\_\_\_\_\_\_\_\_\_\_\_\_\_\_\_ \\
Principal's Quote (Direct Speech) & \_\_\_\_\_\_\_\_\_\_\_\_\_\_\_\_\_\_\_\_\_\_\_\_\_\_\_\_\_\_\_\_\_\_ \\
Next Year's Target & \_\_\_\_\_\_\_\_\_\_\_\_\_\_\_\_\_\_\_\_\_\_\_\_\_\_\_\_\_\_\_\_\_\_ \\
\hline
\end{tabular}
\end{center}
\end{exercice}

\newpage

\coursec{Corrigés des exercices}

\begin{corrige}[Exercice 1 — QCM : Used to / Be used to]
\begin{enumerate}
    \item \textbf{a) used to} — Past habit that no longer exists (\emph{when I was a child}): \emph{used to + base verb}.
    \item \textbf{b) is used to} — Present routine/habit the person is accustomed to: \emph{be used to + -ing} (\emph{walking}).
    \item \textbf{c) used to} — A past situation that has changed (\emph{but now it is free}).
    \item \textbf{a) Did / use to} — Interrogative of \emph{used to} in the past: \emph{Did + subject + use to...?} (After auxiliary \emph{did}, the main verb takes the base form \emph{use to}.)
    \item \textbf{c) are used to} — Present state of being accustomed to something: \emph{be used to + noun} (\emph{the hot climate}).
\end{enumerate}
\textbf{Point de vigilance :} \emph{used to + infinitive} = habitude passée révolue ; \emph{be used to + -ing/noun} = être habitué à. Ne jamais écrire \emph{didn't used to} (écrire \emph{didn't use to}).
\end{corrige}

\begin{corrige}[Exercice 2 — Vrai ou Faux justifié : Tag Questions \& Environnement]
\begin{enumerate}
    \item \textbf{True.} With \emph{Let's...}, the tag is always \emph{shall we?}
    \item \textbf{True.} \emph{Nobody} is negative in meaning, so the tag is positive: \emph{did they?} (\emph{Nobody} takes \emph{they} in the tag.)
    \item \textbf{False.} A \emph{seedling} is a \textbf{young plant}, not a fully grown tree.
    \item \textbf{False.} Deforestation \textbf{destroys} natural habitats; \emph{reforestation} (planting trees) protects them.
    \item \textbf{True.} Positive statement (\emph{waters}, present simple, 3rd person) $\rightarrow$ negative tag with the auxiliary \emph{does}: \emph{doesn't he?}
\end{enumerate}
\end{corrige}

\begin{corrige}[Exercice 3 — Définitions et Vocabulaire : Santé et Noms Composés]
Matching and plurals:
\begin{center}
\begin{tabular}{|c|c|l|}
\hline
\textbf{Term} & \textbf{Definition} & \textbf{Plural} \\
\hline
1. Mosquito net & b & mosquito nets \\
2. Health centre & a & health centres \\
3. First-aid kit & d & first-aid kits \\
4. Communicable disease & c & communicable diseases \\
5. Passer-by & e & passers-by \\
\hline
\end{tabular}
\end{center}
\textbf{Rule reminder :} for hyphenated compounds with a head noun (\emph{passer-by}), the head noun takes the plural: \emph{passers-by}. For simple compounds, only the last word takes \emph{-s}: \emph{mosquito nets}.
\end{corrige}

\begin{corrige}[Exercice 4 — Discours Direct / Indirect : Transformations de base]
\begin{enumerate}
    \item The nurse said (that) she \textbf{was vaccinating} the children \textbf{that day}. (present continuous $\rightarrow$ past continuous ; today $\rightarrow$ that day)
    \item The doctor told me (that) I \textbf{had to} drink clean water. (\emph{must} $\rightarrow$ \emph{had to})
    \item A tourist asked \textbf{where the nearest hospital was}. (wh-question: statement word order, no auxiliary, no question mark ; is $\rightarrow$ was)
    \item The guide said (that) the museum \textbf{opened} at 9 a.m. (present simple $\rightarrow$ past simple)
    \item The teacher ordered (us/them) \textbf{not to touch} the medical waste. (negative imperative $\rightarrow$ \emph{told/ordered + object + not to + infinitive})
\end{enumerate}
\textbf{Point de vigilance :} in reported questions, keep the affirmative word order and remove the question mark.
\end{corrige}

\begin{corrige}[Exercice 5 — Dialogue complet : Visite touristique à Foumban]
\begin{enumerate}
    \item \textbf{used to live} (past situation over centuries)
    \item \textbf{use to receive} (after \emph{Did}, base form \emph{use to})
    \item \textbf{used to host} ; \textbf{didn't use to allow} (past habit that has changed — now allowed)
    \item \textbf{am not used to walking} (present: be used to + -ing)
    \item \textbf{are used to working} (present habit of the artisans)
    \item \textbf{Did / use to make} (question about the past)
    \item \textbf{used to use} (past technique)
    \item \textbf{am used to seeing} (present: accustomed to)
    \item \textbf{are used to buying} (present habit) — \emph{used to buy} also acceptable if the guide speaks of the past, but the present context favours \emph{are used to buying}.
    \item \textbf{am not used to spending} (present: be used to + -ing)
\end{enumerate}
\end{corrige}

\begin{corrige}[Exercice 6 — Tag Questions : Environnement et Santé]
\begin{enumerate}
    \item \textbf{shouldn't we?} (modal \emph{should} $\rightarrow$ negative tag)
    \item \textbf{doesn't it?} (present simple, subject = gerund phrase = \emph{it})
    \item \textbf{do they?} (\emph{Nobody} negative $\rightarrow$ positive tag, pronoun \emph{they})
    \item \textbf{shall we?} (fixed tag after \emph{Let's})
    \item \textbf{haven't they?} (present perfect, auxiliary \emph{have})
    \item \textbf{isn't it?} (verb \emph{be}, present simple)
    \item \textbf{do you?} (\emph{rarely} is negative in meaning $\rightarrow$ positive tag)
    \item \textbf{is there?} (negative statement with \emph{there isn't} $\rightarrow$ positive tag \emph{is there})
    \item \textbf{mustn't they?} (\emph{Everyone} $\rightarrow$ \emph{they} ; modal \emph{must} repeated in the tag)
    \item \textbf{aren't I?} (irregular tag for \emph{I am})
\end{enumerate}
\textbf{Point de vigilance :} mots à valeur négative (\emph{nobody, rarely, never, hardly}) imposent un tag positif ; \emph{I am} donne \emph{aren't I?}.
\end{corrige}

\begin{corrige}[Exercice 7 — Pluriels des noms composés : Spécialité Santé]
\begin{enumerate}
    \item toothbrush $\rightarrow$ \textbf{toothbrushes} (plural of the final noun ; -sh $\rightarrow$ -shes)
    \item health centre $\rightarrow$ \textbf{health centres}
    \item passer-by $\rightarrow$ \textbf{passers-by} (head noun \emph{passer} takes the plural)
    \item first-aid kit $\rightarrow$ \textbf{first-aid kits}
    \item mosquito net $\rightarrow$ \textbf{mosquito nets}
    \item mother-in-law $\rightarrow$ \textbf{mothers-in-law} (head noun \emph{mother})
    \item grown-up $\rightarrow$ \textbf{grown-ups} (no noun inside $\rightarrow$ -s at the end)
    \item blood test $\rightarrow$ \textbf{blood tests}
    \item x-ray $\rightarrow$ \textbf{x-rays}
    \item medical check-up $\rightarrow$ \textbf{medical check-ups}
\end{enumerate}
\textbf{Point de vigilance :} \emph{mothers-in-law}, \emph{passers-by} : le nom principal porte la marque du pluriel ; \emph{grown-ups}, \emph{check-ups} : pas de nom principal, donc -s final.
\end{corrige}

\begin{corrige}[Exercice 8 — Discours Indirect : Déclarations, Questions, Ordres (Médecin-chef)]
\begin{enumerate}
    \item The District Medical Officer said (that) cholera \textbf{spreads / spread} through contaminated water. (general truth: the present \emph{spreads} may be kept ; backshift \emph{spread} also accepted)
    \item He \textbf{warned} us \textbf{not to drink} water from the river without boiling it.
    \item He \textbf{asked} how many cases of measles we \textbf{had recorded} that month. (present perfect $\rightarrow$ past perfect ; this $\rightarrow$ that)
    \item He said (that) the vaccination campaign \textbf{would start} the following Monday. (will $\rightarrow$ would ; next Monday $\rightarrow$ the following Monday)
    \item He \textbf{asked} where the nearest health centre \textbf{was located}. (affirmative word order)
    \item He \textbf{told} parents (that) they \textbf{had to} bring their children for routine immunization. (must $\rightarrow$ had to)
    \item He \textbf{asked if/whether} we \textbf{had distributed} the mosquito nets to the pregnant women. (yes/no question $\rightarrow$ if/whether ; past simple $\rightarrow$ past perfect)
    \item He \textbf{ordered/advised} us \textbf{to wash} our hands with soap before every meal. (imperative $\rightarrow$ to + infinitive)
\end{enumerate}
\textbf{Point de vigilance :} choisir le verbe introducteur adapté : \emph{asked} (question), \emph{ordered/warned/advised + object + (not) to} (ordre/conseil), \emph{said/told} (déclaration).
\end{corrige}

\begin{corrige}[Exercice 9 — Compréhension de lecture : La lutte contre le paludisme au Cameroun]
\textbf{1. True or False (with quotations):}
\begin{itemize}
    \item[a)] \textbf{True.} \emph{``The disease accounts for approximately 30\% of consultations in health facilities and is a leading cause of morbidity and mortality.''}
    \item[b)] \textbf{False.} \emph{``Nationwide campaigns have distributed millions of nets free of charge to households.''} (They are free, not subsidized.)
    \item[c)] \textbf{False.} \emph{``IRS is used in targeted high-transmission areas in the North and Far North regions.''} (Not in all ten regions.)
    \item[d)] \textbf{True.} \emph{``Community Health Workers educate families on environmental hygiene: clearing bushes, draining stagnant water, and covering water storage containers.''}
\end{itemize}
\textbf{2. Open questions:}
\begin{itemize}
    \item[e)] Children under five and pregnant women.
    \item[f)] IPTp = \textbf{Intermittent Preventive Treatment in pregnancy}; it benefits \textbf{pregnant women} (and their unborn babies), offered free during antenatal consultations.
    \item[g)] Any two of: misuse of nets (for fishing or gardening) ; resistance to insecticides ; self-medication with incomplete doses.
\end{itemize}
\textbf{3. Vocabulary in context:}
\begin{itemize}
    \item[h)] \textbf{mortality} (a leading cause of mortality)
    \item[i)] \textbf{the cornerstone}
    \item[j)] \textbf{self-medication}
\end{itemize}
\textbf{Barème indicatif (sur 20) :} Q1 : 4 $\times$ 2 pts (1 pt réponse + 1 pt citation exacte) = 8 pts ; Q2 : 3 $\times$ 2 pts = 6 pts ; Q3 : 3 $\times$ 2 pts = 6 pts.
\textbf{Points de vigilance :} la justification par une \emph{citation exacte} du texte est exigée (entre guillemets) ; une justification paraphrasée ne vaut que la moitié des points.
\end{corrige}

\begin{corrige}[Exercice 10 — Rédaction type Probatoire : Article pour le magazine scolaire]
\textbf{Model answer (about 180 words):}

\textbf{STOP DISEASES BEFORE THEY START!}

Our school is launching a campaign against communicable diseases, and everyone is concerned. Malaria, transmitted by the female Anopheles mosquito, causes fever and headaches. Cholera, spread through contaminated water, leads to severe diarrhoea and dehydration. Tuberculosis is a contagious airborne disease that attacks the lungs.

Fortunately, prevention is possible. Our school nurse declared: \emph{``Sleeping under a treated mosquito net every night is your best protection against malaria.''} She also said: \emph{``Wash your hands with soap before every meal and after using the toilet.''} Moreover, she advised us to drink only boiled or treated water and to keep our environment clean.

Vaccination is equally essential. The nurse told us that immunization protects the whole community, not just the vaccinated person.

Dear students, remember: \textbf{prevention is better than cure}. Sleep under a net, wash your hands, clear stagnant water around your homes, and get vaccinated. Together, let's make our school a healthy school!

\textbf{Barème indicatif (sur 20) :} respect de la tâche et du format article (titre accrocheur) : 3 pts ; contenu (2--3 maladies : causes/symptômes/transmission ; mesures de prévention) : 6 pts ; deux citations au discours direct + une au discours indirect correctement ponctuées et introduites : 4 pts ; cohérence, organisation, appel à l'action : 3 pts ; correction linguistique (grammaire, vocabulaire de l'unité, orthographe) : 4 pts.
\textbf{Points de vigilance :} guillemets et majuscule pour le discours direct ; backshift et suppression des guillemets pour le discours indirect ; rester entre 150 et 200 mots ; employer le vocabulaire de l'unité (\emph{outbreak, contagious, immunization, sanitation, handwashing}).
\end{corrige}

\begin{corrige}[Exercice 11 — Compréhension orale simulée : Campagne de reboisement scolaire]
\begin{center}
\begin{tabular}{|l|l|}
\hline
\textbf{Information Required} & \textbf{Answers} \\
\hline
Name of School & Government Technical High School Buea \\
Date of Campaign & Saturday, June 15th \\
Start Time / End Time & 8:00 a.m. / 1:00 p.m. \\
Number of Student Participants & Over 350 students \\
Number of Teachers Involved & 15 teachers \\
Donor of Seedlings & Ministry of Forestry and Wildlife (MINFOF) \\
Total Number of Seedlings & 2 000 seedlings \\
Breakdown of Species & 1. Prunus africana (1 200) \quad 2. Mahogany (500) \quad 3. Eucalyptus (300) \\
Planting Locations (2) & 1. Along the school boundaries \quad 2. Near the water catchment area \\
Principal's Quote (Direct Speech) & \emph{``We are not just planting trees; we are planting our future.''} \\
Next Year's Target & 3 000 trees \\
\hline
\end{tabular}
\end{center}
\textbf{Vérification du calcul :} $1\,200 + 500 + 300 = 2\,000$ plants, ce qui correspond bien au total annoncé. 20 groupes $\times$ 100 arbres $= 2\,000$ arbres plantés : cohérent.
\textbf{Barème indicatif (sur 10) :} 11 informations $\times$ 1 pt, arrondi à 10 (la citation doit être exacte pour valider le point).
\textbf{Points de vigilance :} noter les chiffres exacts (350 élèves, 15 enseignants) ; reproduire la citation au discours direct avec guillemets ; distinguer le donateur (MINFOF) du lieu de plantation.
\end{corrige}

\newpage

\coursec{Fiche bilan}

\begin{popfigure}
\begin{center}
\begin{center}\tcbox[colback=popdark,colframe=popdark,coltext=white,arc=9pt,boxsep=4pt,left=6pt,right=6pt]{\bfseries\small Environment, Wellbeing \& Health}\end{center}
\vspace{-2pt}
\begin{tcbraster}[raster columns=3,raster equal height=rows,raster column skip=6pt,raster row skip=6pt,enhanced,blankest]
\begin{tcolorbox}[enhanced,colback=white,colframe=popblue,boxrule=0.9pt,arc=3pt,title={Tourist sites},fonttitle=\bfseries\scriptsize,coltitle=white,colbacktitle=popblue,left=4pt,right=4pt,top=2pt,bottom=2pt,boxsep=1pt,toptitle=1.5pt,bottomtitle=1.5pt,halign=left]
\begin{itemize}[nosep,leftmargin=*,itemsep=1pt,font=\scriptsize]
\item used to / didn't use to
\item be used to + noun/-ing
\end{itemize}
\end{tcolorbox}
\begin{tcolorbox}[enhanced,colback=white,colframe=popgreen,boxrule=0.9pt,arc=3pt,title={Protecting nature},fonttitle=\bfseries\scriptsize,coltitle=white,colbacktitle=popgreen,left=4pt,right=4pt,top=2pt,bottom=2pt,boxsep=1pt,toptitle=1.5pt,bottomtitle=1.5pt,halign=left]
\begin{itemize}[nosep,leftmargin=*,itemsep=1pt,font=\scriptsize]
\item Tag questions
\item planting trees, gardening
\end{itemize}
\end{tcolorbox}
\begin{tcolorbox}[enhanced,colback=white,colframe=poporange,boxrule=0.9pt,arc=3pt,title={Health \& diseases},fonttitle=\bfseries\scriptsize,coltitle=white,colbacktitle=poporange,left=4pt,right=4pt,top=2pt,bottom=2pt,boxsep=1pt,toptitle=1.5pt,bottomtitle=1.5pt,halign=left]
\begin{itemize}[nosep,leftmargin=*,itemsep=1pt,font=\scriptsize]
\item Compound nouns
\item prevention vocabulary
\end{itemize}
\end{tcolorbox}
\begin{tcolorbox}[enhanced,colback=white,colframe=poppurple,boxrule=0.9pt,arc=3pt,title={Reported speech 1},fonttitle=\bfseries\scriptsize,coltitle=white,colbacktitle=poppurple,left=4pt,right=4pt,top=2pt,bottom=2pt,boxsep=1pt,toptitle=1.5pt,bottomtitle=1.5pt,halign=left]
\begin{itemize}[nosep,leftmargin=*,itemsep=1pt,font=\scriptsize]
\item statements: tense shift
\end{itemize}
\end{tcolorbox}
\begin{tcolorbox}[enhanced,colback=white,colframe=popgold,boxrule=0.9pt,arc=3pt,title={Reported speech 2},fonttitle=\bfseries\scriptsize,coltitle=white,colbacktitle=popgold,left=4pt,right=4pt,top=2pt,bottom=2pt,boxsep=1pt,toptitle=1.5pt,bottomtitle=1.5pt,halign=left]
\begin{itemize}[nosep,leftmargin=*,itemsep=1pt,font=\scriptsize]
\item questions \& commands
\end{itemize}
\end{tcolorbox}
\begin{tcolorbox}[enhanced,colback=white,colframe=poppink,boxrule=0.9pt,arc=3pt,title={Writing},fonttitle=\bfseries\scriptsize,coltitle=white,colbacktitle=poppink,left=4pt,right=4pt,top=2pt,bottom=2pt,boxsep=1pt,toptitle=1.5pt,bottomtitle=1.5pt,halign=left]
\begin{itemize}[nosep,leftmargin=*,itemsep=1pt,font=\scriptsize]
\item composition on health issues
\end{itemize}
\end{tcolorbox}
\end{tcbraster}

\end{center}
\legende{Carte mentale de la leçon : les six compétences clés autour du thème \emph{Environment, Wellbeing and Health}.}
\end{popfigure}

\begin{aretenir}[L'essentiel de la leçon -- Grammar: the must-know rules]
\begin{center}
\begin{tabularx}{\linewidth}{|l|X|X|}\hline
\rowcolor{popblueL} \textbf{Point} & \textbf{Rule} & \textbf{Example} \\ \hline
\cle{used to} + V & past habit or state, now finished (habitude passée révolue) & People \emph{used to} travel on foot. \\ \hline
\cle{didn't use to} + V & negative of used to (no ``d'' after \emph{did}) & We \emph{didn't use to} have hotels here. \\ \hline
\cle{be used to} + noun/-ing & to be accustomed to (être habitué à) & Tourists \emph{are used to climbing} Mount Cameroon. \\ \hline
\cle{Tag questions} & positive sentence $\to$ negative tag ; negative sentence $\to$ positive tag & You plant trees, \emph{don't you}? / It isn't late, \emph{is it}? \\ \hline
\cle{Compound nouns} & plurals: usually add -s to the \emph{last} word & toothbrush $\to$ toothbrush\emph{es} ; outbreak $\to$ outbreak\emph{s} \\ \hline
\cle{Reported speech} & backshift: present $\to$ past ; past $\to$ past perfect ; will $\to$ would ; can $\to$ could & ``I am sick'' $\to$ He said he \emph{was} sick. \\ \hline
\end{tabularx}
\end{center}
\end{aretenir}

\begin{attention}[Pièges fréquents]
\begin{itemize}
\item \emph{didn't use to} (sans \emph{d}) : on écrit ``We didn't \textbf{use} to...'' et non ``didn't used to''.
\item Ne pas confondre : \emph{I used to live} in Douala (avant, plus maintenant) et \emph{I am used to living} in Douala (j'y suis habitué).
\item Tag question : le tag reprend l'\textbf{auxiliaire} de la phrase principale, jamais un autre : ``She can swim, \emph{can't} she?'' et non ``doesn't she?''.
\item Avec \emph{there is/are} : ``There is a clinic, \emph{isn't there}?''
\end{itemize}
\end{attention}

\begin{propriete}[Reported speech -- time and place changes]
\begin{center}
\begin{tabularx}{\linewidth}{|X|X|X|X|}\hline
\rowcolor{popgreenL} \textbf{Direct} & \textbf{Indirect} & \textbf{Direct} & \textbf{Indirect} \\ \hline
now & then & today & that day \\ \hline
yesterday & the day before & tomorrow & the next day \\ \hline
here & there & this / these & that / those \\ \hline
ago & before & last week & the week before \\ \hline
\end{tabularx}
\end{center}
\end{propriete}

\begin{propriete}[Reporting questions and commands]
\begin{itemize}
\item \textbf{Yes/No questions:} \emph{asked + if/whether} + statement order (ordre sujet--verbe, pas d'inversion, pas de \emph{do/does/did}). ``Do you wash your hands?'' $\to$ She asked \emph{if I washed} my hands.
\item \textbf{Wh-questions:} keep the question word, statement order. ``Where is the clinic?'' $\to$ He asked \emph{where the clinic was}.
\item \textbf{Commands:} \emph{told/asked + object + (not) to} + V. ``Wash your hands!'' $\to$ The nurse told us \emph{to wash} our hands. ``Don't touch the seedlings!'' $\to$ He told us \emph{not to touch} the seedlings.
\end{itemize}
\end{propriete}

\begin{exemplebox}[Key vocabulary]
\begin{itemize}
\item \textbf{Tourism:} tourist site (site touristique), landscape (paysage), wildlife (faune et flore), guide (guide), souvenir (souvenir), to hike (randonner).
\item \textbf{Environment:} to plant trees (planter des arbres), gardening (jardinage), natural habitat (habitat naturel), to protect (protéger), seedling (jeune plant).
\item \textbf{Health:} communicable disease (maladie transmissible), to prevent (prévenir), hygiene (hygiène), symptom (symptôme), vaccine (vaccin), outbreak (épidémie), handwashing (lavage des mains).
\end{itemize}
\end{exemplebox}

\begin{methode}[Express methods]
\begin{enumerate}
\item \textbf{Tag question:} find the auxiliary of the main verb, reverse the polarity (positive $\leftrightarrow$ negative), add the pronoun.
\item \textbf{Reported speech:} (1) identify statement / question / command ; (2) choose \emph{say/ask/tell} ; (3) backshift tenses ; (4) change time and place words ; (5) change pronouns.
\item \textbf{Writing a composition:} introduction (the health issue) $\to$ body (causes, prevention) $\to$ conclusion (advice to the community).
\end{enumerate}
\end{methode}

\begin{exoresolu}[Reported speech -- un exemple complet]
Énoncé. Report the following sentences. (a) The nurse said: ``I vaccinate children every day.'' (b) The teacher asked: ``Did you plant the trees yesterday?'' (c) The guide said to the tourists: ``Don't feed the animals.''
\tcblower
Solution détaillée.
\begin{enumerate}
\item (a) Statement $\to$ \emph{say} + backshift (present simple $\to$ past simple ; every day $\to$ every day, inchangé) : The nurse said (that) she \emph{vaccinated} children every day.
\item (b) Yes/No question $\to$ \emph{asked + if} + statement order + backshift (past simple $\to$ past perfect ; yesterday $\to$ the day before) : The teacher asked \emph{if we had planted} the trees the day before.
\item (c) Negative command $\to$ \emph{told + object + not to} + V : The guide told the tourists \emph{not to feed} the animals.
\end{enumerate}
\end{exoresolu}

\begin{aretenir}[Checklist avant le Bac]
\begin{itemize}
\item[$\square$] I can use \emph{used to} and \emph{didn't use to} for past habits.
\item[$\square$] I can distinguish \emph{used to + V} from \emph{be used to + -ing}.
\item[$\square$] I can build correct tag questions with the right auxiliary and polarity.
\item[$\square$] I know the plural of compound nouns (toothbrushes, outbreaks, grown-ups).
\item[$\square$] I can report statements with the correct tense backshift.
\item[$\square$] I can report yes/no questions with \emph{if/whether} and wh-questions with statement word order.
\item[$\square$] I can report commands with \emph{told/asked + object + to + V}.
\item[$\square$] I can change time and place expressions (now $\to$ then, here $\to$ there, today $\to$ that day).
\item[$\square$] I master vocabulary about tourist sites, tree planting and disease prevention.
\item[$\square$] I can write a short composition giving advice on preventing communicable diseases.
\end{itemize}
\end{aretenir}

\findecours

\end{document}
